\documentclass[11pt,openany]{book}

\usepackage[paperwidth=17cm,paperheight=24cm,top=2.3cm,bottom=2.3cm,inner=2.3cm,outer=1.9cm,headsep=0.6cm]{geometry}
\usepackage{fontspec}
\usepackage{longtable}
\usepackage{xepersian}
\usepackage{xcolor}
\usepackage{graphicx}
\usepackage{listings}
\usepackage{framed}
\usepackage{fancyvrb}
\usepackage{amsmath}
\usepackage{amssymb}
\usepackage{booktabs}
\usepackage{array}
\usepackage{hyperref}
\usepackage{titlesec}
\usepackage{needspace}
\usepackage{fancyhdr}
\usepackage{enumitem}

\makeatletter
\providecommand\UseMathForPositioningText{}
\renewcommand\@pnumwidth{2.3em}
\renewcommand\@tocrmarg{3.3em}
\renewcommand*\l@section{\@dottedtocline{1}{1.5em}{2.8em}}
\makeatother

\settextfont[
Path=./../../fonts/,
Extension=.ttf,
BoldFont = {* Bd},
ItalicFont = {* It},
BoldItalicFont = {* BdIt}
]{XB Zar}
\setlatintextfont{Latin Modern Roman}

\newfontfamily\titlefont[
Path=./../../fonts/,
Extension=.ttf,
Mapping=farsidigits
]{XB Titre}
\newfontfamily\headfont[
Path=./../../fonts/,
Extension=.ttf,
BoldFont = {* Bd},
ItalicFont = {* It},
BoldItalicFont = {* BdIt},
Mapping=farsidigits
]{XM Yekan}
\newfontfamily\salamcodefont[
Path=./../../fonts/,
Extension=.ttf,
BoldFont = {*Bd},
ItalicFont = {*It},
BoldItalicFont = {*BdIt},
Scale = 0.92
]{XB Roya}
\newfontfamily\latincodefont[
Path=./../../fonts/,
Extension=.ttf,
UprightFont = DejaVuSansMono,
BoldFont = DejaVuSansMono-Bold,
ItalicFont = DejaVuSansMono-Oblique,
BoldItalicFont = DejaVuSansMono-BoldOblique,
Scale = 0.85
]{DejaVuSansMono}

\definecolor{salamKeyword}{HTML}{1E40FF}
\definecolor{salamType}{HTML}{267F99}
\definecolor{salamString}{HTML}{A31515}
\definecolor{salamComment}{HTML}{008000}
\definecolor{salamNumber}{HTML}{098658}
\definecolor{brand}{HTML}{0B6E4F}
\definecolor{brandLight}{HTML}{E6F2EE}
\definecolor{ink}{HTML}{1F2933}
\definecolor{labelColor}{HTML}{0B6E4F}
\definecolor{noteBg}{HTML}{EAF3FB}
\definecolor{noteFrame}{HTML}{3E7CB1}
\definecolor{tipBg}{HTML}{ECF8EE}
\definecolor{tipFrame}{HTML}{3C9A5F}
\definecolor{warnBg}{HTML}{FDF1E6}
\definecolor{warnFrame}{HTML}{D9822B}
\definecolor{deepBg}{HTML}{F3EFFA}
\definecolor{deepFrame}{HTML}{8A63C2}
\definecolor{exBg}{HTML}{FFF8E1}
\definecolor{exFrame}{HTML}{C9A227}
\definecolor{sumBg}{HTML}{F1F3F5}
\definecolor{sumFrame}{HTML}{495057}

\newcommand{\salamcodeENfont}{\latincodefont}
\newcommand{\salamcodeFAfont}{\salamcodefont}
\usepackage{listings}
\usepackage{fancyvrb}
\usepackage{xstring}

\providecommand{\salamcodeENfont}{\ttfamily}
\providecommand{\salamcodeFAfont}{\normalfont}

\definecolor{bg}{HTML}{F4F4F2}
\definecolor{kw}{HTML}{1E40FF}   %
\definecolor{ty}{HTML}{267F99}   %
\definecolor{cm}{HTML}{008000}   %
\definecolor{st}{HTML}{A31515}   %
\definecolor{nm}{HTML}{098658}   %
\definecolor{outBg}{RGB}{244,244,244}
\definecolor{outFrame}{RGB}{200,200,200}

\newcommand{\kw}[1]{\textcolor{kw}{\textbf{#1}}}
\newcommand{\tyword}[1]{\textcolor{ty}{#1}}
\newcommand{\cmtword}[1]{\textcolor{cm}{\textit{#1}}}
\newcommand{\stword}[1]{\textcolor{st}{#1}}
\newcommand{\nmword}[1]{\textcolor{nm}{#1}}

\lstdefinelanguage{Salam}{
    morekeywords={func,ret,if,else,until,each,in,mut,const,type,struct,enum,end,%
        import,include,as,true,false,null,this,break,continue,layout,package,print,%
        println,printerr,printerrln,input,defer,operator,extern,interface,%
        pub,component,repeat,impl,on,to,by,main,sizeof,spawn,join,link},
    morekeywords=[2]{void,bool,char,str,int,uint,float,auto,byte,%
        HashMap,Vector,MapIter,Option,Stack,Queue,Thread,Mutex,WaitGroup},
    sensitive=true,
    morecomment=[l]{//},
    morecomment=[n]{/*}{*/},
    morestring=[b]", %
    morestring=[b]',
    morestring=[b]`,
    literate=%
        {i8}{{\tyword{i8}}}2   {i16}{{\tyword{i16}}}3 {i32}{{\tyword{i32}}}3%
        {i64}{{\tyword{i64}}}3 {u8}{{\tyword{u8}}}2   {u16}{{\tyword{u16}}}3%
        {u32}{{\tyword{u32}}}3 {u64}{{\tyword{u64}}}3 {f32}{{\tyword{f32}}}3%
        {f64}{{\tyword{f64}}}3%
        {0}{{\nmword{0}}}1 {1}{{\nmword{1}}}1 {2}{{\nmword{2}}}1%
        {3}{{\nmword{3}}}1 {4}{{\nmword{4}}}1 {5}{{\nmword{5}}}1%
        {6}{{\nmword{6}}}1 {7}{{\nmword{7}}}1 {8}{{\nmword{8}}}1%
        {9}{{\nmword{9}}}1,
}

\lstdefinestyle{salamcardstyle}{
    language=Salam,
    basicstyle=\salamcodeENfont\small,
    backgroundcolor=\color{bg},
    frame=single,
    numbers=left,
    keywordstyle=\color{kw}\bfseries,
    keywordstyle=[2]\color{ty},
    commentstyle=\color{cm}\itshape,
    stringstyle=\color{st},
    showstringspaces=false,
    breaklines=true,
    breakatwhitespace=false,
    tabsize=4,
    columns=fullflexible,
    keepspaces=true,
    upquote=true,
    extendedchars=true,
}

\lstnewenvironment{salamen}[1][]{\setLTR\lstset{style=salamcardstyle,#1}}{}

\lstdefinestyle{shellstyle}{
  basicstyle=\salamcodeENfont\small,
  backgroundcolor=\color{outBg},
  frame=single,
  rulecolor=\color{outFrame},
  framesep=6pt,
  numbers=none,
  showstringspaces=false,
  breaklines=true,
  columns=fullflexible,
  keepspaces=true,
  extendedchars=true,
}
\lstnewenvironment{progout}[1][]{\setLTR\lstset{style=shellstyle,#1}}{}
\lstnewenvironment{shell}[1][]{\setLTR\lstset{style=shellstyle,#1}}{}

\makeatletter

\providecommand{\latincodefont}{\salamcodeENfont}
\providecommand{\lr}[1]{\LRE{#1}}

\@for\salamfa@x:=%
  func,ret,if,else,until,each,in,mut,const,type,struct,enum,end,import,as,true,%
  false,null,this,break,continue,layout,package,print,println,printerr,%
  printerrln,input,defer,operator,extern,interface,pub,component,repeat,impl,%
  on,to,by,main,sizeof,spawn,join,link,%
  روال,برگشت,اگر,وگرنه,تا,بر,ناپایا,پایا,گونه,ساختار,جداشمار,پایان,%
  واردسازی,برگردان,درست,نادرست,پوچ,این,بشکن,گذر,چیدمان,بسته,چاپ,سرچاپ,%
  نادرست‌چاپ,نادرست‌سرچاپ,ورودی,دیرکن,کارور,فراخوانه,میانجی,همگانی,درخط,%
  نادرخط,ناب,نابرگشت,بی‌کاره,بخش,تکرار,کاربست,هر,در,همخوان,و,یا,برابر,%
  نابرابر,ترابرد,ریشه\do{%
    \expandafter\let\csname salamkw@\salamfa@x\endcsname\@empty}

\@for\salamfa@x:=%
  void,bool,char,str,i8,i16,i32,int,i64,u8,u16,u32,uint,u64,f32,float,f64,%
  auto,byte,HashMap,Vector,MapIter,Option,Stack,Queue,Thread,Mutex,WaitGroup,%
  تهی,منطقی,نویسه,رشته,صحیح۸,صحیح۱۶,صحیح۳۲,صحیح,صحیح۶۴,صحیح8,صحیح16,صحیح32,%
  صحیح64,طبیعی۸,طبیعی۱۶,طبیعی۳۲,طبیعی,طبیعی۶۴,طبیعی8,طبیعی16,طبیعی32,طبیعی64,%
  اعشار۳۲,اعشار,اعشار۶۴,اعشار32,اعشار64,خودکار,وکتور,نگاشت\do{%
    \expandafter\let\csname salamty@\salamfa@x\endcsname\@empty}

% ---------------------------------------------------------------------------
% How a line is drawn.
%
% XeTeX only runs the bidi algorithm inside a single word, so a code line is
% never handed to it as running text.  Each line is cut into items (words,
% operator characters, spaces, string literals, comments) and the items are
% then handed to the right-to-left paragraph one box at a time, with the
% direction of each box and the shape of each bracket decided here.
%
%  * A line holding any persian letter reads right to left from its first
%    character to its last: the items come out in reverse order and every
%    operator character is mirrored, so `:=' shows as `=:', `>=' as `=<' and
%    `(' as `)'.  Only a word or a number keeps its own inner direction.
%  * A line with no persian letter keeps each space-separated token exactly as
%    written, but the tokens still run right to left, so the first token of
%    the line sits against the right edge.
%  * Inside a string literal or a comment, neighbouring latin words and the
%    punctuation between them stay together as one left-to-right island, the
%    way ordinary mixed persian/english text reads.
% ---------------------------------------------------------------------------

\newdimen\salamfa@unit \salamfa@unit=0.5em
\newdimen\salamfa@spw  \salamfa@spw=0.4em
\newcount\salamfa@lead
\newcount\salamfa@nlines
\newcount\salamfa@c
\newcount\salamfa@mode      % 0 code, 1 string, 2 line comment, 3 block comment
\newcount\salamfa@qchar
\newcount\salamfa@spc
\newcount\salamfa@ctx       % 0 code, 1 string, 2 comment, 3 program output
\newcount\salamfa@signc
\newif\ifsalamfa@out
\newif\ifsalamfa@lineP
\newif\ifsalamfa@started
\newif\ifsalamfa@esc
\newif\ifsalamfa@dotp
\newif\ifsalamfa@slashp
\newif\ifsalamfa@starp
\newif\ifsalamfa@lastsig
\newif\ifsalamfa@wP
\newif\ifsalamfa@wF
\newif\ifsalamfa@wnum
\newif\ifsalamfa@isl
\newif\ifsalamfa@leadcmt
\newcommand{\salamfa@code}{}
\def\salamfa@SP{\salamfa@SP}
\def\salamfa@END{\salamfa@END}
\protected\def\salamfa@end{}
\def\salamfa@endmac{\salamfa@end}

\def\salamfa@gapp#1#2{\xdef#1{\unexpanded\expandafter{#1}\unexpanded{#2}}}
\def\salamfa@xapp#1#2{\edef\salamfa@t{#2}%
  \expandafter\salamfa@gapp\expandafter#1\expandafter{\salamfa@t}}

% ---- character classes ------------------------------------------------------
% 0 neutral, 1 latin letter or _, 2 ascii digit, 3 persian letter,
% 4 persian digit, 5 neutral drawn with the persian font
\def\salamfa@classof#1{%
  \salamfa@c\z@
  \ifnum`#1<128
    \ifnum`#1>47 \ifnum`#1<58 \salamfa@c\tw@\fi\fi
    \ifnum`#1>64 \ifnum`#1<91 \salamfa@c\@ne\fi\fi
    \ifnum`#1>96 \ifnum`#1<123 \salamfa@c\@ne\fi\fi
    \ifnum`#1=95 \salamfa@c\@ne\fi
  \else\ifnum`#1<1536
    \salamfa@c\@ne
    \ifnum`#1=171 \salamfa@c5 \fi
    \ifnum`#1=187 \salamfa@c5 \fi
  \else\ifnum`#1<1792
    \salamfa@c\thr@@
    \ifnum`#1>1775 \ifnum`#1<1786 \salamfa@c4 \fi\fi
    \ifnum`#1>1631 \ifnum`#1<1642 \salamfa@c4 \fi\fi
    \ifnum`#1=1548 \salamfa@c5 \fi
    \ifnum`#1=1563 \salamfa@c5 \fi
    \ifnum`#1=1567 \salamfa@c5 \fi
    \ifnum`#1>1641 \ifnum`#1<1646 \salamfa@c5 \fi\fi
    \ifnum`#1=1748 \salamfa@c5 \fi
  \else\ifnum`#1<8192
    \salamfa@c\@ne
  \else\ifnum`#1<8304
    \salamfa@c5
    \ifnum`#1=8204 \salamfa@c\thr@@\fi
    \ifnum`#1=8205 \salamfa@c\thr@@\fi
  \else\ifnum`#1<64336
    \salamfa@c5
  \else\ifnum`#1<65280
    \salamfa@c\thr@@
  \else
    \salamfa@c5
  \fi\fi\fi\fi\fi\fi\fi}

% ---- mirrored glyphs ----------------------------------------------------------
\edef\salamfa@lb{\expandafter\@gobble\string\{}
\edef\salamfa@rb{\expandafter\@gobble\string\}}
\def\salamfa@defmir#1#2{\expandafter\xdef\csname salamfa@mir@#1\endcsname{#2}}
\salamfa@defmir{(}{)}\salamfa@defmir{)}{(} % chktex 10
\salamfa@defmir{[}{]}\salamfa@defmir{]}{[}
\salamfa@defmir{<}{>}\salamfa@defmir{>}{<}
\salamfa@defmir{«}{»}\salamfa@defmir{»}{«}
\expandafter\salamfa@defmir\expandafter{\salamfa@lb}{\salamfa@rb}
\expandafter\salamfa@defmir\expandafter{\salamfa@rb}{\salamfa@lb}

% ---- pieces: what finally gets typeset -----------------------------------------
% xepersian's inter-character hooks rewrite a number such as 12.5 for running
% text, which scrambles it here, so they are switched off in left-to-right boxes.
\def\salamfa@pL#1#2{\hbox{\LRE{\XeTeXinterchartokenstate\z@\salamcodeENfont#1{#2}}}}
\def\salamfa@pF#1#2{\hbox{\LRE{\XeTeXinterchartokenstate\z@
  \salamcodeFAfont#1{\salamfa@fadots#2\relax}}}}
\def\salamfa@pP#1#2{\hbox{\RLE{\salamcodeFAfont#1{#2}}}}
\def\salamfa@pI#1{\hbox{\LRE{\XeTeXinterchartokenstate\z@#1}}}
% the persian font turns a dot between persian digits into a decimal comma;
% keep the dot the source actually has
\def\salamfa@fadots#1{%
  \ifx#1\relax\else
    \ifnum`#1=46 \kern\z@.\kern\z@\else#1\fi
    \expandafter\salamfa@fadots
  \fi}
\def\salamfa@pS{\hskip\salamfa@spw\relax}

% ---- lexer ------------------------------------------------------------------------
\def\salamfa@add#1{%
  \ifnum\salamfa@mode=\z@ \salamfa@gapp\salamfa@items{#1}%
  \else \salamfa@gapp\salamfa@sub{#1}\fi}
\def\salamfa@xadd#1{\edef\salamfa@t{#1}\expandafter\salamfa@add\expandafter{\salamfa@t}}

% spaces are held back until something follows them, so trailing ones vanish
% and leading ones become the indentation
\def\salamfa@flushsp{%
  \ifsalamfa@started
    \@whilenum\salamfa@spc>\z@\do{\salamfa@add{\salamfa@iS}\advance\salamfa@spc\m@ne}%
  \else
    \salamfa@lead\salamfa@spc \salamfa@startedtrue
  \fi
  \salamfa@spc\z@}

\def\salamfa@wreset{%
  \gdef\salamfa@wb{}\salamfa@wPfalse\salamfa@wFfalse\salamfa@wnumfalse}

\def\salamfa@flushword{%
  \ifx\salamfa@wb\@empty\else
    \ifsalamfa@wP \def\salamfa@k{P}\else\ifsalamfa@wF \def\salamfa@k{F}\else
      \def\salamfa@k{L}\fi\fi
    \salamfa@xadd{\noexpand\salamfa@iW{\salamfa@wb}{\salamfa@k}%
      {\ifsalamfa@wnum 1\else 0\fi}}%
    \salamfa@lastsigtrue
  \fi
  \salamfa@wreset}

% one operator or punctuation character
\def\salamfa@op#1{%
  \salamfa@flushword \salamfa@flushsp
  \salamfa@classof{#1}%
  \salamfa@xadd{\noexpand\salamfa@iO{\unexpanded{#1}}{\ifnum\salamfa@c=5 P\else L\fi}}%
  \salamfa@lastsigfalse
  \ifnum`#1=41 \salamfa@lastsigtrue\fi
  \ifnum`#1=93 \salamfa@lastsigtrue\fi
  \ifnum`#1=125 \salamfa@lastsigtrue\fi}

% settle what the previous character left open, now that the next one (#1)
% is known: a dot inside a word (p.x, 12.5) and the sign of a number (-5)
\def\salamfa@resolve#1{%
  \ifsalamfa@dotp
    \salamfa@dotpfalse
    \ifx#1\salamfa@SP \salamfa@op.\else
      \salamfa@classof{#1}%
      \ifnum\salamfa@c>\z@ \ifnum\salamfa@c<5
        \g@addto@macro\salamfa@wb{.}\else\salamfa@op.\fi
      \else \salamfa@op.\fi
    \fi
  \fi
  \ifnum\salamfa@signc>\z@
    \ifx#1\salamfa@SP \salamfa@signop\else
      \salamfa@classof{#1}%
      \ifnum\salamfa@c=\tw@ \salamfa@signword\else
      \ifnum\salamfa@c=4 \salamfa@signword\else \salamfa@signop\fi\fi
    \fi
    \salamfa@signc\z@
  \fi}
\def\salamfa@signword{%
  \salamfa@flushsp
  \ifnum\salamfa@signc=45 \gdef\salamfa@wb{-}\else \gdef\salamfa@wb{+}\fi
  \salamfa@wnumtrue}
\def\salamfa@signop{%
  \ifnum\salamfa@signc=45 \salamfa@op-\else \salamfa@op+\fi}

\def\salamfa@wordchar#1{%
  \ifx\salamfa@wb\@empty
    \salamfa@flushsp
    \ifnum\salamfa@c=\tw@ \salamfa@wnumtrue\fi
    \ifnum\salamfa@c=4 \salamfa@wnumtrue\fi
  \fi
  \g@addto@macro\salamfa@wb{#1}%
  \ifnum\salamfa@c=\thr@@ \salamfa@wPtrue \global\salamfa@linePtrue\fi
  \ifnum\salamfa@c=4 \salamfa@wFtrue\fi}

\def\salamfa@char#1{%
  \ifx#1\salamfa@SP
    \salamfa@resolve{#1}%
    \ifsalamfa@slashp \salamfa@slashpfalse \salamfa@op/\fi
    \ifsalamfa@starp \salamfa@starpfalse \salamfa@op*\fi
    \salamfa@flushword
    \advance\salamfa@spc\@ne
  \else
    \ifsalamfa@out
      \salamfa@classof{#1}%
      \ifnum\salamfa@c>\tw@ \global\salamfa@linePtrue\fi
    \fi
    \salamfa@charB{#1}%
  \fi}

\def\salamfa@charB#1{%
  \ifcase\salamfa@mode
    % code: a slash may open a comment
    \ifsalamfa@slashp
      \salamfa@slashpfalse
      \ifnum`#1=47 \salamfa@opencmt2\else
      \ifnum`#1=42 \salamfa@opencmt3\else
        \salamfa@op/\salamfa@charC{#1}\fi\fi
    \else \salamfa@charC{#1}\fi
  \or
    % string literal
    \ifsalamfa@esc
      \salamfa@escfalse \salamfa@c\@ne \salamfa@wordchar{#1}%
    \else\ifnum`#1=\salamfa@qchar
      \salamfa@resolve{#1}\salamfa@op{#1}%
      \global\salamfa@mode\z@
      \salamfa@xadd{\noexpand\salamfa@iT{\unexpanded\expandafter{\salamfa@sub}}}%
      \salamfa@lastsigtrue
    \else\ifnum`#1=92 \ifnum\salamfa@qchar=96 \salamfa@charC{#1}\else
      \salamfa@resolve{#1}\salamfa@esctrue \salamfa@c\@ne \salamfa@wordchar{#1}\fi
    \else \salamfa@charC{#1}\fi\fi\fi
  \or
    \salamfa@charC{#1}%
  \or
    % block comment: look for the closing */
    \ifsalamfa@starp
      \salamfa@starpfalse
      \ifnum`#1=47
        \salamfa@op*\salamfa@op/\salamfa@closecmt
      \else \salamfa@op*\salamfa@charD{#1}\fi
    \else \salamfa@charD{#1}\fi
  \fi}
\def\salamfa@charD#1{%
  \ifnum`#1=42 \salamfa@resolve{#1}\salamfa@flushword\salamfa@starptrue
  \else \salamfa@charC{#1}\fi}

% what every mode shares: words, numbers and plain operators
\def\salamfa@charC#1{%
  \salamfa@resolve{#1}%
  \salamfa@classof{#1}%
  \ifnum\salamfa@c>\z@ \ifnum\salamfa@c<5
    \salamfa@wordchar{#1}%
  \else \salamfa@charE{#1}\fi
  \else \salamfa@charE{#1}\fi}
\def\salamfa@charE#1{%
  \ifnum`#1=46
    \ifx\salamfa@wb\@empty \salamfa@op{#1}\else \salamfa@dotptrue\fi
  \else\ifnum\salamfa@mode=\z@ \salamfa@charF{#1}%
  \else \salamfa@charG{#1}\fi\fi}
\def\salamfa@charF#1{%
  \ifsalamfa@out \salamfa@charG{#1}\else
  \ifnum`#1=34 \salamfa@openstr{#1}\else
  \ifnum`#1=39 \salamfa@openstr{#1}\else
  \ifnum`#1=96 \salamfa@openstr{#1}\else
  \ifnum`#1=47 \salamfa@flushword\salamfa@slashptrue\else
    \salamfa@charG{#1}\fi\fi\fi\fi\fi}
\def\salamfa@charG#1{%
  \ifx\salamfa@wb\@empty
    \ifsalamfa@lastsig \salamfa@op{#1}\else
      \ifnum`#1=45 \salamfa@signc`#1\relax\else
      \ifnum`#1=43 \salamfa@signc`#1\relax\else \salamfa@op{#1}\fi\fi
    \fi
  \else \salamfa@op{#1}\fi}

\def\salamfa@openstr#1{%
  \salamfa@flushword \salamfa@flushsp
  \gdef\salamfa@sub{}%
  \global\salamfa@mode\@ne \salamfa@qchar`#1\relax
  \salamfa@op{#1}}

\def\salamfa@opencmt#1{%
  \salamfa@flushword \salamfa@flushsp
  \gdef\salamfa@sub{}%
  \ifnum#1=\tw@
    \gdef\salamfa@cmark{\salamfa@iO{/}{L}\salamfa@iO{/}{L}}%
  \else
    \gdef\salamfa@cmark{\salamfa@iO{/}{L}\salamfa@iO{*}{L}}%
  \fi
  \global\salamfa@mode#1\relax}
\def\salamfa@emitcmt{%
  \salamfa@flushword
  \salamfa@spc\z@
  \edef\salamfa@t{\noexpand\salamfa@iC{\unexpanded\expandafter{\salamfa@cmark}}%
    {\unexpanded\expandafter{\salamfa@sub}}}%
  \expandafter\salamfa@gapp\expandafter\salamfa@items\expandafter{\salamfa@t}%
  \gdef\salamfa@sub{}\gdef\salamfa@cmark{}}
\def\salamfa@closecmt{%
  \global\salamfa@mode\z@
  \salamfa@emitcmt
  \salamfa@lastsigfalse}

\def\salamfa@lex#1{%
  \ifx#1\salamfa@END\else
    \salamfa@char{#1}%
    \expandafter\salamfa@lex
  \fi}

% cut the detokenized line at its spaces
\def\salamfa@split#1 {%
  \def\salamfa@tmp{#1}%
  \ifx\salamfa@tmp\salamfa@endmac\else
    \g@addto@macro\salamfa@chars{#1\salamfa@SP}%
    \expandafter\salamfa@split
  \fi}

% ---- layout -------------------------------------------------------------------------
% The line itself is a right-to-left paragraph, so pieces are handed over in
% reading order and TeX lays them out from the right edge.  Every piece is a
% box, which is what lets a long line wrap onto the next row in the right
% order.  A left-to-right island goes in as one box.
\def\salamfa@pre#1{\xdef\salamfa@vis{\unexpanded\expandafter{\salamfa@vis}\unexpanded{#1}}}
\def\salamfa@preI#1{\ifx#1\@empty\else
  \edef\salamfa@t{\noexpand\salamfa@pI{\unexpanded\expandafter{#1}}}%
  \expandafter\salamfa@pre\expandafter{\salamfa@t}\fi}

% style of a word, given its text (#1) and whether it is a number (#2)
\def\salamfa@setsty#1#2{%
  \ifcase\salamfa@ctx
    \ifcsname salamkw@#1\endcsname \def\salamfa@sty{\kw}%
    \else\ifcsname salamty@#1\endcsname \def\salamfa@sty{\tyword}%
    \else\ifnum#2=\@ne \def\salamfa@sty{\salamfaNM}%
    \else \def\salamfa@sty{\salamfa@id}\fi\fi\fi
  \or \def\salamfa@sty{\stword}%
  \or \def\salamfa@sty{\cmtword}%
  \or \def\salamfa@sty{\salamfa@id}%
  \fi}
\def\salamfa@setopsty{%
  \ifcase\salamfa@ctx \def\salamfa@sty{\salamfa@id}%
  \or \def\salamfa@sty{\stword}%
  \or \def\salamfa@sty{\cmtword}%
  \or \def\salamfa@sty{\salamfa@id}\fi}

% build \salamfa@piece for a word, or for an operator (#3 = 1 mirrors it)
\def\salamfa@mkW#1#2#3{%
  \salamfa@setsty{#1}{#3}%
  \if#2P\def\salamfa@k{\salamfa@pP}\else
  \if#2F\def\salamfa@k{\salamfa@pF}\else \def\salamfa@k{\salamfa@pL}\fi\fi
  \edef\salamfa@piece{\expandafter\noexpand\salamfa@k
    {\expandafter\noexpand\salamfa@sty}{\unexpanded{#1}}}}
\def\salamfa@mkO#1#2#3{%
  \salamfa@setopsty
  \def\salamfa@m{#1}%
  \ifnum#3=\@ne
    \ifcsname salamfa@mir@#1\endcsname
      \edef\salamfa@m{\csname salamfa@mir@#1\endcsname}\fi
  \fi
  \if#2P\def\salamfa@k{\salamfa@pF}\else \def\salamfa@k{\salamfa@pL}\fi
  \edef\salamfa@piece{\expandafter\noexpand\salamfa@k
    {\expandafter\noexpand\salamfa@sty}{\unexpanded\expandafter{\salamfa@m}}}}

% handlers that append in reading order, inside a left-to-right run
\def\salamfa@hLtr#1{%
  \def\salamfa@iW##1##2##3{\salamfa@mkW{##1}{##2}{##3}%
    \expandafter\salamfa@gapp\expandafter#1\expandafter{\salamfa@piece}}%
  \def\salamfa@iO##1##2{\salamfa@mkO{##1}{##2}0%
    \expandafter\salamfa@gapp\expandafter#1\expandafter{\salamfa@piece}}%
  \def\salamfa@iS{\salamfa@gapp#1{\salamfa@pS}}}

% handlers for a right-to-left code line: everything comes out reversed
\def\salamfa@hRcode{%
  \def\salamfa@iW##1##2##3{\salamfa@mkW{##1}{##2}{##3}%
    \expandafter\salamfa@pre\expandafter{\salamfa@piece}}%
  \def\salamfa@iO##1##2{\salamfa@mkO{##1}{##2}1%
    \expandafter\salamfa@pre\expandafter{\salamfa@piece}}%
  \def\salamfa@iS{\salamfa@pre{\salamfa@pS}}%
  \def\salamfa@iT##1{\begingroup\salamfa@ctx\@ne\salamfa@hRtext
    ##1\salamfa@closeisl\endgroup}%
  \def\salamfa@iC##1##2{\begingroup\salamfa@ctx\tw@\salamfa@hRtext
    ##1##2\salamfa@closeisl\endgroup}}

% handlers for right-to-left running text: latin words, and the punctuation
% between two of them, form left-to-right islands
\def\salamfa@hRtext{%
  \global\salamfa@islfalse
  \gdef\salamfa@islb{}\gdef\salamfa@pendL{}\gdef\salamfa@pendR{}%
  \def\salamfa@iW##1##2##3{\salamfa@mkW{##1}{##2}{##3}%
    \if##2L%
      \ifsalamfa@isl
        \expandafter\salamfa@gapp\expandafter\salamfa@islb\expandafter{\salamfa@pendL}%
        \gdef\salamfa@pendL{}\gdef\salamfa@pendR{}%
      \else \global\salamfa@isltrue \fi
      \expandafter\salamfa@gapp\expandafter\salamfa@islb\expandafter{\salamfa@piece}%
    \else
      \let\salamfa@savepiece\salamfa@piece
      \salamfa@closeisl
      \expandafter\salamfa@pre\expandafter{\salamfa@savepiece}%
    \fi}%
  \def\salamfa@iO##1##2{%
    \ifsalamfa@isl
      \salamfa@mkO{##1}{##2}0%
      \expandafter\salamfa@gapp\expandafter\salamfa@pendL\expandafter{\salamfa@piece}%
      \salamfa@mkO{##1}{##2}1%
      \salamfa@xapp\salamfa@pendR{\noexpand\salamfa@pre{\unexpanded\expandafter{\salamfa@piece}}}%
    \else
      \salamfa@mkO{##1}{##2}1%
      \expandafter\salamfa@pre\expandafter{\salamfa@piece}%
    \fi}%
  \def\salamfa@iS{%
    \ifsalamfa@isl
      \salamfa@gapp\salamfa@pendL{\salamfa@pS}%
      \salamfa@gapp\salamfa@pendR{\salamfa@pre{\salamfa@pS}}%
    \else \salamfa@pre{\salamfa@pS}\fi}}
\def\salamfa@closeisl{%
  \ifsalamfa@isl
    \global\salamfa@islfalse
    \salamfa@preI\salamfa@islb
    \salamfa@pendR
    \gdef\salamfa@islb{}\gdef\salamfa@pendL{}\gdef\salamfa@pendR{}%
  \fi}

% handlers for a latin code line: tokens right to left, each kept as written
\def\salamfa@flushtok{%
  \salamfa@preI\salamfa@tokb
  \gdef\salamfa@tokb{}}
\def\salamfa@hLcode{%
  \gdef\salamfa@tokb{}%
  \salamfa@hLtr\salamfa@tokb
  \def\salamfa@iS{\salamfa@flushtok\salamfa@pre{\salamfa@pS}}%
  \def\salamfa@iT##1{\begingroup\salamfa@ctx\@ne
    \salamfa@hLtr\salamfa@tokb ##1\endgroup}%
  % a comment: the // marker first, then its text as one left-to-right run
  \def\salamfa@iC##1##2{\salamfa@flushtok
    \begingroup\salamfa@ctx\tw@
      \gdef\salamfa@cmtb{}\salamfa@hLtr\salamfa@cmtb ##1%
      \salamfa@preI\salamfa@cmtb
      \gdef\salamfa@cmtb{}%
      \salamfa@leadcmttrue
      \def\salamfa@iS{\ifsalamfa@leadcmt\salamfa@pre{\salamfa@pS}\else
        \salamfa@gapp\salamfa@cmtb{\salamfa@pS}\fi}%
      \def\salamfa@iW####1####2####3{\salamfa@leadcmtfalse
        \salamfa@mkW{####1}{####2}{####3}%
        \expandafter\salamfa@gapp\expandafter\salamfa@cmtb\expandafter{\salamfa@piece}}%
      \def\salamfa@iO####1####2{\salamfa@leadcmtfalse\salamfa@mkO{####1}{####2}0%
        \expandafter\salamfa@gapp\expandafter\salamfa@cmtb\expandafter{\salamfa@piece}}%
      ##2%
      \salamfa@preI\salamfa@cmtb
    \endgroup}}

% ---- arithmetic islands --------------------------------------------------------
% An expression made of numbers and latin names, such as 2026 - 2008 or
% p.x * p.x + p.y, must keep its left-to-right order: read backwards, a minus
% or a division means something else.  Such a run is drawn as one
% left-to-right box, in persian lines and latin lines alike.  It may hold
% spaces, + - * / % ^ and balanced brackets, but it needs at least one of those
% operators, and a bracket glued to a persian name (a call or an index on
% it) never opens one, so an expression led by a persian name still reads
% right to left.
%
% The items of a line are first copied into numbered slots with a join code:
% 0 never joins, 1 operand, 2 operator, 3 space, 4 opening, 5 closing bracket.
\newcount\salamfa@n
\newcount\salamfa@N
\newif\ifsalamfa@prevP
\newif\ifsalamfa@go
\newif\ifsalamfa@hasop
\def\salamfa@J#1{\csname salamfa@J\number#1\endcsname}
\def\salamfa@store#1#2{%
  \global\advance\salamfa@N\@ne
  \expandafter\gdef\csname salamfa@A\the\salamfa@N\endcsname{#1}%
  \expandafter\xdef\csname salamfa@J\the\salamfa@N\endcsname{#2}%
  \expandafter\global\expandafter\let\csname salamfa@I\the\salamfa@N\endcsname\relax}
\def\salamfa@jof#1{%
  \def\salamfa@j{0}%
  \ifnum`#1=43 \def\salamfa@j{2}\fi
  \ifnum`#1=45 \def\salamfa@j{2}\fi
  \ifnum`#1=42 \def\salamfa@j{2}\fi
  \ifnum`#1=47 \def\salamfa@j{2}\fi
  \ifnum`#1=37 \def\salamfa@j{2}\fi
  \ifnum`#1=94 \def\salamfa@j{2}\fi
  \ifnum`#1=41 \def\salamfa@j{5}\fi
  \ifnum`#1=93 \def\salamfa@j{5}\fi
  \ifsalamfa@prevP\else
    \ifnum`#1=40 \def\salamfa@j{4}\fi
    \ifnum`#1=91 \def\salamfa@j{4}\fi
  \fi}
\def\salamfa@collect{%
  \global\salamfa@N\z@ \global\salamfa@prevPfalse
  \def\salamfa@iW##1##2##3{%
    \if##2P\salamfa@store{\salamfa@iW{##1}{##2}{##3}}{0}\global\salamfa@prevPtrue
    \else \salamfa@store{\salamfa@iW{##1}{##2}{##3}}{1}\global\salamfa@prevPfalse\fi}%
  \def\salamfa@iO##1##2{\salamfa@jof{##1}%
    \salamfa@store{\salamfa@iO{##1}{##2}}{\salamfa@j}\global\salamfa@prevPfalse}%
  \def\salamfa@iS{\salamfa@store{\salamfa@iS}{3}\global\salamfa@prevPfalse}%
  \def\salamfa@iT##1{\salamfa@store{\salamfa@iT{##1}}{0}\global\salamfa@prevPfalse}%
  \def\salamfa@iC##1##2{\salamfa@store{\salamfa@iC{##1}{##2}}{0}\global\salamfa@prevPfalse}%
  \salamfa@items}

% find the maximal joinable runs, then cut each one down to islands
\def\salamfa@markislands{%
  \salamfa@n\@ne \salamfa@runloop}
\def\salamfa@runloop{%
  \ifnum\salamfa@n>\salamfa@N\else
    \ifnum\salamfa@J\salamfa@n>\z@
      \@tempcnta\salamfa@n
      \salamfa@runend
      \edef\salamfa@t{\noexpand\salamfa@proc{\the\@tempcnta}{\the\numexpr\salamfa@n-1\relax}}%
      \salamfa@t
    \else \advance\salamfa@n\@ne \fi
    \expandafter\salamfa@runloop
  \fi}
\def\salamfa@runend{%
  \salamfa@gofalse
  \ifnum\salamfa@n>\salamfa@N\else
    \ifnum\salamfa@J\salamfa@n>\z@ \salamfa@gotrue\fi
  \fi
  \ifsalamfa@go \advance\salamfa@n\@ne \expandafter\salamfa@runend\fi}

% slots #1..#2: trim the edges, split at an unmatched bracket, and mark what
% is left as an island if it holds an operator
\newcount\salamfa@kk
\newcount\salamfa@depth
\newcount\salamfa@cut
\def\salamfa@proc#1#2{%
  \begingroup
  \@tempcnta#1\relax \@tempcntb#2\relax
  \loop
    \salamfa@gofalse
    \ifnum\@tempcnta>\@tempcntb\else
      \ifcase\salamfa@J\@tempcnta \or\or \salamfa@gotrue\or \salamfa@gotrue\or\or
        \salamfa@gotrue\fi
    \fi
    \ifsalamfa@go \advance\@tempcnta\@ne \repeat
  \loop
    \salamfa@gofalse
    \ifnum\@tempcnta>\@tempcntb\else
      \ifcase\salamfa@J\@tempcntb \or\or \salamfa@gotrue\or \salamfa@gotrue\or
        \salamfa@gotrue\fi
    \fi
    \ifsalamfa@go \advance\@tempcntb\m@ne \repeat
  \ifnum\@tempcnta>\@tempcntb
    \endgroup
  \else
    \salamfa@cut\z@ \salamfa@hasopfalse
    % a closing bracket with no opener, scanning forward
    \salamfa@depth\z@ \salamfa@kk\@tempcnta
    \loop
      \ifnum\salamfa@J\salamfa@kk=2 \salamfa@hasoptrue\fi
      \ifnum\salamfa@J\salamfa@kk=4 \advance\salamfa@depth\@ne\fi
      \ifnum\salamfa@J\salamfa@kk=5
        \ifnum\salamfa@depth=\z@
          \ifnum\salamfa@cut=\z@ \salamfa@cut\salamfa@kk\fi
        \else \advance\salamfa@depth\m@ne\fi
      \fi
      \advance\salamfa@kk\@ne
      \ifnum\salamfa@kk>\@tempcntb\else \repeat
    % otherwise an opening bracket with no closer, scanning backward
    \ifnum\salamfa@cut=\z@
      \salamfa@depth\z@ \salamfa@kk\@tempcntb
      \loop
        \ifnum\salamfa@J\salamfa@kk=5 \advance\salamfa@depth\@ne\fi
        \ifnum\salamfa@J\salamfa@kk=4
          \ifnum\salamfa@depth=\z@
            \ifnum\salamfa@cut=\z@ \salamfa@cut\salamfa@kk\fi
          \else \advance\salamfa@depth\m@ne\fi
        \fi
        \advance\salamfa@kk\m@ne
        \ifnum\salamfa@kk<\@tempcnta\else \repeat
    \fi
    \ifnum\salamfa@cut=\z@
      \ifsalamfa@hasop
        \expandafter\xdef\csname salamfa@I\the\@tempcnta\endcsname{\the\@tempcntb}%
      \fi
      \endgroup
    \else
      \edef\salamfa@t{%
        \noexpand\salamfa@proc{\the\@tempcnta}{\the\numexpr\salamfa@cut-1\relax}%
        \noexpand\salamfa@proc{\the\numexpr\salamfa@cut+1\relax}{\the\@tempcntb}}%
      \expandafter\endgroup\salamfa@t
    \fi
  \fi}

% lay the slots out with the current handlers, an island as a single run
\def\salamfa@rloop{%
  \ifnum\salamfa@n>\salamfa@N\else
    \expandafter\ifx\csname salamfa@I\the\salamfa@n\endcsname\relax
      \csname salamfa@A\the\salamfa@n\endcsname
      \advance\salamfa@n\@ne
    \else
      \edef\salamfa@e{\csname salamfa@I\the\salamfa@n\endcsname}%
      \begingroup
        \gdef\salamfa@islx{}\salamfa@hLtr\salamfa@islx
        \salamfa@iloop
      \endgroup
      \salamfa@putisland
      \salamfa@n\salamfa@e\relax \advance\salamfa@n\@ne
    \fi
    \expandafter\salamfa@rloop
  \fi}
\def\salamfa@iloop{%
  \ifnum\salamfa@n>\salamfa@e\else
    \csname salamfa@A\the\salamfa@n\endcsname
    \advance\salamfa@n\@ne
    \expandafter\salamfa@iloop
  \fi}

\newcommand{\salamfa@id}[1]{#1}
\newcommand{\salamfaNM}[1]{\textcolor{nm}{#1}}

\newcommand{\salamfa@procline}[1]{%
  \global\advance\salamfa@nlines\@ne
  \begingroup
  \salamfa@lead\z@ \salamfa@spc\z@ \salamfa@startedfalse
  \salamfa@escfalse \salamfa@dotpfalse \salamfa@slashpfalse \salamfa@starpfalse
  \salamfa@lastsigfalse \salamfa@signc\z@ \salamfa@wreset
  \global\salamfa@linePfalse
  \gdef\salamfa@items{}\gdef\salamfa@vis{}\gdef\salamfa@chars{}%
  \ifnum\salamfa@mode=\thr@@
    \salamfa@startedtrue \gdef\salamfa@sub{}\gdef\salamfa@cmark{}%
  \else \global\salamfa@mode\z@ \fi
  \edef\salamfa@line{\detokenize{#1}\space\salamfa@end\space}%
  \expandafter\salamfa@split\salamfa@line
  \expandafter\salamfa@lex\salamfa@chars\salamfa@END
  % close whatever the end of the line leaves open
  \ifsalamfa@slashp \salamfa@op/\fi
  \ifsalamfa@starp \salamfa@op*\fi
  \ifsalamfa@dotp \salamfa@dotpfalse \salamfa@op.\fi
  \ifnum\salamfa@signc>\z@ \salamfa@signop \salamfa@signc\z@ \fi
  \salamfa@flushword
  \ifcase\salamfa@mode
  \or \global\salamfa@mode\z@
      \salamfa@xadd{\noexpand\salamfa@iT{\unexpanded\expandafter{\salamfa@sub}}}%
  \or \global\salamfa@mode\z@ \salamfa@emitcmt
  \or \salamfa@emitcmt
  \fi
  % lay the items out
  \ifsalamfa@out
    \salamfa@ctx\thr@@
    \ifsalamfa@lineP \salamfa@hRtext \salamfa@items \salamfa@closeisl
    \else
      \gdef\salamfa@tokb{}\salamfa@hLtr\salamfa@tokb
      \salamfa@items \salamfa@flushtok
    \fi
  \else
    \salamfa@ctx\z@
    \salamfa@collect \salamfa@markislands
    \ifsalamfa@lineP
      \salamfa@hRcode \def\salamfa@putisland{\salamfa@preI\salamfa@islx}%
    \else
      \salamfa@hLcode \def\salamfa@putisland{%
        \expandafter\salamfa@gapp\expandafter\salamfa@tokb\expandafter{\salamfa@islx}}%
    \fi
    \salamfa@n\@ne \salamfa@rloop
    \ifsalamfa@lineP\else \salamfa@flushtok \fi
  \fi
  \xdef\salamfa@code{\unexpanded\expandafter{\salamfa@code}%
    \ifnum\salamfa@nlines>\@ne \noexpand\\\fi
    \noexpand\strut
    \ifsalamfa@out\else \noexpand\salamfa@num{\the\salamfa@nlines}\fi
    \ifnum\salamfa@lead>\z@
      \noexpand\hspace*{\the\dimexpr\salamfa@lead\salamfa@unit\relax}\fi
    \unexpanded\expandafter{\salamfa@vis}}%
  \endgroup}

\renewcommand{\FancyVerbFormatLine}[1]{\salamfa@procline{#1}}
\newbox\salamfa@scratch

% The line number rides on the line itself, as a zero-width box at its right
% end that hangs out past the frame, so a line that wraps keeps its number on
% its first row and the numbers below stay in step.
\def\salamfa@num#1{%
  \LRE{\hbox to\z@{\hskip\dimexpr2\fboxsep+\fboxrule\relax
    {\normalfont\color{black}#1}\hss}}}

\newcommand{\salamfacardnumgap}{0pt}
\newcommand{\salamfacardbegin}{%
  \gdef\salamfa@code{}\global\salamfa@nlines=0
  \global\salamfa@mode\z@ \salamfa@outfalse
  \VerbatimEnvironment
  \setbox\salamfa@scratch=\vbox\bgroup\begin{Verbatim}}
\newcommand{\salamfacardend}{%
  \end{Verbatim}\egroup
  \par\noindent
  \hfuzz=30pt%
  \begin{LTR}%
    \fcolorbox{black}{bg}{%
      \begin{minipage}[t]{\dimexpr\linewidth-2\fboxsep-2\fboxrule-\salamfacardnumgap\relax}%
        \begin{RTL}\raggedleft\salamcodeFAfont\salamfa@code\end{RTL}%
      \end{minipage}%
    }%
  \end{LTR}%
  \par}
\newenvironment{salamfa}{\salamfacardbegin}{\salamfacardend}

\providecommand{\salamfaoutcolor}{black}
\providecommand{\salamfaoutpad}{2\fboxsep+2\fboxrule}
\providecommand{\salamfaoutheadfont}{\bfseries}
\providecommand{\salamfaoutheadcolor}{black}

\newcommand{\salamfaoutbegin}{%
  \gdef\salamfa@code{}\global\salamfa@nlines=0
  \global\salamfa@mode\z@ \salamfa@outtrue
  \VerbatimEnvironment
  \setbox\salamfa@scratch=\vbox\bgroup\begin{Verbatim}}
\newcommand{\salamfaoutend}{%
  \end{Verbatim}\egroup
  \par\nointerlineskip\noindent
  \begin{LTR}%
    \fcolorbox{\salamfaoutcolor}{white}{%
      \begin{minipage}[t]{\dimexpr\linewidth-(\salamfaoutpad)\relax}%
        \begin{RTL}\raggedleft\salamcodeFAfont\small
          \salamfa@code\end{RTL}%
      \end{minipage}%
    }%
  \end{LTR}%
  \par}
\newenvironment{progoutfa}{\salamfaoutbegin}{\salamfaoutend}
\newenvironment{outputfa}{%
  \par\bigskip\noindent{\salamfaoutheadfont\small\color{\salamfaoutheadcolor}خروجی برنامه}\par\nobreak\vskip1pt
  \salamfaoutbegin}{\salamfaoutend}

\makeatother


\renewcommand{\salamfacardnumgap}{(\fboxsep+3em)}
\renewcommand{\salamfaoutcolor}{sumFrame}
\renewcommand{\salamfaoutpad}{3\fboxsep+2\fboxrule+3em}
\renewcommand{\salamfaoutheadfont}{\headfont}
\renewcommand{\salamfaoutheadcolor}{labelColor}

% Captions and output headings travel with their box. A box's contents are
% built unframed first and measured: if the whole box fits, it is placed
% with whatever heads it; if not, a long box is cut at a line boundary so its
% first part (with the caption) fills this page and the rest continues in its
% own frame on the next; a short box moves to the next page whole.
\makeatletter
\newbox\salambook@box
\newbox\salambook@inner
\newbox\salambook@part
\newdimen\salambook@avail
\newdimen\salambook@headh
\def\salambook@head{}
\def\salambook@framer{}
\def\salambook@framecode#1{\fcolorbox{black}{bg}{#1}}
\def\salambook@frameout#1{\fcolorbox{\salamfaoutcolor}{white}{#1}}
\def\salambook@framed#1{% #1: a box register, used up
  \setbox\salambook@box=\vbox{\hsize\linewidth\parindent\z@\hfuzz=30pt
    \noindent\begin{LTR}\salambook@framer{\box#1}\end{LTR}\par}}
\def\salambook@emithead{\salambook@head\global\let\salambook@head\@empty}
\def\salambook@place{% content in \salambook@inner, frame in \salambook@framer
  \par\penalty\@M % let the page builder catch up before measuring
  \salambook@headh=\z@
  \ifx\salambook@head\@empty\else\salambook@headh=3\baselineskip\fi
  \setbox\salambook@part=\copy\salambook@inner
  \salambook@framed\salambook@part
  \ifdim\pagegoal=\maxdimen \salambook@avail=\maxdimen
  \else \salambook@avail=\dimexpr\pagegoal-\pagetotal\relax \fi
  \ifdim\dimexpr\ht\salambook@box+\dp\salambook@box+\salambook@headh\relax<\salambook@avail
    \salambook@emithead\noindent\box\salambook@box\par
  \else
    \ifdim\ht\salambook@inner>8\baselineskip
      \ifdim\salambook@avail>\dimexpr\salambook@headh+4\baselineskip+2\fboxsep+2\fboxrule+6pt\relax
        \salambook@split
      \else \salambook@whole \fi
    \else \salambook@whole \fi
  \fi
  \addvspace{\medskipamount}}
\def\salambook@whole{%
  \Needspace{\dimexpr\ht\salambook@box+\dp\salambook@box+\salambook@headh\relax}%
  \salambook@emithead\noindent\box\salambook@box\par}
\def\salambook@split{%
  \splittopskip=\ht\strutbox \splitmaxdepth=\dp\strutbox \vbadness=\@M
  \dimen@=\dimexpr\salambook@avail-\salambook@headh-2\fboxsep-2\fboxrule-1.5\baselineskip\relax
  \ifdim\dimen@>\dimexpr\ht\salambook@inner-3\baselineskip\relax
    \dimen@=\dimexpr\ht\salambook@inner-3\baselineskip\relax\fi
  \setbox\salambook@part=\vsplit\salambook@inner to \dimen@
  \setbox\salambook@part=\vbox{\unvbox\salambook@part}%
  \salambook@framed\salambook@part
  \salambook@emithead\noindent\box\salambook@box\par
  \setbox\salambook@inner=\vbox{\unvbox\salambook@inner}%
  \salambook@framed\salambook@inner
  \noindent\box\salambook@box\par}
\def\salambook@content#1#2{% #1 width, #2 extra font setup
  \setbox\salambook@inner=\vbox{\hsize#1\parindent\z@
    \begin{RTL}\raggedleft\salamcodeFAfont#2\salamfa@code\end{RTL}}}
\renewcommand{\salamfacardend}{%
  \end{Verbatim}\egroup
  \par
  \salambook@content{\dimexpr\linewidth-2\fboxsep-2\fboxrule-\salamfacardnumgap\relax}{}%
  \let\salambook@framer\salambook@framecode
  \salambook@place}
\renewcommand{\salamfaoutend}{%
  \end{Verbatim}\egroup
  \par
  \salambook@content{\dimexpr\linewidth-(\salamfaoutpad)\relax}{\small}%
  \let\salambook@framer\salambook@frameout
  \salambook@place}
\renewenvironment{outputfa}{%
  \par\bigskip
  \gdef\salambook@head{\noindent{\salamfaoutheadfont\small\color{\salamfaoutheadcolor}خروجی برنامه}\par\nobreak\vskip1pt}%
  \salamfaoutbegin}{\salamfaoutend}
% Inline code with operators, brackets or decimals: laid out by the same
% engine as the listings (so it reads the same way), kept on one line, and
% set in the plain inline-code style.
\newcommand{\scode}[1]{\mbox{\begingroup
  \def\%{\@percentchar}\def\textbackslash{\@backslashchar}\edef\salambook@sc{#1}%
  \gdef\salamfa@code{}\global\salamfa@nlines\z@ \global\salamfa@mode\z@ \salamfa@outfalse
  \def\salamfa@num##1{}\let\strut\relax
  \let\textcolor\@secondoftwo \def\color##1{}%
  \def\kw##1{##1}\def\tyword##1{##1}\def\salamfaNM##1{##1}%
  \def\stword##1{##1}\def\cmtword##1{##1}\def\opword##1{##1}%
  \expandafter\salamfa@procline\expandafter{\salambook@sc}%
  \salamcodeFAfont\small\salamfa@code
  \endgroup}}
% A solution heading in appendix C is emitted with the listing it heads.
\newcommand{\solutionhead}[1]{\par\gdef\salambook@head{\subsection*{#1}}}
% A hand-built output box (the chapter 10 calendar) in the standard frame.
\newcommand{\salamoutbox}[1]{\par\bigskip
  \gdef\salambook@head{\noindent{\salamfaoutheadfont\small\color{\salamfaoutheadcolor}خروجی برنامه}\par\nobreak\vskip1pt}%
  \setbox\salambook@inner=\vbox{\hsize\dimexpr\linewidth-(\salamfaoutpad)\relax\parindent\z@
    \begin{RTL}\raggedleft\salamcodeFAfont\small #1\end{RTL}}%
  \let\salambook@framer\salambook@frameout
  \salambook@place}
\let\salambook@section\section
\renewcommand{\section}{\par\Needspace{6\baselineskip}\salambook@section}
\let\salambook@subsection\subsection
\renewcommand{\subsection}{\par\Needspace{6\baselineskip}\salambook@subsection}
\makeatother

\hfuzz=18pt
\newcommand{\fcode}[1]{{\salamcodefont\small\spaceskip=0.33em plus 0.1em minus 0.05em #1}}
\newcommand{\code}[1]{\lr{{\latincodefont\small #1}}}
\newcommand{\kwd}[1]{{\salamcodefont\small\bfseries\color{salamKeyword}#1}}

\newcounter{program}[chapter]
\renewcommand{\theprogram}{\thechapter.\arabic{program}}
\makeatletter
\newcommand{\program}[1]{%
  \par\medskip\refstepcounter{program}%
  \xdef\salamprogramnum{\theprogram}%
  \gdef\salambook@head{\noindent{\headfont\small\color{labelColor}برنامه\enskip\salamprogramnum}%
    {\small\ \ #1}\par\nobreak\vskip2pt}}
\makeatother

\newenvironment{salamadm}[3]{%
  \par\smallskip\noindent
  \def\FrameCommand{{\color{#2}\vrule width 3pt}\fboxsep=8pt\colorbox{#1}}%
  \MakeFramed{\hsize=\dimexpr\linewidth-3pt\relax\linewidth=\hsize\FrameRestore}%
  \hfuzz=7pt%
  \noindent{\headfont\small\color{#2}#3}\par\nobreak\vskip3pt\noindent\ignorespaces
}{%
  \par\endMakeFramed\smallskip
}

\newenvironment{note}[1][نکته]{\salamadm{noteBg}{noteFrame}{#1}}{\endsalamadm}
\newenvironment{tip}[1][پیشنهاد]{\salamadm{tipBg}{tipFrame}{#1}}{\endsalamadm}
\newenvironment{warn}[1][مراقب باشید]{\salamadm{warnBg}{warnFrame}{#1}}{\endsalamadm}
\newenvironment{deep}[1][بیشتر بدانید]{\salamadm{deepBg}{deepFrame}{#1}}{\endsalamadm}
\newenvironment{example}[1][مثال]{\salamadm{exBg}{exFrame}{#1}}{\endsalamadm}
\newenvironment{summary}[1][چکیده‌ی فصل]{\salamadm{sumBg}{sumFrame}{#1}}{\endsalamadm}
\newenvironment{goals}[1][در این فصل چه می‌آموزید؟]{\salamadm{brandLight}{brand}{#1}}{\endsalamadm}

\newcounter{exercise}[chapter]
\renewcommand{\theexercise}{\thechapter.\arabic{exercise}}
\newcommand{\exercise}{\par\medskip\refstepcounter{exercise}%
  \noindent{\headfont\small\color{exFrame}تمرین \theexercise}\ \ \ignorespaces}

\titleformat{\chapter}[display]
  {\normalfont\titlefont\color{brand}}
  {\filleft\fontsize{60}{60}\selectfont\thechapter}
  {8pt}
  {\filleft\fontsize{24}{30}\selectfont}
  [\vspace{4pt}{\color{brand}\titlerule[1.2pt]}\vspace{6pt}]
\titlespacing*{\chapter}{0pt}{20pt}{28pt}

\titleformat{\section}
  {\normalfont\headfont\large\bfseries\color{ink}}{\thesection}{0.8em}{}
\titleformat{\subsection}
  {\normalfont\headfont\normalsize\bfseries\color{ink}}{\thesubsection}{0.8em}{}

\makeatletter
\newcommand{\beginappendices}{\clearpage
  \setcounter{chapter}{0}\setcounter{section}{0}%
  \gdef\@chapapp{\appendixname}%
  \renewcommand{\thechapter}{\Alph{chapter}}%
  \renewcommand{\theHchapter}{app.\arabic{chapter}}%
  \renewcommand{\theHsection}{app.\arabic{chapter}.\arabic{section}}}
\makeatother

\setlength{\headheight}{18pt}
\pagestyle{fancy}
\fancyhf{}
\fancyhead[LO,RE]{\small\headfont\rightmark}
\fancyhead[RO,LE]{\small\headfont\leftmark}
\fancyfoot[C]{\thepage}
\renewcommand{\headrulewidth}{0.4pt}

\hypersetup{
  colorlinks=true,
  hypertexnames=false,
  linkcolor=brand,
  urlcolor=salamString,
  citecolor=salamType,
  pdftitle={آموزش برنامه‌نویسی با زبان سلام},
  pdfauthor={سیدعلی محمدیه، مهدی اسماعیلی},
  bookmarksnumbered=true,
}

\setlist{itemsep=2pt,topsep=4pt}
\linespread{1.2}
\emergencystretch=3em
\raggedbottom

\begin{document}

\frontmatter
\begin{titlepage}
\centering
\vspace*{1.5cm}
{\titlefont\color{brand}\fontsize{34}{44}\selectfont آموزش برنامه‌نویسی\\[6pt] با زبان سلام\par}
\vspace{0.9cm}
{\headfont\large از نخستین برنامه تا نوشتن برنامه‌های کوچک و کاربردی\par}
\vspace{1.6cm}
{\color{brand}\rule{0.55\textwidth}{1pt}\par}
\vspace{1.2cm}
{\large\bfseries سیدعلی محمدیه\par}
{\normalsize طراح زبان برنامه‌نویسی سلام\par}
\vspace{0.7cm}
{\large\bfseries دکتر مهدی اسماعیلی\par}
{\normalsize عضو هیئت علمی دانشگاه آزاد اسلامی واحد کاشان\par}
\vspace{1.2cm}
{\color{brand}\rule{0.55\textwidth}{1pt}\par}
\vfill
{\headfont\large انتشارات دانشگاه آزاد اسلامی\par}
\vspace{0.3cm}
{\normalsize چاپ نخست\par}
\end{titlepage}

\thispagestyle{empty}
\vspace*{1cm}
\noindent
\begin{tabular}{@{}p{3cm}p{10cm}@{}}
عنوان: & آموزش برنامه‌نویسی با زبان سلام \\[3pt]
نویسندگان: & سیدعلی محمدیه، مهدی اسماعیلی \\[3pt]
ناشر: & انتشارات دانشگاه آزاد اسلامی \\[3pt]
نوبت چاپ: & نخست \\[3pt]
شمارگان: & ۱٬۰۰۰ نسخه \\[3pt]
قطع: & وزیری \\[3pt]
\end{tabular}

\vspace{1.5cm}
\noindent\small
این کتاب نخستین کتاب رسمی آموزش زبان برنامه‌نویسی سلام است. همه‌ی برنامه‌های
آن با نسخه‌ی رسمی زبان آزموده شده‌اند و می‌توانید بی‌آنکه چیزی روی رایانه
نصب کنید، آن‌ها را در ویرایشگر برخط سلام اجرا کنید:
\begin{center}
\url{https://editor.salamlang.ir}
\end{center}
\noindent حق چاپ و نشر این اثر برای ناشر محفوظ است. نمونه‌کدهای کتاب را
می‌توانید آزادانه در یادگیری و آموزش به کار ببرید.
\normalsize
\clearpage

\chapter*{پیشگفتار}
\addcontentsline{toc}{chapter}{پیشگفتار}
\markboth{پیشگفتار}{پیشگفتار}

برنامه‌نویسی یعنی گفت‌وگو با رایانه. شما چیزی می‌خواهید، آن را با زبانی
که رایانه می‌فهمد می‌نویسید و رایانه، بی‌چون‌وچرا و با سرعتی باورنکردنی،
همان را انجام می‌دهد. سال‌هاست که این گفت‌وگو تنها به زبان انگلیسی ممکن
بوده است؛ واژه‌هایی مانند \code{if} و \code{while} و \code{print} برای
میلیون‌ها نفر نخستین سد راه یادگیری بوده‌اند، نه به این خاطر که
برنامه‌نویسی دشوار است، بلکه به این خاطر که پیش از رسیدن به اندیشه‌ی
برنامه‌نویسی، باید از دیوار یک زبان بیگانه بالا می‌رفتند.

زبان برنامه‌نویسی سلام برای برداشتن همین دیوار ساخته شده است. در سلام
می‌نویسید \fcode{اگر}، می‌نویسید \fcode{تکرار}، می‌نویسید \fcode{سرچاپ}؛ نام
متغیرها و تابع‌هایتان را به فارسی می‌گذارید و برنامه‌ای که می‌نویسید مانند
یک متن فارسی خوانده می‌شود. سلام یک زبان واقعی و کامل است، نه یک اسباب‌بازی
آموزشی: با آن می‌توان برنامه‌های بزرگ نوشت، صفحه‌ها و وب‌سایت‌ها ساخت و
حتی خود مترجم زبان سلام نیز با زبان سلام نوشته شده است. سلام املای انگلیسی
هم دارد و برنامه‌ای که به فارسی می‌نویسید، همان قدرت و همان سرعت را
دارد که اگر به انگلیسی می‌نوشتید. اما این کتاب درباره‌ی سلام فارسی است؛ از
این‌جا به بعد، همه چیز به زبان خودمان خواهد بود.

\section*{این کتاب برای چه کسی است؟}

برای هر کسی که می‌خواهد برنامه‌نویسی را از صفر بیاموزد. هیچ پیش‌نیازی لازم
نیست: نه آشنایی قبلی با برنامه‌نویسی، نه ریاضیات پیشرفته و نه زبان
انگلیسی. تنها چیزی که لازم دارید یک رایانه یا حتی یک گوشی با مرورگر است و
کمی حوصله. دانش‌آموزان دبیرستان، دانشجویان سال اول، معلمان، و هر بزرگسال
کنجکاوی که همیشه دلش می‌خواسته بداند «برنامه‌نویسی یعنی چه» مخاطب این کتاب
هستند. اگر پیش‌تر با زبان دیگری برنامه نوشته‌اید هم می‌توانید این کتاب را
بخوانید، اما احتمالاً بخش‌های نخست را سریع‌تر پشت سر خواهید گذاشت.

\section*{چیزی روی رایانه نصب نکنید}

یکی از تصمیم‌های مهم ما در نوشتن این کتاب این بود که خواننده هیچ‌گاه مجبور
نشود با «خط فرمان»، «ترمینال» یا نصب نرم‌افزار درگیر شود. همه‌ی برنامه‌های
کتاب را می‌توانید در ویرایشگر برخط سلام بنویسید و اجرا کنید:
\begin{center}
\url{https://editor.salamlang.ir}
\end{center}
این ویرایشگر در مرورگر شما اجرا می‌شود، به هیچ سرویس دیگری وابسته نیست و
همان چیزی است که در فصل نخست با آن آشنا می‌شوید. کافی است نشانی را باز
کنید، زبان فارسی را برگزینید و بنویسید.

\section*{در این کتاب چه چیزهایی می‌آموزید؟}

مسیر این کتاب همان مسیر آشنای همه‌ی کتاب‌های خوب برنامه‌نویسی است، با این
تفاوت که در هر گام کد فارسی می‌نویسید:
\begin{itemize}
  \item در فصل~\ref{ch:first} نخستین برنامه‌ی خود را می‌نویسید و یاد
        می‌گیرید چگونه چیزی را روی صفحه چاپ کنید.
  \item فصل‌های~\ref{ch:variables} تا~\ref{ch:operators} به متغیر، حافظه،
        گونه‌های داده و محاسبه می‌پردازند؛ یعنی مصالح ساختمانی هر برنامه.
  \item در فصل~\ref{ch:conditions} برنامه‌ها یاد می‌گیرند تصمیم بگیرند و
        در فصل~\ref{ch:loops} یاد می‌گیرند کاری را بارها تکرار کنند.
  \item فصل~\ref{ch:functions} به شما نشان می‌دهد چگونه برنامه را به
        قطعه‌های کوچک و قابل استفاده‌ی دوباره تقسیم کنید.
  \item فصل‌های~\ref{ch:arrays} و~\ref{ch:strings} درباره‌ی نگه‌داری
        فهرستی از داده‌ها و کار با متن هستند.
  \item و فصل~\ref{ch:projects} همه‌ی این‌ها را در چند برنامه‌ی کوچک و
        کاربردی کنار هم می‌گذارد.
\end{itemize}
ما عمداً از الگوریتم‌های مشهور و برنامه‌های طولانی پرهیز کرده‌ایم. هدف این
کتاب این است که اصول پایه را چنان خوب یاد بگیرید که بعد از آن هر چیزی را
بتوانید خودتان بیاموزید. مثال‌ها همه ساده و ملموس‌اند: محاسبه‌ی معدل، صورت‌حساب
خرید، جدول تبدیل دما، شمارش معکوس و چیزهایی از این دست.

\section*{چگونه این کتاب را بخوانیم؟}

برنامه‌نویسی با خواندن یاد گرفته نمی‌شود؛ با نوشتن یاد گرفته می‌شود. از شما
می‌خواهیم هر برنامه‌ای را که در کتاب می‌بینید، خودتان در ویرایشگر تایپ کنید
(نه اینکه کپی کنید)، اجرایش کنید و خروجی را با آنچه در کتاب آمده مقایسه
کنید. سپس آن را دست‌کاری کنید: عددی را عوض کنید، خطی را حذف کنید، چیزی
اضافه کنید و ببینید چه می‌شود. بهترین معلم شما همین آزمایش‌کردن است.

در سراسر کتاب چند نوع جعبه می‌بینید. جعبه‌های \textbf{نکته} چیزی را روشن‌تر
می‌کنند، جعبه‌های \textbf{پیشنهاد} راه بهتری برای انجام کارها نشان می‌دهند،
جعبه‌های \textbf{مراقب باشید} از اشتباه‌های رایج خبر می‌دهند و جعبه‌های
\textbf{بیشتر بدانید} برای خوانندگان کنجکاو نوشته شده‌اند و می‌توانید در
خواندن نخست از آن‌ها بگذرید. پایان هر فصل تمرین‌هایی دارد؛ پاسخ شماری از
آن‌ها را در پیوست پایانی کتاب آورده‌ایم، اما خواهش می‌کنیم پیش از نگاه
کردن به پاسخ، خودتان تلاش کنید.

\section*{سخن نویسندگان}

این کتاب حاصل همکاری خالق زبان سلام و یک مدرس دانشگاه است؛ یکی از درون
زبان به آن نگاه کرده و دیگری از نگاه کلاس درس و دانشجویی که برای نخستین بار
با برنامه‌نویسی روبه‌رو می‌شود. امیدواریم این دو نگاه در کنار هم کتابی
ساخته باشد که هم درست باشد و هم مهربان. اگر در خواندن کتاب به اشکالی
برخوردید یا پیشنهادی داشتید، خوشحال می‌شویم آن را از راه صفحه‌ی پروژه‌ی
سلام در گیت‌هاب با ما در میان بگذارید.

\begin{flushleft}
سیدعلی محمدیه\\
مهدی اسماعیلی
\end{flushleft}
\clearpage


\tableofcontents
\clearpage

\mainmatter
\chapter{نخستین برنامه}
\label{ch:first}

\begin{goals}
\begin{itemize}
  \item برنامه چیست و رایانه چگونه آن را اجرا می‌کند
  \item آشنایی با ویرایشگر برخط سلام و اجرای نخستین برنامه
  \item چاپ کردن متن و عدد روی صفحه با \fcode{سرچاپ} و \fcode{چاپ}
  \item نوشتن توضیح در میان برنامه
  \item خواندن پیام خطا و رفع ساده‌ترین اشتباه‌ها
\end{itemize}
\end{goals}

\section{برنامه چیست؟}

فرض کنید می‌خواهید به دوستی که هرگز چای دم نکرده است، یاد بدهید چای دم
کند. نمی‌توانید بگویید «یک چای خوب درست کن»؛ باید قدم‌به‌قدم بگویید: کتری
را پر از آب کن، روی اجاق بگذار، صبر کن تا بجوشد، چای خشک را در قوری بریز
و همین‌طور تا آخر. رایانه هم درست چنین دوستی است: بسیار سریع، بسیار دقیق
و بسیار کم‌هوش. هر کاری را که به او بگویید انجام می‌دهد، اما فقط به شرطی که
دقیقاً بگویید چه کند.

\textbf{برنامه} همین فهرست دستورهای قدم‌به‌قدم است، و \textbf{برنامه‌نویسی}
هنر نوشتن آن. زبانی که این دستورها را با آن می‌نویسیم
\textbf{زبان برنامه‌نویسی} نام دارد. زبان برنامه‌نویسی، برخلاف زبان انسان‌ها،
واژه‌های اندکی دارد و قاعده‌هایش بسیار سخت‌گیرانه است؛ اما در عوض هیچ
ابهامی در آن نیست. زبان این کتاب، \textbf{سلام} است: زبانی که دستورهایش به
فارسی نوشته می‌شود.

\begin{note}
وقتی برنامه‌ای می‌نویسید، متنی که نوشته‌اید را \textbf{کد} یا
\textbf{کد منبع} می‌نامند. رایانه این متن را می‌خواند، دستورهایش را یکی‌یکی
انجام می‌دهد و نتیجه را نشان می‌دهد. به این کار \textbf{اجرا کردن} برنامه
می‌گویند.
\end{note}

\section{ویرایشگر برخط سلام}

برای نوشتن و اجرا کردن برنامه‌های این کتاب به هیچ نصب و تنظیمی نیاز
ندارید. مرورگر خود را باز کنید و به این نشانی بروید:
\begin{center}
\url{https://editor.salamlang.ir}
\end{center}
صفحه‌ای که می‌بینید «زمین بازی سلام» نام دارد و از سه بخش اصلی ساخته شده
است:
\begin{enumerate}
  \item \textbf{نوار بالا:} در این نوار می‌توانید زبان را انتخاب کنید. گزینه‌ی
        \textbf{فارسی} را برگزینید تا هم ظاهر صفحه فارسی شود و هم ویرایشگر
        آماده‌ی دریافت کد فارسی. در همین نوار دو حالت \textbf{برنامه} و
        \textbf{چیدمان} وجود دارد؛ ما در سراسر این کتاب در حالت
        \textbf{برنامه} کار می‌کنیم. دکمه‌ی \textbf{اجرا} هم همین‌جاست.
  \item \textbf{ویرایشگر:} پنجره‌ای که کد را در آن می‌نویسید. شماره‌ی خط‌ها
        در کنار آن دیده می‌شود و واژه‌های مهم زبان رنگی می‌شوند.
  \item \textbf{خروجی:} پنجره‌ای که نتیجه‌ی اجرای برنامه در آن نمایش داده
        می‌شود. اگر برنامه اشتباهی داشته باشد، پیام خطا هم همین‌جا ظاهر
        می‌شود.
\end{enumerate}

\begin{tip}
در نوار بالای ویرایشگر فهرستی از برنامه‌های نمونه هم هست. بد نیست همین حالا
یکی از آن‌ها را باز کنید و دکمه‌ی اجرا را بزنید تا ببینید همه چیز درست کار
می‌کند. نگران نباشید اگر هنوز چیزی از کد سر در نمی‌آورید؛ تا پایان این
کتاب همه‌ی آن‌ها را خواهید فهمید.
\end{tip}

\section{سلام، دنیا!}

سنتی قدیمی میان برنامه‌نویسان هست: نخستین برنامه در هر زبان تازه، برنامه‌ای
است که جمله‌ی «سلام، دنیا!» را روی صفحه می‌نویسد. ما هم همین سنت را نگه
می‌داریم. متن ویرایشگر را پاک کنید و این سه خط را در آن بنویسید:

\program{سلام، دنیا!}
\begin{salamfa}
روال ریشه:
    سرچاپ "سلام، دنیا!"
پایان
\end{salamfa}

حالا دکمه‌ی \textbf{اجرا} را بزنید. (اگر دکمه‌ی \textbf{خودکار} در نوار بالا روشن
باشد، که به طور پیش‌فرض روشن است، ویرایشگر همراه با نوشتن شما برنامه را خودش
اجرا می‌کند.) در پنجره‌ی خروجی این را می‌بینید:

\begin{outputfa}
سلام، دنیا!
\end{outputfa}

تبریک! شما همین حالا نخستین برنامه‌ی خود را نوشتید و اجرا کردید. حالا
ببینیم این سه خط چه می‌گویند.

\subsection{روال ریشه: نقطه‌ی شروع}

خط نخست، \scode{روال ریشه:}، به سلام می‌گوید: «اجرای برنامه از این‌جا شروع
می‌شود.» هر برنامه‌ی سلام باید دقیقاً یک \fcode{روال ریشه} داشته باشد.
واژه‌ی \fcode{پایان} در خط آخر پایان آن را نشان می‌دهد. فعلاً این دو خط را
مانند جلد یک دفتر ببینید: هر چه می‌نویسیم، میان این دو خط می‌نویسیم. در فصل
\ref{ch:functions} خواهیم دید که واژه‌ی «روال» چیزهای بیشتری در چنته دارد.

به دونقطه‌ی پایان خط نخست دقت کنید. در سلام هر جا که یک بلوک (یک دسته
دستور که با هم هستند) آغاز می‌شود، خط با \code{:} تمام می‌شود و بلوک با
\fcode{پایان} بسته می‌شود. این الگو را بارها و بارها خواهید دید.

\subsection{تورفتگی: نظم دیداری برنامه}

دقت کنید که خط میانی کمی جلوتر از دو خط دیگر نوشته شده است. به این
جلوآمدگی \textbf{تورفتگی} می‌گویند و نشان می‌دهد این خط \emph{درون} روال ریشه است. رسم بر این است که هر چه درون یک بلوک است را با چهار فاصله تو
ببریم. سلام برای اجرای برنامه به این فاصله‌ها نیازی ندارد و اگر آن‌ها را
ننویسید هم برنامه اجرا می‌شود؛ اما برنامه‌ای که تورفتگی درستی ندارد، مثل
کتابی است که پاراگراف‌بندی ندارد: می‌شود خواند، اما سخت.

\subsection{چاپ: نوشتن روی صفحه}

خط میانی همان جایی است که کار انجام می‌شود. دستور \kwd{سرچاپ} هر چه را که
جلویش بنویسید روی پنجره‌ی خروجی می‌نویسد و سپس به خط بعد می‌رود. متنی که
می‌خواهیم چاپ شود را میان دو علامت نقل قول \code{"} می‌گذاریم. به این متن
داخل نقل قول \textbf{رشته} می‌گویند؛ در فصل \ref{ch:strings} بسیار با آن
سروکار خواهیم داشت. خود علامت‌های نقل قول چاپ نمی‌شوند؛ آن‌ها فقط به سلام
می‌گویند متن از کجا شروع می‌شود و کجا پایان می‌یابد.

\begin{tip}
علامت استاندارد نقل قول \code{"} است که در صفحه‌کلید معمولاً با نگه‌داشتن
شیفت و زدن کلید \code{'} نوشته می‌شود. اگر تایپ آن در متن فارسی دست‌وپاگیر
بود، می‌توانید به جای آن از گیومه‌ی فارسی \code{«...»} هم استفاده کنید:
\fcode{سرچاپ «سلام، دنیا!»} درست همان کاری را می‌کند که
\fcode{سرچاپ "سلام، دنیا!"} می‌کند. با این‌حال \code{"} رایج‌تر و استاندارد
است و در این کتاب همه‌جا همان را به کار می‌بریم. البته می‌توانید \emph{درون}
رشته هر علامتی، از جمله گیومه‌ی فارسی، بنویسید: \scode{سرچاپ "او گفت: «سلام»"}
کاملاً درست است.
\end{tip}

\section{چند دستور پشت سر هم}

یک برنامه به ندرت فقط یک دستور دارد. دستورها را زیر هم می‌نویسیم و سلام
آن‌ها را به همان ترتیب، از بالا به پایین، اجرا می‌کند:

\program{چند خط چاپ}
\begin{salamfa}
روال ریشه:
    سرچاپ "به کتاب آموزش سلام خوش آمدید."
    سرچاپ "این دومین خط است."
    سرچاپ "و این سومین خط."
پایان
\end{salamfa}

\begin{outputfa}
به کتاب آموزش سلام خوش آمدید.
این دومین خط است.
و این سومین خط.
\end{outputfa}

هر \fcode{سرچاپ} یک خط تازه در خروجی می‌سازد. ترتیب مهم است: اگر جای دو دستور
را عوض کنید، جای دو خط در خروجی هم عوض می‌شود. همین حالا امتحان کنید.

\section{چاپ عددها و محاسبه}

\fcode{سرچاپ} فقط برای متن نیست؛ عدد را هم چاپ می‌کند. عددها را بدون نقل قول
می‌نویسیم:

\begin{salamfa}
روال ریشه:
    سرچاپ ۱۴۰۴
    سرچاپ ۲ + ۳
    سرچاپ ۱۰ * ۴
پایان
\end{salamfa}

\begin{outputfa}
1404
5
40
\end{outputfa}

به خط سوم دقت کنید: نوشتیم \scode{۲ + ۳} اما سلام \code{5} را چاپ کرد. یعنی
سلام پیش از چاپ، حساب را انجام داد. علامت \code{*} در برنامه‌نویسی همان ضرب
است (چون علامت × روی صفحه‌کلید نیست). در فصل \ref{ch:operators} به طور کامل
درباره‌ی محاسبه صحبت می‌کنیم.

حالا ببینید اگر همان \scode{۲ + ۳} را داخل نقل قول بگذاریم چه می‌شود:

\begin{salamfa}
روال ریشه:
    سرچاپ "2 + 3"
    سرچاپ ۲ + ۳
پایان
\end{salamfa}

\begin{outputfa}
2 + 3
5
\end{outputfa}

هر چه داخل نقل قول باشد \emph{متن} است و سلام به محتوایش کاری ندارد؛ آن را
همان‌طور که هست چاپ می‌کند. اما بیرون نقل قول، \scode{۲ + ۳} یک محاسبه است.
این تفاوت میان «متن» و «عدد» یکی از مهم‌ترین چیزهایی است که در این کتاب یاد
می‌گیرید.

\section{چاپ چند چیز در یک خط}

اگر بخواهیم در یک خط هم متن و هم عدد چاپ کنیم، آن‌ها را با ویرگول
انگلیسی \code{,} از هم جدا می‌کنیم. سلام آن‌ها را پشت سر هم، با یک فاصله
میانشان، چاپ می‌کند:

\program{متن و عدد در یک خط}
\begin{salamfa}
روال ریشه:
    سرچاپ "حاصل جمع ۲ و ۳ برابر است با", ۲ + ۳
    سرچاپ "سن سارا:", ۱۷, "سال"
پایان
\end{salamfa}

\begin{outputfa}
حاصل جمع ۲ و ۳ برابر است با 5
سن سارا: 17 سال
\end{outputfa}

\begin{note}
درون رشته‌ها می‌توانید از رقم‌های فارسی استفاده کنید و همان‌طور که هستند
چاپ می‌شوند. اما خود سلام نتیجه‌ی محاسبه‌ها را همیشه با رقم‌های انگلیسی
چاپ می‌کند، حتی اگر عددهای کد را با رقم فارسی نوشته باشید؛ در خروجی
بالا هم به همین دلیل «۵» به شکل انگلیسی \code{5} چاپ شده است. سلام برای
نوشتن یک عدد در کد هر دو رقم را می‌پذیرد: \scode{چاپ ۲ + ۳} و
\scode{چاپ 2 + 3} هر دو همان نتیجه را می‌دهند. در این کتاب عددهای کد را
با رقم فارسی می‌نویسیم، چون کتاب فارسی است؛ فقط یادتان باشد که خروجی
خودِ سلام همیشه انگلیسی می‌ماند.
\end{note}

\section{چاپ: نوشتن بدون رفتن به خط بعد}

دستور \kwd{چاپ} درست مثل \fcode{سرچاپ} است، با یک تفاوت: بعد از نوشتن، به
خط بعد نمی‌رود. یعنی دستور بعدی از همان‌جا ادامه می‌دهد:

\program{تفاوت چاپ و سرچاپ}
\begin{salamfa}
روال ریشه:
    چاپ "الف"
    چاپ "ب"
    چاپ "پ"
    سرچاپ ""
    سرچاپ "ت"
    سرچاپ "ث"
پایان
\end{salamfa}

\begin{outputfa}
الفبپ
ت
ث
\end{outputfa}

سه \fcode{چاپ} نخست همه در یک خط نوشتند. سپس \fcode{سرچاپ ""} یک رشته‌ی
خالی چاپ کرد؛ یعنی چیزی ننوشت اما به خط بعد رفت. \fcode{چاپ} در فصل‌های
بعد، وقتی بخواهیم چیزی را قطعه‌قطعه در یک خط بسازیم، به کارمان می‌آید. تا
آن زمان بیشتر از \fcode{سرچاپ} استفاده می‌کنیم.

\begin{tip}
\fcode{سرچاپ} را می‌توانید بدون هیچ چیز هم بنویسید. \fcode{سرچاپ} به تنهایی
یک خط خالی در خروجی می‌سازد و برای فاصله انداختن میان بخش‌های خروجی مفید
است.
\end{tip}

\section{توضیح نوشتن در برنامه}

گاهی می‌خواهیم در میان کد، برای خودمان یا دیگران یادداشتی بگذاریم:
این خط چه کار می‌کند، چرا این‌طور نوشته شده، یا چه چیزی هنوز کامل نیست.
به این یادداشت‌ها \textbf{توضیح} می‌گویند. سلام هر چه را بعد از دو خط مورب
\code{//} بیاید تا پایان همان خط نادیده می‌گیرد:

\program{توضیح در برنامه}
\begin{salamfa}
// این برنامه یک پیام خوش‌آمد چاپ می‌کند
روال ریشه:
    سرچاپ "خوش آمدید"    // چاپ پیام اصلی
    // چاپ "این خط اجرا نمی‌شود"
    سرچاپ "روز خوبی داشته باشید"
پایان
\end{salamfa}

\begin{outputfa}
خوش آمدید
روز خوبی داشته باشید
\end{outputfa}

خط چهارم را ببینید: با اینکه یک دستور \fcode{چاپ} است، چون جلویش \code{//}
گذاشته‌ایم اجرا نشده است. برنامه‌نویسان از این ترفند برای «خاموش کردن»
موقت یک خط استفاده می‌کنند، بی‌آنکه آن را پاک کنند.

اگر توضیح چند خط باشد، می‌توانید آن را میان \code{/*} و \code{*/} بنویسید:

\begin{salamfa}
/*
   این برنامه را سارا نوشته است.
   کار آن: چاپ یک شعر کوتاه.
*/
روال ریشه:
    سرچاپ "بنی آدم اعضای یکدیگرند"
    سرچاپ "که در آفرینش ز یک گوهرند"
پایان
\end{salamfa}

\begin{outputfa}
بنی آدم اعضای یکدیگرند
که در آفرینش ز یک گوهرند
\end{outputfa}

\begin{tip}
توضیح خوب چیزی را می‌گوید که از خود کد پیدا نیست: «چرا» را می‌گوید، نه
«چه» را. نوشتن \fcode{// چاپ پیام} بالای \fcode{سرچاپ "پیام"} کمکی به کسی
نمی‌کند؛ اما نوشتن \fcode{// این عدد را از جدول قیمت‌های ۱۴۰۴ برداشته‌ایم}
ارزشمند است. در این کتاب، برای آموزش، بیشتر از حد معمول توضیح می‌نویسیم.
\end{tip}

\section{وقتی اشتباه می‌کنیم}

همه اشتباه می‌کنند؛ برنامه‌نویسان باتجربه فقط سریع‌تر اشتباهشان را پیدا
می‌کنند. بیایید عمداً اشتباه کنیم تا ببینیم سلام چه می‌گوید. علامت نقل قول
پایانی را جا می‌اندازیم:

\begin{salamfa}
روال ریشه:
    سرچاپ "سلام، دنیا!
پایان
\end{salamfa}

با زدن دکمه‌ی اجرا، به جای خروجی، پیامی مانند این می‌بینید:

\begin{outputfa}
خطا: رشته‌ی پایان‌نیافته
 --> main.salam:2:11
\end{outputfa}

پیام خطا دو چیز مهم به ما می‌گوید: \textbf{چه} اشتباهی رخ داده («رشته‌ی
پایان‌نیافته»، یعنی رشته‌ای که نقل قول پایانی ندارد) و \textbf{کجا}
(خط ۲، ستون ۱۱؛ همان‌جا که رشته آغاز شده است). \code{main.salam} نامی است
که ویرایشگر برای پرونده‌ی برنامه‌ی شما به کار می‌برد. ویرایشگر زیر پیام، همان
خط از برنامه را هم نشان می‌دهد و با علامت \code{\textasciicircum} به جای
اشتباه اشاره می‌کند. عادت کنید پیام‌های خطا را با دقت بخوانید؛ آن‌ها دشمن شما نیستند،
راهنمای شما هستند. در پیوست \ref{app:errors} فهرستی از خطاهای رایج و معنی
آن‌ها را آورده‌ایم.

چند اشتباه بسیار رایج در همین نخستین قدم‌ها:
\begin{itemize}
  \item فراموش کردن \fcode{پایان} در پایان روال ریشه.
  \item فراموش کردن دونقطه‌ی بعد از \fcode{روال ریشه}.
  \item فراموش کردن ویرگول میان چیزهایی که با یک \fcode{سرچاپ} چاپ می‌شوند.
  \item نوشتن نام دستورها با املای اشتباه، مثلاً \fcode{سرچاب} به جای
        \fcode{سرچاپ}.
\end{itemize}

\begin{deep}
چگونه رایانه‌ای که فقط صفر و یک می‌فهمد، متن فارسی ما را اجرا می‌کند؟
پاسخ برنامه‌ای است به نام \textbf{مترجم} (یا کامپایلر). مترجم سلام کد شما
را می‌خواند، اول بررسی می‌کند که قاعده‌های زبان را رعایت کرده باشد (همان‌جا
که خطاها را می‌بینید) و سپس آن را به دستورهایی تبدیل می‌کند که پردازنده‌ی
رایانه مستقیم می‌فهمد. در ویرایشگر برخط، همین مترجم درون مرورگر شما کد را
بررسی می‌کند و سپس برنامه‌ای به نام \textbf{مفسر} دستورها را یکی‌یکی اجرا
می‌کند؛ نتیجه همان است، فقط راه رسیدن به آن فرق دارد. جالب است بدانید که مترجم سلام، خودش با زبان سلام نوشته شده است.
\end{deep}

\section{یک برنامه‌ی کامل‌تر}

بیایید هر چه در این فصل آموختیم را در یک برنامه کنار هم بگذاریم: کارت
شناسایی کوچکی که خودتان را معرفی می‌کند.

\program{کارت معرفی}
\begin{salamfa}
// کارت معرفی
روال ریشه:
    سرچاپ "================="
    سرچاپ "نام: سارا احمدی"
    سرچاپ "شهر: کاشان"
    سرچاپ "سن:", ۱۴۰۵ - ۱۳۸۸, "سال"
    سرچاپ "علاقه‌مندی: برنامه‌نویسی"
    سرچاپ "================="
پایان
\end{salamfa}

\begin{outputfa}
=================
نام: سارا احمدی
شهر: کاشان
سن: 17 سال
علاقه‌مندی: برنامه‌نویسی
=================
\end{outputfa}

خط «سن» را ببینید: به جای نوشتن عدد ۱۷، تفریق دو سال را نوشتیم و سلام خودش
حساب کرد. این یک عادت خوب است: هر جا عددی از محاسبه به دست می‌آید، بگذارید
رایانه حساب کند، نه شما.

\begin{summary}
\begin{itemize}
  \item برنامه فهرستی از دستورهاست که رایانه به ترتیب، از بالا به پایین،
        اجرا می‌کند.
  \item هر برنامه‌ی سلام با \scode{روال ریشه:} شروع و با \fcode{پایان}
        پایان می‌یابد. دستورها میان این دو خط و با تورفتگی نوشته می‌شوند.
  \item \fcode{سرچاپ} چیزی را می‌نویسد و به خط بعد می‌رود؛ \fcode{چاپ}
        می‌نویسد اما در همان خط می‌ماند.
  \item متن را میان \code{"} می‌نویسیم و عدد را بدون آن. چند چیز را با
        \code{,} از هم جدا می‌کنیم.
  \item هر چه بعد از \code{//} بیاید توضیح است و اجرا نمی‌شود.
  \item پیام خطا می‌گوید چه اشتباهی در کدام خط رخ داده است. آن را با دقت
        بخوانید.
\end{itemize}
\end{summary}

\section*{تمرین‌ها}

\exercise برنامه‌ای بنویسید که نام، نام خانوادگی و شهر شما را هر یک در یک
خط جداگانه چاپ کند.

\exercise برنامه‌ای بنویسید که با سه دستور \fcode{چاپ} و یک دستور
\fcode{سرچاپ}، جمله‌ی «من برنامه‌نویس هستم» را در یک خط چاپ کند.

\exercise برنامه‌ای بنویسید که این محاسبه‌ها را انجام دهد و نتیجه‌ی هر یک
را در یک خط، همراه با توضیحی کوتاه، چاپ کند: حاصل جمع ۱۲۵ و ۳۷۵؛
حاصل ضرب ۲۴ در ۷؛ تعداد ساعت‌های یک هفته.

\exercise برنامه‌ی زیر سه اشتباه دارد. بدون اجرا کردن، آن‌ها را پیدا کنید و
سپس با اجرا کردن پاسخ خود را بررسی کنید:
\begin{salamfa}
روال ریشه
    print "سلام"
    چاپ "خداحافظ
پایان
\end{salamfa}

\exercise برنامه‌ای بنویسید که یک شعر یا ضرب‌المثل چهارخطی دلخواه را چاپ
کند و بالای برنامه، در یک توضیح چندخطی، نام شاعر یا منبع آن را بنویسید.

\chapter{متغیر و حافظه}\label{ch:variables}

\begin{goals}
\begin{itemize}
  \item حافظه‌ی رایانه چگونه چیزها را نگه می‌دارد
  \item تعریف متغیر با \code{:=} و تغییر مقدار آن با \code{=}
  \item تفاوت مقدار ثابت و مقدار تغییرپذیر (\fcode{ناپایا})
  \item قاعده‌های نام‌گذاری و انتخاب نام خوب
  \item چند الگوی پرکاربرد: شمارنده، جمع‌کردن، جابه‌جایی
\end{itemize}
\end{goals}

\section{جعبه‌هایی با برچسب}

در فصل پیش هر عددی را همان‌جا که لازم بود نوشتیم: \scode{سرچاپ ۱۴۰۵ - ۱۳۸۸}.
اما برنامه‌های واقعی باید چیزهایی را \emph{به خاطر بسپارند}: نام کاربر، مبلغ
خرید، امتیاز بازی، تعداد دفعاتی که کاری انجام شده. جایی که برنامه این
چیزها را نگه می‌دارد، \textbf{حافظه}‌ی رایانه است.

حافظه را مثل قفسه‌ای بزرگ از جعبه‌های کوچک تصور کنید. در هر جعبه می‌توان
یک چیز گذاشت: یک عدد، یک متن، و یا چیزهایی که بعدها می‌بینیم. برای اینکه
بعداً بتوانیم جعبه را پیدا کنیم، روی آن برچسبی با یک نام می‌چسبانیم. به
چنین جعبه‌ای \textbf{متغیر} می‌گویند: جایی نام‌دار در حافظه که یک مقدار در
آن نگه داشته می‌شود.

\section{تعریف متغیر}

در سلام، برای ساختن یک جعبه‌ی تازه و گذاشتن چیزی در آن، از علامت \code{:=}
استفاده می‌کنیم. آن را «دونقطه مساوی» بخوانید، یا بهتر از آن: «تعریف کن
برابر با»:

\program{نخستین متغیرها}
\begin{salamfa}
تابع ریشه:
    نام := "سارا"
    سن := ۱۷
    سرچاپ "نام:", نام
    سرچاپ "سن:", سن
پایان
\end{salamfa}

\begin{outputfa}
نام: سارا
سن: 17
\end{outputfa}

خط \scode{نام := "سارا"} سه کار می‌کند: جعبه‌ای می‌سازد، برچسب
\fcode{نام} را روی آن می‌چسباند و متن \fcode{"سارا"} را در آن می‌گذارد. از
این پس، هر جا در برنامه \fcode{نام} بنویسیم، سلام به آن جعبه نگاه می‌کند و
مقدارش را برمی‌دارد. به همین دلیل \scode{سرچاپ "نام:", نام} عبارت «نام:
سارا» را چاپ کرد، نه «نام: نام».

\begin{note}
به تفاوت \fcode{"نام"} و \fcode{نام} دقت کنید. اولی، با نقل قول، یک متن
است که همان‌طور که هست چاپ می‌شود. دومی، بدون نقل قول، نام یک متغیر است و
سلام به جای آن مقدار درون جعبه را می‌گذارد.
\end{note}

متغیرها را می‌توان در محاسبه به کار برد، درست مثل عددها:

\program{محاسبه با متغیرها}
\begin{salamfa}
تابع ریشه:
    قیمت کتاب := ۱۲۰۰۰۰
    تعداد := ۳
    سرچاپ "مبلغ کل:", قیمت کتاب * تعداد, "تومان"
پایان
\end{salamfa}

\begin{outputfa}
مبلغ کل: 360000 تومان
\end{outputfa}

مزیت این کار روشن است: اگر فردا قیمت کتاب عوض شد، فقط یک خط را عوض می‌کنیم
و همه‌ی محاسبه‌هایی که از \fcode{قیمت کتاب} استفاده می‌کنند خودبه‌خود درست
می‌شوند.

\section{نام‌گذاری متغیرها}

در برنامه‌ی بالا نام یکی از متغیرها \fcode{قیمت کتاب} بود؛ دو واژه با یک
فاصله میانشان. این یکی از ویژگی‌های خوشایند سلام است: نام متغیر می‌تواند
از چند واژه ساخته شود، درست مثل زبان فارسی. قاعده‌های نام‌گذاری این‌هاست:
\begin{itemize}
  \item نام می‌تواند از حرف‌های فارسی، حرف‌های انگلیسی، رقم‌ها و فاصله
        ساخته شود.
  \item نام نمی‌تواند با رقم شروع شود: \fcode{۲نفر} نادرست است، اما
        \fcode{نفر۲} درست است.
  \item نام نمی‌تواند یکی از واژه‌های خود زبان باشد: نمی‌توانید متغیری به
        نام \fcode{سرچاپ} یا \fcode{پایان} یا \fcode{اگر} بسازید. این قاعده
        برای هر یک از واژه‌های یک نام چندواژه‌ای هم برقرار است؛ چون
        \fcode{هر}، \fcode{در}، \fcode{تا} و \fcode{و} واژه‌های
        زبان‌اند، نام‌هایی مثل \fcode{قیمت هر واحد} یا \fcode{بطری در جعبه}
        پذیرفته نمی‌شوند. فهرست این واژه‌ها را در پیوست~\ref{app:keywords}
        می‌بینید.
  \item حرف بزرگ و کوچک در حرف‌های انگلیسی فرق دارد؛ \fcode{x} و
        \fcode{X} دو متغیر جداگانه‌اند.
\end{itemize}

مهم‌تر از قاعده‌ها، \emph{عادت‌های خوب} است. نام متغیر باید بگوید درونش چه
چیزی است. مقایسه کنید:

\begin{salamfa}
تابع ریشه:
    الف := ۱۲۰۰۰۰
    ب := ۳
    سرچاپ الف * ب
پایان
\end{salamfa}

با همان برنامه‌ی پیشین. هر دو همان عدد را حساب می‌کنند، اما برنامه‌ی پیشین را یک هفته بعد هم
می‌توان فهمید و این یکی را نه. برای نام‌های موقت و کم‌اهمیت (مثلاً در یک
محاسبه‌ی دو خطی) نام‌های کوتاه اشکالی ندارد، اما برای هر چیزی که در برنامه
می‌ماند، نام کامل و گویا انتخاب کنید.

\begin{tip}
در سلام می‌توانید نام متغیرها را انگلیسی هم بنویسید؛ مثلاً در برخی از
نمونه‌های ویرایشگر برخط، شمارنده‌ی حلقه \fcode{i} نام دارد. اما در این
کتاب همه‌ی نام‌ها را فارسی می‌نویسیم تا برنامه یک‌دست و روان خوانده شود.
\end{tip}

\section{تغییر دادن مقدار: مقدار ثابت و متغیر}

حالا نکته‌ای که سلام را از بسیاری زبان‌ها متمایز می‌کند. جعبه‌ای که با
\code{:=} می‌سازید، به طور پیش‌فرض \textbf{مهر و موم شده} است: مقداری که
اول در آن گذاشته‌اید، تا پایان همان می‌ماند. تلاش برای عوض کردن آن خطاست:

\begin{salamfa}
تابع ریشه:
    سن := ۱۷
    سن = ۱۸
    سرچاپ سن
پایان
\end{salamfa}

\begin{outputfa}
خطا[خطا۰۱۳]: نمی‌توان متغیر تغییرناپذیر 'سن' را دوباره مقداردهی کرد؛
برای این کار آن را 'ناپایا' تعریف کنید،
یا به‌جای آن به عناصر یا فیلدهایش مقدار دهید
 --> main.salam:3:5
\end{outputfa}

\fcode{خطا۰۱۳} درون کروشه، شماره‌ی این نوع خطاست؛ هر نوع خطا در سلام شماره‌ی
خودش را دارد و با آن می‌توانید درباره‌اش جست‌وجو کنید. پیام خطا خودش راه حل را می‌گوید. اگر می‌خواهید مقدار جعبه‌ای در طول
برنامه عوض شود، هنگام تعریف، پیش از نامش واژه‌ی \kwd{ناپایا} را بنویسید:

\program{متغیر تغییرپذیر}
\begin{salamfa}
تابع ریشه:
    ناپایا سن := ۱۷
    سرچاپ "امسال:", سن
    سن = ۱۸
    سرچاپ "سال بعد:", سن
پایان
\end{salamfa}

\begin{outputfa}
امسال: 17
سال بعد: 18
\end{outputfa}

به دو علامت متفاوت دقت کنید:
\begin{itemize}
  \item \code{:=} فقط \emph{یک بار}، هنگام ساختن جعبه، به کار می‌رود.
  \item \code{=} برای \emph{عوض کردن} مقدار جعبه‌ای است که از پیش وجود
        دارد و با \fcode{ناپایا} تعریف شده است.
\end{itemize}
اگر برای متغیری که پیش‌تر ساخته‌اید دوباره \code{:=} بنویسید، سلام خطا
می‌دهد که این نام از قبل تعریف شده است.

\begin{note}[چرا سلام این‌قدر سخت‌گیر است؟]
شاید بپرسید چرا باید زحمت نوشتن \fcode{ناپایا} را بکشیم. پاسخ این است که
بیشتر مقدارهای یک برنامه هرگز عوض نمی‌شوند: قیمت، نام، تعداد روزهای هفته.
وقتی این‌ها را بدون \fcode{ناپایا} تعریف می‌کنید، سلام تضمین می‌کند که هیچ
جای برنامه، حتی به اشتباه، آن‌ها را عوض نمی‌کند. و وقتی جایی \fcode{ناپایا}
می‌بینید، بلافاصله می‌فهمید «این یکی تغییر می‌کند، حواسم باشد». این عادت
کوچک، جلوی بسیاری از اشتباه‌های بزرگ را می‌گیرد.
\end{note}

\section{سلام هر متغیر را جدی می‌گیرد}

نکته‌ی دیگری که در همین ابتدا باید بدانید: اگر متغیری تعریف کنید و هیچ‌جا
از آن استفاده نکنید، سلام خطا می‌دهد:

\begin{salamfa}
تابع ریشه:
    نام := "سارا"
    سن := ۱۷
    سرچاپ "نام:", نام
پایان
\end{salamfa}

\begin{outputfa}
خطا[خطا۰۵۹]: متغیر استفاده‌نشده 'سن': پس از تعریف هیچ‌جا به کار نرفته است؛
تعریف آن را حذف کنید یا مقدارش را همان‌جا که لازم است به کار ببرید؛
آغاز کردن نام با '_' این خطا را پنهان می‌کند اما توصیه نمی‌شود،
چون کد مرده در برنامه می‌ماند
 --> main.salam:3:5
\end{outputfa}

دلیلش ساده است: متغیری که استفاده نمی‌شود یا اضافی است، یا نشانه‌ی اینکه
جایی چیزی را فراموش کرده‌اید. در هر دو حالت بهتر است بدانید. راه حل هم
ساده است: یا از آن استفاده کنید، یا خطش را پاک کنید.

\section{کپی شدن مقدارها}

وقتی مقدار یک متغیر را در متغیر دیگری می‌گذارید، \emph{یک نسخه} از آن مقدار
در جعبه‌ی دوم قرار می‌گیرد. دو جعبه به هم وصل نیستند:

\program{کپی شدن مقدار}
\begin{salamfa}
تابع ریشه:
    ناپایا الف := ۱۰
    ب := الف
    الف = ۲۰
    سرچاپ "الف:", الف
    سرچاپ "ب:", ب
پایان
\end{salamfa}

\begin{outputfa}
الف: 20
ب: 10
\end{outputfa}

با اینکه \fcode{ب} از روی \fcode{الف} ساخته شد، عوض شدن \fcode{الف} هیچ
تأثیری روی \fcode{ب} نگذاشت. در لحظه‌ی \scode{ب := الف}، عدد ۱۰ کپی
شد و رفت در جعبه‌ی \fcode{ب}؛ از آن به بعد هر جعبه راه خودش را می‌رود.

\section{الگوهای پرکاربرد}

چند کار هست که برنامه‌نویسان هر روز با متغیرها انجام می‌دهند. بیایید آن‌ها
را یکی‌یکی ببینیم.

\subsection{شمارنده}

شاید عجیب‌ترین خط برای تازه‌کارها این باشد:

\begin{salamfa}
شمارنده = شمارنده + ۱
\end{salamfa}

از نگاه ریاضی، این جمله بی‌معنی است؛ هیچ عددی با یک واحد بیشتر از خودش
برابر نیست. اما به یاد بیاورید که \code{=} در برنامه‌نویسی «برابر است» نیست،
بلکه «بگذار در» است. این خط می‌گوید: مقدار کنونی \fcode{شمارنده} را بردار،
یکی به آن اضافه کن و نتیجه را دوباره در همان جعبه بگذار. اول سمت راست
حساب می‌شود، بعد نتیجه در سمت چپ گذاشته می‌شود.

\program{شمارنده}
\begin{salamfa}
تابع ریشه:
    ناپایا شمارنده := ۰
    سرچاپ شمارنده
    شمارنده = شمارنده + ۱
    سرچاپ شمارنده
    شمارنده = شمارنده + ۱
    سرچاپ شمارنده
پایان
\end{salamfa}

\begin{outputfa}
0
1
2
\end{outputfa}

این کار آن‌قدر رایج است که سلام برایش نوشتار کوتاه‌تری دارد.
\scode{شمارنده += ۱} دقیقاً همان \scode{شمارنده = شمارنده + ۱}
است، و \scode{شمارنده++} همان \scode{شمارنده += ۱}. هر سه شکل
درست‌اند و در این کتاب هر سه را خواهید دید:

\begin{salamfa}
تابع ریشه:
    ناپایا شمارنده := ۰
    شمارنده = شمارنده + ۱
    شمارنده += ۱
    شمارنده++
    سرچاپ شمارنده
پایان
\end{salamfa}

\begin{outputfa}
3
\end{outputfa}

\subsection{جمع کردن}

همین الگو برای جمع زدن هم به کار می‌رود. متغیری با مقدار صفر می‌سازیم و هر
چیزی را که باید جمع شود به آن اضافه می‌کنیم:

\program{جمع خریدها}
\begin{salamfa}
تابع ریشه:
    ناپایا مجموع := ۰
    مجموع += ۴۵۰۰۰    // نان
    مجموع += ۱۲۰۰۰۰   // پنیر
    مجموع += ۳۲۰۰۰    // شیر
    سرچاپ "جمع خرید:", مجموع, "تومان"
پایان
\end{salamfa}

\begin{outputfa}
جمع خرید: 197000 تومان
\end{outputfa}

\subsection{جابه‌جایی دو مقدار}

فرض کنید می‌خواهید محتوای دو جعبه را با هم عوض کنید. نمی‌شود اول \fcode{الف}
را در \fcode{ب} گذاشت و بعد \fcode{ب} را در \fcode{الف}، چون بعد از قدم
اول، مقدار قبلی \fcode{ب} از دست رفته است. راه حل، یک جعبه‌ی سوم است:

\program{جابه‌جایی}
\begin{salamfa}
تابع ریشه:
    ناپایا الف := ۵
    ناپایا ب := ۹
    سرچاپ "پیش از جابه‌جایی:", الف, ب
    موقت := الف
    الف = ب
    ب = موقت
    سرچاپ "پس از جابه‌جایی:", الف, ب
پایان
\end{salamfa}

\begin{outputfa}
پیش از جابه‌جایی: 5 9
پس از جابه‌جایی: 9 5
\end{outputfa}

قدم به قدم دنبال کنید: \fcode{موقت} یک نسخه از ۵ را نگه می‌دارد، \fcode{الف}
مقدار ۹ می‌گیرد، و در آخر \fcode{ب} مقدار ۵ را از \fcode{موقت} پس می‌گیرد.
دقت کنید که \fcode{موقت} را بدون \fcode{ناپایا} تعریف کردیم، چون فقط یک بار
مقدار می‌گیرد و هرگز عوض نمی‌شود.

\section{پایدارها}

بعضی مقدارها نه تنها در طول اجرای برنامه عوض نمی‌شوند، بلکه اصلاً ذاتشان
ثابت است: تعداد روزهای هفته، نرخ مالیات، عدد پی. برای این‌ها سلام واژه‌ی
\kwd{پایا} را دارد:

\program{پایا}
\begin{salamfa}
تابع ریشه:
    پایا نرخ := ۹
    قیمت := ۵۰۰۰۰۰
    مالیات := (قیمت * نرخ) / ۱۰۰
    سرچاپ "مالیات:", مالیات
    سرچاپ "قیمت با مالیات:", قیمت + مالیات
پایان
\end{salamfa}

\begin{outputfa}
مالیات: 45000
قیمت با مالیات: 545000
\end{outputfa}

از نگاه برنامه، \scode{پایا نرخ := ۹} با \scode{نرخ := ۹}
تفاوت چندانی ندارد؛ هر دو تغییرناپذیرند. تفاوت در پیام است: \fcode{پایا}
به خواننده می‌گوید «این یک عدد مهم است که با قصد این‌جا گذاشته‌ایم». دو
نکته‌ی کوچک: نام پایدار باید یک واژه باشد (برخلاف متغیرها که می‌توانند
چندواژه‌ای باشند)، و پایدارها را می‌توان بیرون از \fcode{تابع ریشه} و در
بالای برنامه هم تعریف کرد؛ آن وقت در همه جای برنامه در دسترس‌اند. در فصل~\ref{ch:functions} این کار را خواهیم دید.

\begin{deep}[متغیر واقعاً کجاست؟]
حافظه‌ی رایانه واقعاً مثل قفسه‌ای از خانه‌های شماره‌دار است؛ به هر خانه یک
\textbf{نشانی} می‌گویند و در هر خانه یک عدد کوچک جا می‌شود. وقتی می‌نویسید
\scode{سن := ۱۷}، مترجم سلام چند خانه‌ی خالی پیدا می‌کند، نشانی آن‌ها
را به نام \fcode{سن} گره می‌زند و عدد ۱۷ را در آن‌ها می‌نویسد. بعد از آن،
هر جا \fcode{سن} بنویسید، به آن نشانی ترجمه می‌شود. شما هرگز لازم نیست
نشانی‌ها را ببینید؛ نام‌ها برای همین اختراع شده‌اند. متن‌ها مثل
\fcode{"سارا"} کمی بیشتر جا می‌گیرند (هر حرف چند خانه)، اما اصل کار همان
است.
\end{deep}

\section{یک برنامه‌ی کامل‌تر}

صورت‌حساب یک خرید اینترنتی را با متغیرها بنویسیم:

\program{صورت‌حساب خرید}
\begin{salamfa}
تابع ریشه:
    قیمت واحد := ۴۵۰۰۰
    تعداد := ۴
    هزینه ارسال := ۳۰۰۰۰
    جمع کالا := قیمت واحد * تعداد
    سرچاپ "جمع کالاها:", جمع کالا
    سرچاپ "هزینه ارسال:", هزینه ارسال
    سرچاپ "مبلغ قابل پرداخت:", جمع کالا + هزینه ارسال
پایان
\end{salamfa}

\begin{outputfa}
جمع کالاها: 180000
هزینه ارسال: 30000
مبلغ قابل پرداخت: 210000
\end{outputfa}

هیچ‌کدام از متغیرها با \fcode{ناپایا} تعریف نشده‌اند، چون هیچ‌کدام در طول
برنامه عوض نمی‌شوند. این برنامه را با نوشتن مستقیم عددها هم می‌شد نوشت،
اما آن وقت اگر تعداد از ۴ به ۵ تغییر می‌کرد، باید چند جا را دست می‌زدیم و
احتمالاً یکی را جا می‌انداختیم.

\begin{summary}
\begin{itemize}
  \item متغیر جایی نام‌دار در حافظه است که یک مقدار را نگه می‌دارد.
  \item \scode{نام := مقدار} متغیر تازه‌ای می‌سازد که دیگر عوض نمی‌شود.
        \scode{ناپایا نام := مقدار} متغیری می‌سازد که می‌توان بعداً با
        \code{=} مقدارش را عوض کرد.
  \item نام‌ها می‌توانند چندواژه‌ای باشند، اما با رقم شروع نمی‌شوند و
        واژه‌های زبان نیستند. نام گویا انتخاب کنید.
  \item متغیری که تعریف شده اما استفاده نشده، خطاست.
  \item \scode{ب := الف} یک نسخه از مقدار می‌سازد؛ دو متغیر به هم وصل
        نیستند.
  \item \scode{شمارنده += ۱} و \scode{شمارنده++} شکل‌های کوتاه
        \scode{شمارنده = شمارنده + ۱} هستند.
  \item \fcode{پایا} برای مقدارهایی است که ذاتاً ثابت‌اند.
\end{itemize}
\end{summary}

\section*{تمرین‌ها}
\markright{تمرین‌ها}

\exercise سه متغیر برای طول، عرض و ارتفاع یک اتاق (به متر) تعریف کنید و
مساحت کف و حجم اتاق را حساب و چاپ کنید.

\exercise متغیری به نام \fcode{موجودی} با مقدار اولیه‌ی ۵۰۰٬۰۰۰ بسازید.
سپس با سه دستور جداگانه، ۱۲۰٬۰۰۰ تومان برداشت، ۲۰۰٬۰۰۰ تومان واریز و
۷۵٬۰۰۰ تومان برداشت کنید و بعد از هر عملیات موجودی را چاپ کنید. (در
برنامه، عددها را بدون جداکننده‌ی هزارگان بنویسید: \fcode{۵۰۰۰۰۰}.)

\Needspace{12\baselineskip}
\exercise برنامه‌ی زیر را بدون اجرا کردن دنبال کنید و بگویید چه چاپ
می‌کند. سپس با اجرا کردن پاسخ خود را بررسی کنید:
\begin{salamfa}
تابع ریشه:
    ناپایا الف := ۳
    ناپایا ب := الف * ۲
    الف = ب + ۱
    ب = الف - ب
    سرچاپ الف, ب
پایان
\end{salamfa}

\exercise برنامه‌ی زیر دو خطا دارد. آن‌ها را پیدا و برطرف کنید:
\begin{salamfa}
تابع ریشه:
    امتیاز := ۱۰
    امتیاز = امتیاز + ۵
    نام بازیکن := "رضا"
    سرچاپ "امتیاز:", امتیاز
پایان
\end{salamfa}

\exercise سه متغیر \fcode{اول}، \fcode{دوم} و \fcode{سوم} با مقدارهای ۱، ۲
و ۳ بسازید و بدون نوشتن مستقیم عددها، مقدارها را طوری بچرخانید که
\fcode{اول} برابر ۲، \fcode{دوم} برابر ۳ و \fcode{سوم} برابر ۱ شود.

\chapter{گونه‌های داده}
\label{ch:types}

\begin{goals}
\begin{itemize}
  \item چهار گونه‌ی داده‌ی پایه: عدد صحیح، عدد اعشاری، رشته و منطقی
  \item سلام چگونه گونه‌ی هر مقدار را تشخیص می‌دهد
  \item تقسیم عددهای صحیح و تفاوت آن با تقسیم اعشاری
  \item تبدیل گونه با \fcode{برگردان}
  \item چسباندن متن و عدد به هم
\end{itemize}
\end{goals}

\section{هر مقدار یک گونه دارد}

در فصل پیش دیدیم که در جعبه‌ی یک متغیر می‌توان عدد یا متن گذاشت. اما رایانه
عدد و متن را یکسان نگه نمی‌دارد و یکسان با آن‌ها رفتار نمی‌کند. برای عددها
می‌توان جمع و ضرب حساب کرد، برای متن نه؛ متن را می‌توان حرف‌به‌حرف خواند،
عدد را نه. به همین دلیل هر مقدار در سلام یک \textbf{گونه} دارد و سلام
همیشه می‌داند در هر جعبه چه گونه‌ای داده است.

در این کتاب با چهار گونه‌ی پایه سروکار داریم:

\begin{center}
\begin{tabular}{lll}
\toprule
گونه & نام در سلام & نمونه \\
\midrule
عدد صحیح & \fcode{صحیح} & \fcode{۱۷}، \scode{-۴}، \fcode{۰} \\
عدد اعشاری & \fcode{اعشار۶۴} & \scode{۲.۵}، \scode{-۰.۷۵}، \scode{۳.۰} \\
رشته (متن) & \fcode{رشته} & \fcode{"سلام"}، \fcode{"17"}، \fcode{""} \\
منطقی & \fcode{منطقی} & \fcode{درست}، \fcode{نادرست} \\
\bottomrule
\end{tabular}
\end{center}

سلام از روی شکل مقداری که می‌نویسید، گونه‌اش را خودش تشخیص می‌دهد. لازم
نیست گونه را بنویسید:

\program{چهار گونه‌ی داده}
\begin{salamfa}
روال ریشه:
    تعداد := ۱۲          // عدد صحیح
    قیمت := ۲.۵          // عدد اعشاری
    نام := "کتاب"        // رشته
    موجود := درست        // منطقی
    سرچاپ تعداد, قیمت, نام, موجود
پایان
\end{salamfa}

\begin{outputfa}
12 2.5 کتاب true
\end{outputfa}

قاعده‌ی تشخیص ساده است: عددی بدون نقطه‌ی اعشار، صحیح است؛ عددی با نقطه،
اعشاری است؛ هر چه میان نقل قول باشد، رشته است؛ و \fcode{درست} و
\fcode{نادرست} منطقی‌اند. به آخرین مقدار خروجی دقت کنید: سلام مقدارهای
منطقی را به شکل \code{true} و \code{false} چاپ می‌کند.

\section{عدد صحیح}

عدد صحیح همان عددهای بی‌اعشار است: تعداد دانش‌آموزان، شماره‌ی صفحه، سال.
می‌تواند منفی هم باشد. جمع، تفریق و ضرب دو عدد صحیح، همیشه یک عدد صحیح
می‌دهد. اما تقسیم داستان دیگری دارد:

\program{تقسیم عددهای صحیح}
\begin{salamfa}
روال ریشه:
    سرچاپ ۷ / ۲
    سرچاپ ۹ / ۳
    سرچاپ ۱ / ۴
پایان
\end{salamfa}

\begin{outputfa}
3
3
0
\end{outputfa}

هفت تقسیم بر دو، سه و نیم است، اما سلام \code{3} چاپ کرد. چرا؟ چون هر دو
عدد صحیح بودند و \textbf{تقسیم دو عدد صحیح در سلام همیشه یک عدد صحیح
می‌دهد}: بخش اعشاری دور ریخته می‌شود. به همین دلیل یک تقسیم بر چهار، صفر
شد. این رفتار در آغاز عجیب به نظر می‌رسد، اما بسیار پرکاربرد است. اگر ۷
سیب را میان ۲ نفر تقسیم کنید، هر کس ۳ سیب کامل می‌گیرد و ۱ سیب باقی
می‌ماند. سلام برای همین «باقی‌مانده» هم علامتی دارد: \code{\%}.

\program{خارج قسمت و باقی‌مانده}
\begin{salamfa}
روال ریشه:
    سیب := ۷
    نفر := ۲
    سرچاپ "سهم هر نفر:", سیب / نفر
    سرچاپ "باقی‌مانده:", سیب % نفر
پایان
\end{salamfa}

\begin{outputfa}
سهم هر نفر: 3
باقی‌مانده: 1
\end{outputfa}

کارور باقی‌مانده در فصل‌های بعد بارها به کارمان می‌آید؛ مثلاً برای فهمیدن
اینکه عددی زوج است یا فرد.

\begin{deep}[عدد صحیح تا کجا بزرگ می‌شود؟]
هر عدد صحیح در سلام در جعبه‌ای با اندازه‌ی مشخص نگه داشته می‌شود و به همین
دلیل حدی دارد: از حدود منفی دو میلیارد تا مثبت دو میلیارد. برای همه‌ی
برنامه‌های این کتاب این حد بیش از اندازه کافی است، اما اگر روزی خواستید
جمعیت جهان یا بودجه‌ی کشور را به ریال حساب کنید، به آن برمی‌خورید. سلام
برای این کارها گونه‌های بزرگ‌تر هم دارد، مثل \fcode{صحیح۶۴}، اما در این کتاب
به آن‌ها نیازی نخواهیم داشت.
\end{deep}

\section{عدد اعشاری}

وقتی به عددهایی با بخش اعشاری نیاز داریم (قیمت به دلار، وزن، معدل، دما)
از عدد اعشاری استفاده می‌کنیم. کافی است در نوشتن عدد نقطه بگذاریم. حتی
\scode{۳.۰} هم یک عدد اعشاری است، هرچند مقدارش سه باشد:

\program{محاسبه با عددهای اعشاری}
\begin{salamfa}
روال ریشه:
    سرچاپ ۷.۰ / ۲.۰
    سرچاپ ۲.۵ * ۴.۰
    سرچاپ ۱۰.۰ / ۴.۰
    سرچاپ ۱.۰ / ۳.۰
پایان
\end{salamfa}

\begin{outputfa}
3.5
10
2.5
0.3333333333333333
\end{outputfa}

دو نکته در این خروجی هست. نخست، \scode{۲.۵ * ۴.۰} برابر \scode{۱۰.۰} است، اما
سلام آن را \code{10} چاپ کرد: سلام هنگام چاپ عددهای اعشاری، صفرهای بی‌فایده‌ی
پایان عدد را نمی‌نویسد. دوم، یک سوم را با شانزده رقم اعشار چاپ کرد، چون
یک سوم اعشار بی‌پایان دارد و رایانه فقط تا حدی از آن را نگه می‌دارد.

\begin{warn}
عددهای اعشاری در رایانه \emph{تقریبی} نگه داشته می‌شوند. برنامه‌ی زیر را
اجرا کنید:
\begin{salamfa}
روال ریشه:
    سرچاپ ۰.۱ + ۰.۲
پایان
\end{salamfa}
\begin{outputfa}
0.30000000000000004
\end{outputfa}
این اشکال سلام نیست؛ همه‌ی رایانه‌های جهان همین را می‌گویند، چون عدد ۰٫۱
در مبنای دو (که رایانه با آن کار می‌کند) اعشار بی‌پایان دارد، درست مثل یک
سوم در مبنای ده. برای کارهای روزمره این خطای بسیار کوچک هیچ اهمیتی
ندارد، اما به خاطر داشته باشید که هرگز نباید دو عدد اعشاری را برای
«دقیقاً برابر بودن» مقایسه کنید.
\end{warn}

\section{وقتی صحیح و اعشاری قاطی می‌شوند}

اگر یک عدد صحیح و یک عدد اعشاری را با هم جمع یا تقسیم کنید، سلام عدد صحیح
را خودش به اعشاری تبدیل می‌کند و نتیجه اعشاری است: \scode{۷ / ۲.۰} برابر
\code{3.5} است. اما وقتی \emph{هر دو} طرف صحیح‌اند، تقسیم صحیح انجام می‌شود.
اگر در این حالت نتیجه‌ی اعشاری بخواهید، باید دست‌کم یکی از دو طرف را به
اعشاری تبدیل کنید. برای این کار واژه‌ی \kwd{برگردان} هست که یک مقدار را به
گونه‌ی دیگری تبدیل می‌کند:

\program{تبدیل صحیح به اعشاری}
\begin{salamfa}
روال ریشه:
    مجموع := ۱۷
    تعداد := ۴
    میانگین := (مجموع برگردان اعشار۶۴) / (تعداد برگردان اعشار۶۴)
    سرچاپ "میانگین:", میانگین
پایان
\end{salamfa}

\begin{outputfa}
میانگین: 4.25
\end{outputfa}

اگر تبدیل را انجام نمی‌دادیم، \scode{مجموع / تعداد} تقسیم دو عدد
صحیح بود و \code{4} می‌داد. پرانتزها الزامی نیستند، اما روشن می‌کنند که
تبدیل روی کدام مقدار انجام می‌شود و خواندن عبارت را آسان‌تر می‌کنند.

تبدیل در جهت عکس هم ممکن است. وقتی عددی اعشاری را با \fcode{برگردان صحیح} به صحیح
تبدیل می‌کنید، بخش اعشاری آن دور ریخته می‌شود (گرد نمی‌شود، بریده می‌شود):

\program{تبدیل اعشاری به صحیح}
\begin{salamfa}
روال ریشه:
    قیمت دقیق := ۳.۷۵
    قیمت رند := قیمت دقیق برگردان صحیح
    سرچاپ قیمت رند
    منفی := -۳.۷۵
    سرچاپ منفی برگردان صحیح
پایان
\end{salamfa}

\begin{outputfa}
3
-3
\end{outputfa}

\begin{note}
\fcode{برگردان} روی هر مقداری کار می‌کند، چه متغیر باشد و چه عددی که همان‌جا
نوشته‌اید: \scode{۳.۷۵ برگردان صحیح} هم \code{3} می‌دهد.
\end{note}

\section{رشته}

رشته یعنی متن: دنباله‌ای از حرف‌ها، رقم‌ها و علامت‌ها که میان نقل قول
نوشته می‌شود. رشته می‌تواند خالی باشد (\code{""}) و می‌تواند رقم داشته باشد،
اما رقمِ درون رشته عدد نیست:

\program{رشته یا عدد؟}
\begin{salamfa}
روال ریشه:
    سرچاپ ۱۲ + ۳
    سرچاپ "12" + ۳
    سرچاپ "12" + "3"
پایان
\end{salamfa}

\begin{outputfa}
15
123
123
\end{outputfa}

خط‌های دوم و سوم خروجی را ببینید. علامت \code{+} برای رشته‌ها معنی «چسباندن» دارد،
نه جمع کردن. \scode{"12" + ۳} یعنی رشته‌ی «۱۲» را به عدد ۳ بچسبان، که
می‌شود رشته‌ی «۱۲۳». اگر یکی از دو طرف \code{+} رشته باشد، سلام طرف دیگر را
هم به رشته تبدیل می‌کند و دو رشته را به هم می‌چسباند. این ویژگی برای ساختن
پیام‌ها بسیار مفید است:

\program{ساختن پیام}
\begin{salamfa}
روال ریشه:
    نام := "سارا"
    سن := ۱۷
    پیام := نام + " " + سن + " سال دارد."
    سرچاپ پیام
پایان
\end{salamfa}

\begin{outputfa}
سارا 17 سال دارد.
\end{outputfa}

به فاصله‌ها دقت کنید: چسباندن، خودش فاصله‌ای نمی‌گذارد، پس هر جا فاصله
می‌خواستیم، رشته‌ای با یک فاصله (\code{" "}) اضافه کردیم. مقایسه کنید با
\fcode{سرچاپ نام, سن} که میان چیزهایش خودش فاصله می‌گذارد. هر دو راه درست‌اند؛
با \code{+} کنترل بیشتری بر شکل پیام دارید و با ویرگول راحت‌ترید.

اگر رشته‌ای رقم دارد و می‌خواهید با آن حساب کنید، باید آن را به عدد تبدیل
کنید. برای این کار رشته‌ها یک قابلیت به نام \fcode{به صحیح} دارند که با
نقطه به آن‌ها می‌چسبد:

\program{تبدیل رشته به عدد}
\begin{salamfa}
روال ریشه:
    متن := "12"
    عدد := متن.به صحیح()
    سرچاپ عدد + ۳
پایان
\end{salamfa}

\begin{outputfa}
15
\end{outputfa}

رقم‌های درون رشته می‌توانند فارسی باشند یا انگلیسی: \scode{"۱۲".به صحیح()} هم
همان ۱۲ را می‌دهد. به همین شکل \fcode{به اعشار} یک رشته را به عدد اعشاری تبدیل می‌کند. در فصل
\ref{ch:strings} با قابلیت‌های دیگر رشته‌ها آشنا می‌شویم.

\section{منطقی}

گونه‌ی منطقی فقط دو مقدار دارد: \fcode{درست} و \fcode{نادرست}. شاید کوچک به
نظر برسد، اما همه‌ی تصمیم‌گیری‌های برنامه (فصل \ref{ch:conditions}) بر
پایه‌ی همین دو مقدار است. بیشتر وقت‌ها مقدار منطقی را مستقیم نمی‌نویسیم،
بلکه از یک مقایسه به دست می‌آید:

\program{مقدارهای منطقی}
\begin{salamfa}
روال ریشه:
    سن := ۲۰
    بزرگسال است := سن >= ۱۸
    سرچاپ بزرگسال است
    سرچاپ سن == ۱۸
    سرچاپ "سلام" == "سلام"
پایان
\end{salamfa}

\begin{outputfa}
true
false
true
\end{outputfa}

علامت \code{>=} یعنی «بزرگ‌تر یا مساوی» و علامت \code{==} (دو مساوی پشت
هم) یعنی «آیا برابر است؟». دقت کنید که \code{==} با \code{=} فرق دارد: اولی
می‌پرسد، دومی مقدار می‌دهد. در فصل بعد همه‌ی علامت‌های مقایسه را می‌بینیم.

\section{گونه‌ی یک متغیر عوض نمی‌شود}

وقتی متغیری با یک گونه ساخته می‌شود، تا آخر همان گونه را دارد. نمی‌توان در
جعبه‌ای که برای عدد ساخته شده، متن گذاشت:

\begin{salamfa}
روال ریشه:
    ناپایا مقدار := ۱۰
    مقدار = "ده"
    سرچاپ مقدار
پایان
\end{salamfa}

\begin{outputfa}
خطا[خطا۰۰۲]: نمی‌توان 'رشته' را به 'صحیح' نسبت داد
 --> main.salam:3:5
\end{outputfa}

پیام می‌گوید مقداری از گونه‌ی \fcode{رشته} را نمی‌توان در متغیری از گونه‌ی
\fcode{صحیح} گذاشت. این سخت‌گیری هم مثل
بقیه‌ی سخت‌گیری‌های سلام به سود شماست: وقتی متغیری به نام \fcode{تعداد}
می‌بینید، مطمئن هستید که همیشه یک عدد در آن است.

\section{یک برنامه‌ی کامل‌تر}

محاسبه‌ی معدل سه درس با واحدهای متفاوت، که همه‌ی گونه‌های این فصل را کنار
هم می‌آورد:

\program{معدل وزنی}
\begin{salamfa}
روال ریشه:
    // نمره‌ها اعشاری‌اند، واحدها صحیح
    نمره ریاضی := ۱۷.۵
    واحد ریاضی := ۳
    نمره فیزیک := ۱۵.۰
    واحد فیزیک := ۲
    نمره ادبیات := ۱۸.۲۵
    واحد ادبیات := ۲

    مجموع واحد := واحد ریاضی + واحد فیزیک + واحد ادبیات
    مجموع وزنی := (نمره ریاضی * (واحد ریاضی برگردان اعشار۶۴)) +
                  (نمره فیزیک * (واحد فیزیک برگردان اعشار۶۴)) +
                  (نمره ادبیات * (واحد ادبیات برگردان اعشار۶۴))
    معدل := مجموع وزنی / (مجموع واحد برگردان اعشار۶۴)

    سرچاپ "مجموع واحدها:", مجموع واحد
    سرچاپ "معدل:", معدل
    سرچاپ "قبول؟", معدل >= ۱۲.۰
پایان
\end{salamfa}

\begin{outputfa}
مجموع واحدها: 7
معدل: 17
قبول؟ true
\end{outputfa}

به دو چیز دقت کنید. نخست اینکه یک محاسبه‌ی طولانی را در چند خط شکسته‌ایم؛
وقتی خطی به کاروری مثل \code{+} ختم می‌شود، سلام می‌فهمد که عبارت هنوز
تمام نشده و آن را در خط بعد دنبال می‌کند. دوم اینکه معدل واقعاً ۱۷٫۰ است (مجموع وزنی ۱۱۹ تقسیم بر ۷) و
سلام مثل همیشه آن را \code{17} چاپ کرده است.

\begin{summary}
\begin{itemize}
  \item هر مقدار یک گونه دارد: \fcode{صحیح} (عدد بی‌اعشار)، \fcode{اعشار۶۴}
        (عدد با اعشار)، \fcode{رشته} (متن) و \fcode{منطقی} (درست یا
        نادرست). سلام گونه را از شکل مقدار تشخیص می‌دهد.
  \item تقسیم دو عدد صحیح، بخش اعشاری را دور می‌ریزد؛ اگر یک طرف اعشاری
        باشد، نتیجه اعشاری است. \code{\%}
        باقی‌مانده‌ی تقسیم را می‌دهد.
  \item عددهای اعشاری تقریبی‌اند؛ آن‌ها را برای برابری دقیق مقایسه نکنید.
  \item \fcode{مقدار برگردان گونه} یک مقدار را به گونه‌ی دیگری تبدیل می‌کند.
        تبدیل اعشاری به صحیح، اعشار را می‌برد.
  \item \code{+} میان رشته و هر چیز دیگر، به معنی چسباندن است.
        \scode{متن.به صحیح()} رشته را به عدد تبدیل می‌کند.
  \item گونه‌ی یک متغیر بعد از ساخته شدن عوض نمی‌شود.
\end{itemize}
\end{summary}

\section*{تمرین‌ها}
\markright{تمرین‌ها}

\exercise بدون اجرا کردن بگویید هر یک از این دستورها چه چاپ می‌کند و سپس
بررسی کنید: \scode{سرچاپ ۱۰ / ۳}، \scode{سرچاپ ۱۰.۰ / ۳.۰}، \scode{سرچاپ ۱۰ \% ۳}،
\scode{سرچاپ "10" + ۳}، \scode{سرچاپ ۱۰ + ۳}.

\exercise دمای هوا به سلسیوس را در متغیری بگذارید و آن را به فارنهایت تبدیل
و چاپ کنید. فرمول: فارنهایت برابر است با سلسیوس ضرب در ۹ تقسیم بر ۵،
به علاوه‌ی ۳۲. برنامه را یک بار با عدد اعشاری و یک بار با عدد صحیح
بنویسید و تفاوت خروجی‌ها را توضیح دهید.

\exercise تعداد کل دقیقه‌های یک فیلم (مثلاً ۱۳۵) را در متغیری بگذارید و
با تقسیم و باقی‌مانده، آن را به شکل «۲ ساعت و ۱۵ دقیقه» چاپ کنید.

\exercise سه متغیر با مقدارهای \code{"3"}، \code{۳} و \scode{۳.۰} بسازید.
آن‌ها را چاپ کنید، سپس هر یک را با عدد ۲ جمع کنید و نتیجه را چاپ کنید.
پیش از اجرا، خروجی را حدس بزنید.

\exercise قیمت یک کالا به دلار (اعشاری) و نرخ دلار به تومان (صحیح) را در
دو متغیر بگذارید و قیمت کالا به تومان را به شکل یک عدد صحیح چاپ کنید.

\chapter{کارور‌ها و محاسبه}
\label{ch:operators}

\begin{goals}
\begin{itemize}
  \item کارور‌های حسابی و ترتیب انجام آن‌ها
  \item کارور‌های مقایسه و نتیجه‌ی منطقی آن‌ها
  \item ترکیب شرط‌ها با \fcode{و}، \fcode{یا} و \fcode{وارونه}
  \item شکل‌های کوتاه مقداردهی
  \item چند محاسبه‌ی روزمره با همین ابزارها
\end{itemize}
\end{goals}

\section{کارور چیست؟}

علامت‌هایی مثل \code{+} و \code{*} که با آن‌ها محاسبه می‌کنیم،
\textbf{کارور} نام دارند و چیزهایی که روی آن‌ها کار می‌کنند
\textbf{عملوند}. در \scode{۲ + ۳}، علامت \code{+} کارور است و \code{۲} و
\code{۳} عملوندها. تا این‌جا با چند کارور آشنا شده‌اید؛ در این فصل همه‌ی
کارور‌های پایه‌ی سلام را یک‌جا می‌بینیم.

\section{کارور‌های حسابی}

\begin{center}
\begin{tabular}{clc}
\toprule
کارور & معنی & نمونه \\
\midrule
\code{+} & جمع & \scode{۷ + ۲} برابر \code{۹} \\
\code{-} & تفریق & \scode{۷ - ۲} برابر \code{۵} \\
\code{*} & ضرب & \scode{۷ * ۲} برابر \code{۱۴} \\
\code{/} & تقسیم & \scode{۷ / ۲} برابر \code{۳} \\
\code{\%} & باقی‌مانده & \scode{۷ \% ۲} برابر \code{۱} \\
\code{**} & توان & \scode{۲ ** ۳} برابر \code{۸} \\
\bottomrule
\end{tabular}
\end{center}

\program{کارور‌های حسابی}
\begin{salamfa}
روال ریشه:
    سرچاپ ۷ + ۲, ۷ - ۲, ۷ * ۲
    سرچاپ ۷ / ۲, ۷ % ۲
    سرچاپ ۲ ** ۳, ۲ ** ۱۰
    سرچاپ ۱۰ - ۱۵
پایان
\end{salamfa}

\begin{outputfa}
9 5 14
3 1
8 1024
-5
\end{outputfa}

کارور توان، \code{**}، همیشه نتیجه‌ی
اعشاری می‌دهد، حتی اگر هر دو عملوند صحیح باشند. علامت منفی را می‌توان
پیش از عدد یا متغیر گذاشت تا قرینه‌اش را بدهد: \scode{-سن}.

\section{ترتیب انجام کارور‌ها}

وقتی چند کارور در یک عبارت باشند، سلام همان قاعده‌ای را رعایت می‌کند که در
ریاضی مدرسه یاد گرفته‌اید: اول توان، بعد ضرب و تقسیم و باقی‌مانده، بعد جمع
و تفریق. کارور‌های هم‌رتبه به همان ترتیبی که در عبارت آمده‌اند، یکی پس از
دیگری، انجام می‌شوند. تنها استثنا توان است که از راست به چپ حساب می‌شود:
\scode{۲ ** ۲ ** ۳} یعنی \scode{۲ ** (۲ ** ۳)} و حاصلش ۲۵۶ است، نه ۶۴.
نمونه‌ی زیر ترتیب بقیه‌ی کارور‌ها را نشان می‌دهد:

\program{ترتیب کارور‌ها}
\begin{salamfa}
روال ریشه:
    سرچاپ ۲ + ۳ * ۴
    سرچاپ (۲ + ۳) * ۴
    سرچاپ ۱۰ - ۴ - ۳
    سرچاپ ۱۰۰ / ۱۰ / ۲
    سرچاپ ۲ * ۳ ** ۲
پایان
\end{salamfa}

\begin{outputfa}
14
20
3
5
18
\end{outputfa}

\scode{۲ + ۳ * ۴} برابر \code{۱۴} شد نه \code{۲۰}، چون ضرب زودتر انجام
می‌شود. \scode{۱۰ - ۴ - ۳} برابر \code{۳} است، چون اول \scode{۱۰ - ۴} حساب
می‌شود و بعد \scode{۶ - ۳}. هر جا شک دارید یا ترتیب دیگری
می‌خواهید، پرانتز بگذارید. پرانتز اضافی هرگز ضرری ندارد و اغلب خواندن را
راحت‌تر می‌کند.

\begin{tip}
عبارت‌های خیلی طولانی را به چند متغیر با نام‌های گویا بشکنید. به جای
\scode{سرچاپ (الف + ب) * ج / ۱۰۰ - د}،
بنویسید \scode{جمع := الف + ب}
و سپس \mbox{\scode{تخفیف := جمع * ج / ۱۰۰}} و در آخر
\scode{سرچاپ تخفیف - د}. خواننده (و خودتان یک هفته بعد) از شما ممنون
خواهد شد.
\end{tip}

\section{کارهایی که با تقسیم و باقی‌مانده می‌شود کرد}

تقسیم صحیح و باقی‌مانده، با هم، ابزاری قدرتمند برای «خرد کردن» عددها
هستند. چند نمونه:

\program{تبدیل ثانیه به ساعت و دقیقه}
\begin{salamfa}
روال ریشه:
    مجموع ثانیه := ۳۷۲۵
    ساعت := مجموع ثانیه / ۳۶۰۰
    دقیقه := (مجموع ثانیه % ۳۶۰۰) / ۶۰
    ثانیه := مجموع ثانیه % ۶۰
    سرچاپ ساعت, "ساعت و", دقیقه, "دقیقه و", ثانیه, "ثانیه"
پایان
\end{salamfa}

\begin{outputfa}
1 ساعت و 2 دقیقه و 5 ثانیه
\end{outputfa}

هر ساعت ۳۶۰۰ ثانیه است، پس تقسیم صحیح بر ۳۶۰۰ تعداد ساعت‌های کامل را
می‌دهد. باقی‌مانده‌ی همین تقسیم، ثانیه‌هایی است که یک ساعت کامل نساخته‌اند و
همان را بر ۶۰ تقسیم می‌کنیم تا دقیقه‌ها به دست بیاید.

\program{رقم‌های یک عدد}
\begin{salamfa}
روال ریشه:
    عدد := ۴۷۲
    سرچاپ "یکان:", عدد % ۱۰
    سرچاپ "دهگان:", (عدد / ۱۰) % ۱۰
    سرچاپ "صدگان:", عدد / ۱۰۰
پایان
\end{salamfa}

\begin{outputfa}
یکان: 2
دهگان: 7
صدگان: 4
\end{outputfa}

باقی‌مانده‌ی تقسیم بر ۱۰، همیشه رقم یکان است. تقسیم بر ۱۰ یک رقم را از
سمت راست می‌اندازد. با ترکیب این دو می‌توان به هر رقمی رسید.

\program{درصد و تخفیف}
\begin{salamfa}
روال ریشه:
    قیمت := ۲۵۰۰۰۰
    درصد تخفیف := ۱۵
    مبلغ تخفیف := (قیمت * درصد تخفیف) / ۱۰۰
    سرچاپ "تخفیف:", مبلغ تخفیف
    سرچاپ "قیمت نهایی:", قیمت - مبلغ تخفیف
پایان
\end{salamfa}

\begin{outputfa}
تخفیف: 37500
قیمت نهایی: 212500
\end{outputfa}

دقت کنید که اول ضرب کردیم و بعد تقسیم. اگر می‌نوشتیم
\scode{قیمت * (درصد تخفیف / ۱۰۰)}، تقسیم صحیح ۱۵ بر ۱۰۰ صفر می‌شد و کل
تخفیف از بین می‌رفت. با عددهای صحیح، تا جای ممکن تقسیم را آخر انجام دهید.

\section{شکل‌های کوتاه مقداردهی}

در فصل \ref{ch:variables} با \code{+=} و \code{++} آشنا شدید. کارور‌های حسابی دیگر هم شکل
کوتاه دارند:

\begin{center}
\begin{tabular}{ll}
\toprule
شکل کوتاه & معادل \\
\midrule
\scode{الف += ۵} & \scode{الف = الف + ۵} \\
\scode{الف -= ۵} & \scode{الف = الف - ۵} \\
\scode{الف *= ۵} & \scode{الف = الف * ۵} \\
\scode{الف /= ۵} & \scode{الف = الف / ۵} \\
\scode{الف \%= ۵} & \scode{الف = الف \% ۵} \\
\scode{الف++} & \scode{الف = الف + ۱} \\
\scode{الف--} & \scode{الف = الف - ۱} \\
\bottomrule
\end{tabular}
\end{center}

توان هم شکل کوتاه \code{**=} دارد، اما چون حاصل توان همیشه اعشاری است، آن را
فقط روی متغیری از گونه‌ی \fcode{اعشار۶۴} می‌توان به کار برد.

\program{شکل‌های کوتاه}
\begin{salamfa}
روال ریشه:
    ناپایا امتیاز := ۱۰۰
    امتیاز += ۲۰
    سرچاپ امتیاز
    امتیاز -= ۵۰
    سرچاپ امتیاز
    امتیاز *= ۲
    سرچاپ امتیاز
    امتیاز /= ۴
    سرچاپ امتیاز
    امتیاز++
    سرچاپ امتیاز
    امتیاز--
    سرچاپ امتیاز
پایان
\end{salamfa}

\begin{outputfa}
120
70
140
35
36
35
\end{outputfa}

\section{کارور‌های مقایسه}

کارور‌های مقایسه دو مقدار را با هم می‌سنجند و نتیجه‌ی منطقی می‌دهند:
\fcode{درست} یا \fcode{نادرست}.

\begin{center}
\begin{tabular}{cl}
\toprule
کارور & معنی \\
\midrule
\code{==} & برابر است با \\
\code{!=} & برابر نیست با \\
\code{<} & کوچک‌تر از \\
\code{>} & بزرگ‌تر از \\
\code{<=} & کوچک‌تر یا مساوی \\
\code{>=} & بزرگ‌تر یا مساوی \\
\bottomrule
\end{tabular}
\end{center}

\program{مقایسه}
\begin{salamfa}
روال ریشه:
    سن := ۲۰
    سرچاپ سن == ۲۰
    سرچاپ سن != ۲۰
    سرچاپ سن < ۱۸
    سرچاپ سن >= ۱۸
    سرچاپ "سیب" == "سیب"
    سرچاپ "سیب" == "گلابی"
پایان
\end{salamfa}

\begin{outputfa}
true
false
false
true
true
false
\end{outputfa}

\begin{tip}
به جای \code{==} و \code{!=} می‌توانید واژه‌های \fcode{برابر} و \fcode{نابرابر} را
هم بنویسید: \fcode{سرچاپ سن برابر ۲۰} درست همان کار
\scode{سرچاپ سن == ۲۰} را می‌کند. در این کتاب بیشتر همان علامت‌ها را به
کار می‌بریم.
\end{tip}

\begin{warn}
اشتباه \code{=} با \code{==} یکی از رایج‌ترین اشتباه‌های همه‌ی
برنامه‌نویسان جهان است. \code{=} مقدار می‌دهد، \code{==} می‌پرسد «آیا
برابرند؟». اگر جایی که باید مقایسه کنید \code{=} بنویسید، سلام خطا
می‌دهد، که خبر خوبی است.
\end{warn}

نتیجه‌ی مقایسه یک مقدار منطقی معمولی است و می‌توان آن را در متغیری نگه
داشت. این کار وقتی شرطی پیچیده است یا چند بار استفاده می‌شود، برنامه را
خواناتر می‌کند:

\begin{salamfa}
روال ریشه:
    دما := ۳۸
    هوا گرم است := دما > ۳۰
    سرچاپ "هوا گرم است؟", هوا گرم است
پایان
\end{salamfa}

\begin{outputfa}
هوا گرم است؟ true
\end{outputfa}

\section{ترکیب شرط‌ها: و، یا، نقیض}

بسیاری از تصمیم‌ها به بیش از یک شرط بستگی دارند: «اگر بزرگسال است \emph{و}
بلیت دارد»، «اگر جمعه است \emph{یا} تعطیل رسمی است». سلام برای ترکیب
شرط‌ها سه کارور دارد که هر سه واژه‌ی فارسی‌اند:

\begin{center}
\begin{tabular}{cll}
\toprule
کارور & معنی & نتیجه درست است وقتی که \\
\midrule
\fcode{و} & و & هر دو طرف درست باشند \\
\fcode{یا} & یا & دست‌کم یکی از دو طرف درست باشد \\
\fcode{وارونه} & نقیض (نه) & طرف راستش نادرست باشد \\
\bottomrule
\end{tabular}
\end{center}

\program{ترکیب شرط‌ها}
\begin{salamfa}
روال ریشه:
    سن := ۲۰
    بلیت دارد := نادرست
    سرچاپ (سن >= ۱۸) و بلیت دارد
    سرچاپ (سن >= ۱۸) یا بلیت دارد
    سرچاپ وارونه بلیت دارد
    سرچاپ (سن > ۱۲) و (سن < ۶۵)
پایان
\end{salamfa}

\begin{outputfa}
false
true
true
true
\end{outputfa}

خط آخر را ببینید: برای پرسیدن «آیا سن میان ۱۲ و ۶۵ است؟» باید دو مقایسه‌ی
جداگانه را با \fcode{و} به هم وصل کرد. نوشتن \scode{۱۲ < سن < ۶۵}، آن‌طور
که در ریاضی می‌نویسیم، در سلام معنی ندارد.

\begin{note}
علامت‌های \code{\&\&}، \code{||} و \code{!} که در بسیاری از زبان‌های دیگر به کار
می‌روند در سلام خطا هستند؛ همیشه \fcode{و}، \fcode{یا} و \fcode{وارونه} بنویسید.
\end{note}

\Needspace{16\baselineskip}
\begin{deep}[جدول درستی]
اگر بخواهید همه‌ی حالت‌های \fcode{و} و \fcode{یا} را یک‌جا ببینید،
جدول زیر کمکتان می‌کند. آن را حفظ نکنید؛ فقط با معنی روزمره‌ی «و» و «یا»
مقایسه‌اش کنید و می‌بینید که کاملاً همان است.
\begin{center}
\begin{tabular}{cccc}
\toprule
الف & ب & \fcode{الف و ب} & \fcode{الف یا ب} \\
\midrule
درست & درست & درست & درست \\
درست & نادرست & نادرست & درست \\
نادرست & درست & نادرست & درست \\
نادرست & نادرست & نادرست & نادرست \\
\bottomrule
\end{tabular}
\end{center}
تنها جایی که «یا»ی برنامه‌نویسی با «یا»ی روزمره فرق دارد، وقتی است که هر
دو طرف درست باشند: در گفت‌وگوی روزمره «چای یا قهوه؟» یعنی یکی از آن دو،
اما \fcode{یا} در سلام وقتی هر دو درست باشند هم درست است.
\end{deep}

\section{یک برنامه‌ی کامل‌تر}

فرض کنید می‌خواهید بدانید برای یک مهمانی چند کارتن نوشیدنی بخرید. هر
کارتن ۶ بطری دارد و هر مهمان ۲ بطری می‌نوشد:

\program{خرید برای مهمانی}
\begin{salamfa}
روال ریشه:
    تعداد مهمان := ۱۷
    بطری هرنفر := ۲
    بطری کارتن := ۶

    بطری لازم := تعداد مهمان * بطری هرنفر
    کارتن کامل := بطری لازم / بطری کارتن
    بطری اضافه := بطری لازم % بطری کارتن
    // اگر بطری اضافه داریم، یک کارتن دیگر لازم است
    کارتن لازم := کارتن کامل + (بطری اضافه > ۰ ؟ ۱ : ۰)

    سرچاپ "بطری لازم:", بطری لازم
    سرچاپ "کارتن لازم:", کارتن لازم
    سرچاپ "بطری باقی‌مانده:", (کارتن لازم * بطری کارتن) - بطری لازم
پایان
\end{salamfa}

\begin{outputfa}
بطری لازم: 34
کارتن لازم: 6
بطری باقی‌مانده: 2
\end{outputfa}

در این برنامه یک کارور تازه هم هست: \scode{شرط ؟ الف : ب}. این کارور
می‌گوید «اگر شرط درست است مقدار الف، وگرنه مقدار ب». این‌جا اگر بطری
اضافه‌ای مانده باشد، یک کارتن بیشتر می‌خریم. (شاید بپرسید چرا ننوشتیم
\fcode{بطری هر نفر}؛ پاسخ در فصل \ref{ch:variables} است: واژه‌ی \fcode{هر}
یکی از واژه‌های خود زبان است و نمی‌تواند بخشی از یک نام باشد.) در فصل بعد، که به تصمیم‌گیری
می‌پردازد، این کارور را بیشتر می‌بینیم.

\begin{summary}
\begin{itemize}
  \item کارور‌های حسابی: \code{+ - * / \%} و توان
        \code{**}. تقسیم صحیح اعشار را
        دور می‌ریزد و \code{\%} باقی‌مانده را می‌دهد.
  \item ترتیب: اول توان، بعد ضرب و تقسیم، بعد جمع و تفریق. با پرانتز
        ترتیب را عوض کنید.
  \item شکل‌های کوتاه: \code{+=}، \code{-=}، \code{*=}، \code{/=}،
        \code{\%=}، \code{++}، \code{--}.
  \item مقایسه: \code{== != < > <= >=} که نتیجه‌ی منطقی می‌دهند.
  \item ترکیب شرط‌ها با \fcode{و}، \fcode{یا} و \fcode{وارونه}.
  \item \scode{شرط ؟ الف : ب} یکی از دو مقدار را بر پایه‌ی شرط انتخاب
        می‌کند.
\end{itemize}
\end{summary}

\section*{تمرین‌ها}
\markright{تمرین‌ها}

\exercise بدون اجرا کردن، نتیجه‌ی هر عبارت را بگویید و سپس بررسی کنید:
\scode{۵ + ۲ * ۳}، \scode{(۵ + ۲) * ۳}، \scode{۲۰ / ۳ * ۳}، \scode{۲۰ \% ۳}،
\scode{۲ ** ۲ ** ۳}.

\exercise عددی چهاررقمی را در متغیری بگذارید و مجموع رقم‌هایش را با تقسیم
و باقی‌مانده حساب و چاپ کنید.

\exercise قیمت یک کالا و مبلغی که مشتری پرداخته را در دو متغیر بگذارید و
باقی‌مانده را به کمترین تعداد اسکناس ۱۰٬۰۰۰ تومانی، ۵٬۰۰۰ تومانی و
۱٬۰۰۰ تومانی تقسیم کنید و تعداد هر یک را چاپ کنید.

\exercise سال تولد کسی را در متغیری بگذارید و با یک عبارت منطقی چاپ کنید
که آیا او امسال (۱۴۰۵) به سن رأی دادن (۱۸ سال) رسیده است یا نه.

\exercise سه متغیر منطقی به نام‌های \fcode{باران می‌بارد}،
\fcode{چتر دارم} و \fcode{ماشین دارم} بسازید و عبارتی بنویسید که بگوید
آیا «بدون خیس شدن» می‌توان بیرون رفت: یعنی یا باران نمی‌بارد، یا چتر یا
ماشین دارم. برنامه را با چند ترکیب مختلف از مقدارها امتحان کنید.

\chapter{تصمیم‌گیری}\label{ch:conditions}

\begin{goals}
\begin{itemize}
  \item اجرای دستورها فقط وقتی شرطی برقرار است، با \fcode{اگر}
  \item انتخاب میان دو راه با \fcode{وگرنه} و میان چند راه با
        \fcode{وگرنه شرط}
  \item شرط‌های تودرتو و ترکیب شرط‌ها
  \item انتخاب مقدار با کارور \scode{؟ :}
  \item انتخاب میان حالت‌های مشخص با \fcode{همخوان}
\end{itemize}
\end{goals}

\section{برنامه‌ای که تصمیم می‌گیرد}

همه‌ی برنامه‌هایی که تا این‌جا نوشتیم یک راه بیشتر نداشتند: از خط اول شروع
می‌کردند و خط به خط تا آخر می‌رفتند. اما برنامه‌های واقعی باید بتوانند
تصمیم بگیرند. «اگر نمره ده یا بیشتر بود، بنویس قبول.» «اگر رمز درست بود،
وارد شو.» «اگر موجودی کافی نبود، هشدار بده.» این «اگر»ها دقیقاً همان
واژه‌ای هستند که در سلام به کار می‌رود.

\section{اگر}

ساده‌ترین شکل تصمیم‌گیری این است: اگر شرطی برقرار بود، یک یا چند دستور
اجرا شود، وگرنه هیچ:

\program{ساده‌ترین اگر}
\begin{salamfa}
روال ریشه:
    نمره := ۱۴
    اگر نمره >= ۱۰:
        سرچاپ "قبول"
    پایان
    سرچاپ "پایان برنامه"
پایان
\end{salamfa}

\begin{outputfa}
قبول
پایان برنامه
\end{outputfa}

ساختار \fcode{اگر} همان الگویی است که در \fcode{روال ریشه} دیدید: یک خط که
با دونقطه تمام می‌شود، چند دستور با تورفتگی، و \fcode{پایان}. آنچه بین
\fcode{اگر} و دونقطه می‌آید، \textbf{شرط} است: یک عبارت منطقی که یا درست
است یا نادرست. اگر درست باشد، دستورهای درون بلوک اجرا می‌شوند؛ اگر نادرست
باشد، سلام از روی کل بلوک می‌پرد و از خط بعد از \fcode{پایان} ادامه می‌دهد.

نمره را در برنامه‌ی بالا به \code{۸} تغییر دهید و دوباره اجرا کنید. این بار
فقط «پایان برنامه» چاپ می‌شود، چون شرط برقرار نبود و دستور \fcode{سرچاپ
"قبول"} اصلاً اجرا نشد.

\begin{note}
شرط \fcode{اگر} می‌تواند هر چیزی باشد که یک مقدار منطقی می‌دهد: یک مقایسه
مثل \scode{نمره >= ۱۰}، یک متغیر منطقی مثل \fcode{بلیت دارد}، یا
ترکیبی از این‌ها با \fcode{و} و \fcode{یا}. عدد به تنهایی شرط نیست؛ نوشتن
\scode{اگر نمره:} خطاست.
\end{note}

\section{اگر و وگرنه}

بسیاری از تصمیم‌ها دو راه دارند: یا این، یا آن. واژه‌ی \kwd{وگرنه} راه دوم
را مشخص می‌کند:

\program{اگر و وگرنه}
\begin{salamfa}
روال ریشه:
    نمره := ۸
    اگر نمره >= ۱۰:
        سرچاپ "قبول"
    وگرنه:
        سرچاپ "مردود"
    پایان
پایان
\end{salamfa}

\begin{outputfa}
مردود
\end{outputfa}

دقت کنید که \fcode{وگرنه} خودش دونقطه دارد، اما \fcode{پایان} جداگانه ندارد؛
یک \fcode{پایان} در پایان، کل ساختار «اگر، وگرنه» را می‌بندد. از دو بلوک،
همیشه دقیقاً یکی اجرا می‌شود: نه هر دو، نه هیچ‌کدام.

\program{زوج یا فرد}
\begin{salamfa}
روال ریشه:
    عدد := ۲۷
    اگر (عدد % ۲) == ۰:
        سرچاپ عدد, "زوج است"
    وگرنه:
        سرچاپ عدد, "فرد است"
    پایان
پایان
\end{salamfa}

\begin{outputfa}
27 فرد است
\end{outputfa}

عدد زوج آن است که باقی‌مانده‌ی تقسیمش بر ۲ صفر باشد. این یکی از
پرکاربردترین شرط‌ها در برنامه‌نویسی است و بارها آن را خواهید دید.

\section{بیش از دو راه: وگرنه شرط}

گاهی راه‌ها بیش از دو تاست. مثلاً می‌خواهیم نمره را به «عالی»، «خوب»،
«قابل قبول» و «مردود» تقسیم کنیم. می‌توان بعد از \fcode{وگرنه} یک شرط تازه
نوشت:

\program{زنجیره‌ی شرط‌ها}
\begin{salamfa}
روال ریشه:
    نمره := ۱۷
    اگر نمره >= ۱۸:
        سرچاپ "عالی"
    وگرنه نمره >= ۱۵:
        سرچاپ "خوب"
    وگرنه نمره >= ۱۰:
        سرچاپ "قابل قبول"
    وگرنه:
        سرچاپ "مردود"
    پایان
پایان
\end{salamfa}

\begin{outputfa}
خوب
\end{outputfa}

سلام شرط‌ها را از بالا به پایین می‌سنجد و به \emph{اولین} شرطی که درست بود
بسنده می‌کند؛ بقیه اصلاً بررسی نمی‌شوند. نمره‌ی ۱۷ از ۱۸ کمتر است، پس شرط
اول نادرست است؛ اما از ۱۵ بیشتر است، پس «خوب» چاپ می‌شود و سلام از بقیه‌ی
شرط‌ها می‌گذرد. \fcode{وگرنه}‌ی آخر که شرط ندارد، وقتی اجرا می‌شود که
هیچ‌کدام از شرط‌ها برقرار نباشند. نوشتن آن اختیاری است.

\begin{warn}
ترتیب شرط‌ها در زنجیره مهم است. اگر برنامه‌ی بالا را با شرط
\scode{نمره >= ۱۰} شروع می‌کردیم، نمره‌ی ۱۷ هم «قابل قبول» می‌شد و
هیچ‌وقت به شرط‌های بعدی نمی‌رسید. شرط‌ها را از خاص‌ترین به عام‌ترین
بنویسید.
\end{warn}

\section{اگر در یک خط}

وقتی بلوک \fcode{اگر} فقط یک دستور کوتاه دارد، می‌توان همه را در یک خط
نوشت. \fcode{پایان} همچنان لازم است:

\program{اگر تک‌خطی}
\begin{salamfa}
روال ریشه:
    دما := ۳۵
    اگر دما > ۳۰: سرچاپ "هوا گرم است" پایان
    اگر دما < ۱۰: سرچاپ "هوا سرد است" پایان
    اگر (دما >= ۱۰) و (دما <= ۳۰): سرچاپ "هوا معتدل است" پایان
پایان
\end{salamfa}

\begin{outputfa}
هوا گرم است
\end{outputfa}

این شکل برای شرط‌های کوتاه و ساده مناسب است. به محض این که بلوک بیش از
یک دستور داشت یا شرط طولانی شد، به شکل چندخطی برگردید.

\section{شرط درون شرط}

درون بلوک \fcode{اگر} می‌توان هر دستوری نوشت، از جمله یک \fcode{اگر} دیگر.
به این کار \textbf{تودرتو} کردن می‌گویند:

\program{شرط تودرتو}
\begin{salamfa}
روال ریشه:
    سن := ۲۰
    بلیت دارد := درست
    اگر سن >= ۱۸:
        اگر بلیت دارد:
            سرچاپ "خوش آمدید"
        وگرنه:
            سرچاپ "لطفاً بلیت تهیه کنید"
        پایان
    وگرنه:
        سرچاپ "ورود برای زیر ۱۸ سال ممکن نیست"
    پایان
پایان
\end{salamfa}

\begin{outputfa}
خوش آمدید
\end{outputfa}

به تورفتگی‌ها دقت کنید: \fcode{اگر} درونی یک پله جلوتر است و \fcode{پایان}
هر بلوک، دقیقاً زیر \fcode{اگر} خودش قرار دارد. بدون تورفتگی درست، فهمیدن
این که کدام \fcode{پایان} کدام بلوک را می‌بندد، تقریباً ناممکن می‌شود.

بسیاری از شرط‌های تودرتو را می‌توان با \fcode{و} و \fcode{یا} صاف کرد.
همان برنامه، بی‌تودرتویی:

\begin{salamfa}
روال ریشه:
    سن := ۲۰
    بلیت دارد := درست
    اگر (سن >= ۱۸) و بلیت دارد:
        سرچاپ "خوش آمدید"
    وگرنه سن < ۱۸:
        سرچاپ "ورود برای زیر ۱۸ سال ممکن نیست"
    وگرنه:
        سرچاپ "لطفاً بلیت تهیه کنید"
    پایان
پایان
\end{salamfa}

\begin{outputfa}
خوش آمدید
\end{outputfa}

هیچ‌کدام از دو شکل «درست‌تر» نیست. قاعده‌ی سرانگشتی: هر کدام را که با یک
بار خواندن بهتر می‌فهمید، همان را بنویسید.

\section{بزرگ‌ترین از میان چند عدد}

یک الگوی رایج که با \fcode{اگر} و یک متغیر تغییرپذیر ساخته می‌شود: پیدا
کردن بزرگ‌ترین مقدار. اول فرض می‌کنیم اولی بزرگ‌ترین است، بعد بقیه را با آن
می‌سنجیم:

\program{بزرگ‌ترین سه عدد}
\begin{salamfa}
روال ریشه:
    الف := ۱۲
    ب := ۲۵
    ج := ۱۹
    ناپایا بزرگترین := الف
    اگر ب > بزرگترین: بزرگترین = ب پایان
    اگر ج > بزرگترین: بزرگترین = ج پایان
    سرچاپ "بزرگ‌ترین:", بزرگترین
پایان
\end{salamfa}

\begin{outputfa}
بزرگ‌ترین: 25
\end{outputfa}

با هر مقداری برای \fcode{الف}، \fcode{ب} و \fcode{ج} امتحان کنید. این
الگو در فصل‌های~\ref{ch:loops} و~\ref{ch:arrays}، وقتی به جای سه عدد صدها عدد
داشته باشیم، دوباره به کارمان می‌آید.

\section{انتخاب مقدار: کارور ؟}

گاهی نمی‌خواهیم دستوری اجرا شود، بلکه می‌خواهیم بر پایه‌ی شرطی، یکی از دو
\emph{مقدار} را انتخاب کنیم. در فصل پیش کارور \scode{شرط ؟ الف : ب} را
دیدید که دقیقاً همین کار را می‌کند:

\program{کارور انتخاب}
\begin{salamfa}
روال ریشه:
    سن := ۱۵
    قیمت بلیت := سن < ۱۲ ؟ ۵۰۰۰۰ : ۱۲۰۰۰۰
    سرچاپ "قیمت بلیت:", قیمت بلیت
    سرچاپ سن >= ۱۸ ؟ "بزرگسال" : "نوجوان"
پایان
\end{salamfa}

\begin{outputfa}
قیمت بلیت: 120000
نوجوان
\end{outputfa}

خط \scode{قیمت بلیت := سن < ۱۲ ؟ ۵۰۰۰۰ : ۱۲۰۰۰۰} را بخوانید:
«قیمت بلیت برابر است با: اگر سن کمتر از ۱۲ است، ۵۰ هزار، وگرنه ۱۲۰ هزار».
همین کار را می‌شد با یک \fcode{اگر} و \fcode{وگرنه} و یک متغیر تغییرپذیر
کرد، اما این شکل کوتاه‌تر است. علامت پرسش می‌تواند فارسی (\fcode{؟}) یا
انگلیسی (\code{?}) باشد.

\begin{tip}
کارور \scode{؟ :} را فقط برای انتخاب‌های ساده به کار ببرید. به محض این که
دو یا سه تای آن‌ها تودرتو شدند، به \fcode{اگر} برگردید؛ خوانایی مهم‌تر از
کوتاهی است.
\end{tip}

\section{همخوان: انتخاب میان حالت‌های مشخص}

وقتی یک مقدار را با چند حالت \emph{مشخص} مقایسه می‌کنیم (شماره‌ی روز
هفته، شماره‌ی ماه، کد یک گزینه)، زنجیره‌ی \fcode{اگر} و \fcode{وگرنه}
طولانی و یکنواخت می‌شود. سلام برای این کار واژه‌ی \kwd{همخوان} را دارد:

\program{نام روز هفته}
\begin{salamfa}
روال ریشه:
    روز := ۵
    همخوان روز:
        ۱: سرچاپ "شنبه" پایان
        ۲: سرچاپ "یکشنبه" پایان
        ۳: سرچاپ "دوشنبه" پایان
        ۴: سرچاپ "سه‌شنبه" پایان
        ۵: سرچاپ "چهارشنبه" پایان
        ۶: سرچاپ "پنجشنبه" پایان
        ۷: سرچاپ "جمعه" پایان
        وگرنه: سرچاپ "شماره‌ی روز نادرست است" پایان
    پایان
پایان
\end{salamfa}

\begin{outputfa}
چهارشنبه
\end{outputfa}

\fcode{همخوان} مقدار \fcode{روز} را با هر یک از حالت‌ها مقایسه می‌کند و
بلوک همان حالتی را که برابر بود اجرا می‌کند. هر حالت یک بلوک کوچک است که
با دونقطه شروع و با \fcode{پایان} بسته می‌شود؛ می‌توان آن را در یک خط یا
چند خط نوشت. \fcode{وگرنه} برای وقتی است که هیچ حالتی برابر نباشد.

\fcode{همخوان} را جایی هم می‌توان به کار برد که یک \emph{مقدار} لازم است،
نه اجرای یک دستور. حالت‌ها دقیقاً به همان شکل نوشته می‌شوند و مقدار هر
حالت همان عبارتی است که بلوکش می‌دهد. چند مقدار می‌توانند با ویرگول یک
حالت مشترک باشند:

\program{فصل هر ماه}
\begin{salamfa}
روال ریشه:
    ماه := ۴
    فصل := همخوان ماه:
        ۱, ۲, ۳: "بهار" پایان
        ۴, ۵, ۶: "تابستان" پایان
        ۷, ۸, ۹: "پاییز" پایان
        ۱۰, ۱۱, ۱۲: "زمستان" پایان
        وگرنه: "نامعتبر" پایان
    پایان
    سرچاپ "ماه", ماه, "در فصل", فصل, "است"
پایان
\end{salamfa}

\begin{outputfa}
ماه 4 در فصل تابستان است
\end{outputfa}

این شکل برای «جدول‌های تبدیل» (عدد به نام، کد به پیام) بسیار خوش‌دست است.
اگر چنین جدولی را با \fcode{اگر} می‌نوشتیم، دوازده شرط و دوازده
\fcode{سرچاپ} لازم بود.

\section{یک برنامه‌ی کامل‌تر}

شاخص توده‌ی بدنی از تقسیم وزن (کیلوگرم) بر مجذور قد (متر) به دست می‌آید.
این برنامه شاخص را حساب و دسته‌بندی می‌کند:

\program{شاخص توده‌ی بدنی}
\begin{salamfa}
روال ریشه:
    وزن := ۷۰.۰
    قد := ۱.۷۵
    شاخص := وزن / (قد * قد)
    سرچاپ "شاخص توده‌ی بدنی:", شاخص

    اگر شاخص < ۱۸.۵:
        سرچاپ "وضعیت: کم‌وزن"
    وگرنه شاخص < ۲۵.۰:
        سرچاپ "وضعیت: طبیعی"
    وگرنه شاخص < ۳۰.۰:
        سرچاپ "وضعیت: اضافه‌وزن"
    وگرنه:
        سرچاپ "وضعیت: چاق"
    پایان
پایان
\end{salamfa}

\begin{outputfa}
شاخص توده‌ی بدنی: 22.857142857142858
وضعیت: طبیعی
\end{outputfa}

به این نکته دقت کنید که در زنجیره‌ی شرط‌ها لازم نبود بنویسیم
\begin{center}\scode{شاخص >= ۱۸.۵ و شاخص < ۲۵.۰}\end{center}
چون اگر به شرط دوم رسیده‌ایم، یعنی شرط اول نادرست بوده و شاخص حتماً از
۱۸٫۵ کمتر نیست. زنجیره، خودش این را تضمین می‌کند.

\begin{summary}
\begin{itemize}
  \item \scode{اگر شرط:} بلوکی را فقط وقتی اجرا می‌کند که شرط درست باشد.
        \scode{وگرنه:} راه دوم را می‌دهد و \scode{وگرنه شرط:} راه‌های
        بیشتر. یک \fcode{پایان} کل ساختار را می‌بندد.
  \item در زنجیره، اولین شرط درست برنده است؛ ترتیب مهم است.
  \item بلوک‌های کوتاه را می‌توان در یک خط نوشت، با \fcode{پایان} در پایان.
  \item \fcode{اگر} می‌تواند درون \fcode{اگر} باشد؛ تورفتگی را رعایت کنید.
        بسیاری از تودرتویی‌ها با \fcode{و} و \fcode{یا} ساده می‌شوند.
  \item \scode{شرط ؟ الف : ب} یکی از دو مقدار را انتخاب می‌کند.
  \item \fcode{همخوان} برای مقایسه‌ی یک مقدار با چند حالت مشخص است؛ هر
        حالت بلوکی است که با دونقطه باز و با \fcode{پایان} بسته می‌شود و
        هم به شکل دستور و هم به شکل مقدار به کار می‌آید.
\end{itemize}
\end{summary}

\section*{تمرین‌ها}
\markright{تمرین‌ها}

\exercise عددی را در متغیری بگذارید و چاپ کنید که مثبت، منفی یا صفر است.

\exercise سه ضلع یک مثلث را در سه متغیر بگذارید و چاپ کنید که مثلث
متساوی‌الاضلاع (هر سه ضلع برابر)، متساوی‌الساقین (دو ضلع برابر) یا مختلف‌الاضلاع
است.

\exercise ساعت روز (۰ تا ۲۳) را در متغیری بگذارید و پیام مناسب چاپ کنید:
پیش از ۱۲ «صبح بخیر»، از ۱۲ تا پیش از ۱۷ «ظهر بخیر»، از ۱۷ تا پیش از ۲۰
«عصر بخیر» و از ۲۰ به بعد «شب بخیر».

\exercise با \fcode{همخوان}، شماره‌ی یک ماه را به تعداد روزهایش تبدیل کنید:
شش ماه اول ۳۱ روز، پنج ماه بعد ۳۰ روز و اسفند ۲۹ روز.

\exercise وزن یک بسته‌ی پستی را در متغیری بگذارید و هزینه‌ی ارسال را حساب
کنید: تا ۱ کیلوگرم ۳۰ هزار تومان، تا ۵ کیلوگرم ۵۰ هزار تومان و بیش از آن،
۵۰ هزار تومان به‌علاوه‌ی ۱۵ هزار تومان برای هر کیلوگرمِ بیش از ۵. اگر بسته به شهر دیگری می‌رود (یک متغیر
منطقی)، ۲۰ درصد به هزینه اضافه شود.

\exercise برنامه‌ی زیر چه چاپ می‌کند؟ اول حدس بزنید، بعد اجرا کنید و اگر
حدس شما اشتباه بود، دلیلش را پیدا کنید:
\begin{salamfa}
روال ریشه:
    عدد := ۱۵
    اگر عدد > ۱۰:
        سرچاپ "بزرگ‌تر از ده"
    وگرنه عدد > ۵:
        سرچاپ "بزرگ‌تر از پنج"
    پایان
    اگر (عدد % ۵) == ۰: سرچاپ "مضرب پنج" پایان
    اگر (عدد % ۳) == ۰: سرچاپ "مضرب سه" پایان
پایان
\end{salamfa}

\chapter{حلقه‌ها}
\label{ch:loops}

\begin{goals}
\begin{itemize}
  \item تکرار یک کار به تعداد مشخص با \fcode{تکرار}
  \item شمارش از یک عدد تا عدد دیگر، رو به بالا، رو به پایین و با گام
  \item تکرار تا وقتی که شرطی برقرار است، با \fcode{تا}
  \item شکستن حلقه با \fcode{بشکن} و پریدن از یک دور با \fcode{گذر}
  \item حلقه درون حلقه
\end{itemize}
\end{goals}

\section{چرا حلقه؟}

فرض کنید می‌خواهید عددهای ۱ تا ۵ را چاپ کنید. با آنچه تا حالا یاد
گرفته‌اید، پنج دستور \fcode{سرچاپ} می‌نویسید. حالا اگر بخواهید تا ۱۰۰۰ را چاپ
کنید چه؟ یا اگر تعداد از پیش معلوم نباشد؟ رایانه‌ها دقیقاً برای همین
کارهای تکراری ساخته شده‌اند، و \textbf{حلقه} ابزاری است که به آن‌ها
می‌گوید: «این کار را بارها انجام بده.»

\section{تکرار به تعداد مشخص}

ساده‌ترین حلقه‌ی سلام \kwd{تکرار} است. جلویش می‌نویسید چند بار، و سلام بلوک
را همان تعداد بار اجرا می‌کند:

\program{ساده‌ترین حلقه}
\begin{salamfa}
روال ریشه:
    تکرار ۳:
        سرچاپ "سلام!"
    پایان
    سرچاپ "پایان"
پایان
\end{salamfa}

\begin{outputfa}
سلام!
سلام!
سلام!
پایان
\end{outputfa}

ساختار حلقه همان ساختار آشناست: خطی که با دونقطه تمام می‌شود، بلوک با
تورفتگی و \fcode{پایان}. به هر بار اجرای بلوک یک \textbf{دور} می‌گوییم. این
حلقه سه دور دارد. عدد جلوی \fcode{تکرار} می‌تواند یک متغیر یا محاسبه هم
باشد: \fcode{تکرار تعداد مهمان:} یا \fcode{تکرار ن * ۲:}.

\section{شمارنده‌ی حلقه}

بیشتر وقت‌ها می‌خواهیم بدانیم در کدام دور هستیم. با نوشتن \kwd{در} و یک
نام، سلام در هر دور شماره‌ی آن دور را در یک متغیر می‌گذارد. شماره از
\emph{صفر} شروع می‌شود:

\program{شمارنده‌ی حلقه}
\begin{salamfa}
روال ریشه:
    تکرار ۵ در شماره:
        سرچاپ "دور", شماره
    پایان
پایان
\end{salamfa}

\begin{outputfa}
دور 0
دور 1
دور 2
دور 3
دور 4
\end{outputfa}

پنج دور، با شماره‌های ۰ تا ۴. شروع از صفر شاید عجیب باشد، اما در
برنامه‌نویسی یک رسم بسیار قدیمی است و در فصل \ref{ch:arrays} خواهید دید که
چرا مفید است. متغیر \fcode{شماره} فقط درون حلقه وجود دارد و بعد از
\fcode{پایان} دیگر در دسترس نیست.

\section{شمارش از یک عدد تا عدد دیگر}

اگر بخواهید از ۱ بشمارید نه از ۰، یا از ۱۰ تا ۲۰، از شکل \fcode{تکرار آغاز
تا پایان} استفاده کنید. هر دو سر بازه شمرده می‌شوند:

\program{شمارش از یک تا پنج}
\begin{salamfa}
روال ریشه:
    تکرار ۱ تا ۵ در شماره:
        سرچاپ شماره, "ضرب در ۳ می‌شود", شماره * ۳
    پایان
پایان
\end{salamfa}

\begin{outputfa}
1 ضرب در ۳ می‌شود 3
2 ضرب در ۳ می‌شود 6
3 ضرب در ۳ می‌شود 9
4 ضرب در ۳ می‌شود 12
5 ضرب در ۳ می‌شود 15
\end{outputfa}

اگر عدد آغاز از پایان بزرگ‌تر باشد، سلام رو به پایین می‌شمارد:

\program{شمارش معکوس}
\begin{salamfa}
روال ریشه:
    تکرار ۵ تا ۱ در ثانیه:
        سرچاپ ثانیه
    پایان
    سرچاپ "پرتاب!"
پایان
\end{salamfa}

\begin{outputfa}
5
4
3
2
1
پرتاب!
\end{outputfa}

و با \kwd{هر} می‌توان به جای یکی‌یکی، چندتاچندتا شمرد:

\program{شمارش با گام}
\begin{salamfa}
روال ریشه:
    ناپایا خط := ""
    تکرار ۰ تا ۲۰ هر ۵ در عدد:
        خط = خط + عدد + " "
    پایان
    سرچاپ خط
پایان
\end{salamfa}

\begin{outputfa}
0 5 10 15 20
\end{outputfa}

در این برنامه به جای چاپ هر عدد در یک خط، همه را در یک رشته کنار هم
چیدیم و در آخر یک بار چاپ کردیم. این ترفند (شروع با رشته‌ی خالی و چسباندن
قطعه‌ها) را بارها خواهید دید. گام باید عددی مثبت باشد؛ برای شمارش معکوس با
گام، کافی است آغاز را بزرگ‌تر از پایان بگذارید.

\section{جمع کردن در حلقه}

حالا الگوی «جمع کردن» فصل \ref{ch:variables} را با حلقه ترکیب کنیم. یک
متغیر تغییرپذیر با مقدار صفر می‌سازیم و در هر دور چیزی به آن اضافه
می‌کنیم:

\program{مجموع یک تا صد}
\begin{salamfa}
روال ریشه:
    ناپایا مجموع := ۰
    تکرار ۱ تا ۱۰۰ در عدد:
        مجموع += عدد
    پایان
    سرچاپ "مجموع ۱ تا ۱۰۰:", مجموع
پایان
\end{salamfa}

\begin{outputfa}
مجموع ۱ تا ۱۰۰: 5050
\end{outputfa}

دقت کنید که \fcode{مجموع} حتماً باید \fcode{ناپایا} باشد، چون در هر دور
عوض می‌شود، و حتماً باید \emph{بیرون} حلقه تعریف شود؛ اگر درون حلقه
تعریفش می‌کردیم، در هر دور یک جعبه‌ی تازه با مقدار صفر ساخته می‌شد و
هیچ‌وقت چیزی جمع نمی‌شد.

حلقه و شرط با هم، ترکیب بسیار قدرتمندی می‌سازند. مثلاً جمع و شمارش فقط
عددهای زوج:

\program{شمارش و جمع زوج‌ها}
\begin{salamfa}
روال ریشه:
    ناپایا مجموع زوج := ۰
    ناپایا تعداد زوج := ۰
    تکرار ۱ تا ۲۰ در عدد:
        اگر (عدد % ۲) == ۰:
            مجموع زوج += عدد
            تعداد زوج++
        پایان
    پایان
    سرچاپ "تعداد:", تعداد زوج, "مجموع:", مجموع زوج
پایان
\end{salamfa}

\begin{outputfa}
تعداد: 10 مجموع: 110
\end{outputfa}

\section{تکرار تا وقتی که: تا}

همه‌ی حلقه‌های بالا از قبل می‌دانستند چند دور دارند. اما گاهی نمی‌دانیم:
«آن‌قدر پس‌انداز کن تا به یک میلیون برسی.» «آن‌قدر عدد را بر ده تقسیم کن تا
صفر شود.» برای این حلقه‌ها واژه‌ی \kwd{تا} هست. جلویش یک شرط می‌آید و
تا وقتی شرط درست است، بلوک تکرار می‌شود:

\program{پس‌انداز تا رسیدن به هدف}
\begin{salamfa}
روال ریشه:
    ناپایا پس انداز := ۱۰۰۰۰۰
    ناپایا ماه := ۰
    تا پس انداز < ۱۰۰۰۰۰۰:
        پس انداز += ۱۵۰۰۰۰
        ماه++
    پایان
    سرچاپ "بعد از", ماه, "ماه پس‌انداز به", پس انداز, "رسید"
پایان
\end{salamfa}

\begin{outputfa}
بعد از 6 ماه پس‌انداز به 1000000 رسید
\end{outputfa}

سلام پیش از \emph{هر} دور شرط را می‌سنجد. تا وقتی پس‌انداز از یک میلیون
کمتر است، ۱۵۰ هزار به آن اضافه می‌کند و یک ماه می‌شمارد. به محض این که شرط
نادرست شد، حلقه تمام می‌شود و برنامه از بعد از \fcode{پایان} ادامه می‌دهد.
اگر شرط از همان اول نادرست باشد، بلوک حتی یک بار هم اجرا نمی‌شود.

\program{شمردن رقم‌های یک عدد}
\begin{salamfa}
روال ریشه:
    عدد := ۴۸۲۱۳
    ناپایا باقی := عدد
    ناپایا تعداد رقم := ۰
    تا باقی > ۰:
        باقی /= ۱۰
        تعداد رقم++
    پایان
    سرچاپ عدد, "دارای", تعداد رقم, "رقم است"
پایان
\end{salamfa}

\begin{outputfa}
48213 دارای 5 رقم است
\end{outputfa}

هر تقسیم صحیح بر ۱۰، یک رقم از سمت راست عدد می‌اندازد: ۴۸۲۱۳، ۴۸۲۱، ۴۸۲،
۴۸، ۴، ۰. پنج بار طول کشید تا به صفر برسد، پس عدد پنج رقم دارد. متغیر
\fcode{باقی} را برای این ساختیم که خود \fcode{عدد} دست‌نخورده بماند و در
آخر بتوانیم چاپش کنیم.

\begin{warn}[حلقه‌ی بی‌پایان]
در حلقه‌ی \fcode{تا}، چیزی درون بلوک باید شرط را سرانجام نادرست کند؛
وگرنه حلقه تا ابد می‌چرخد. به این وضعیت \textbf{حلقه‌ی بی‌پایان} می‌گویند و
یکی از رایج‌ترین اشتباه‌های تازه‌کارهاست. سلام در این راه دو بار به کمکتان
می‌آید. اگر خط \fcode{باقی \lr{/=} ۱۰} را از برنامه‌ی بالا حذف کنید، دیگر
هیچ‌جا \fcode{باقی} عوض نمی‌شود و سلام پیش از اجرا خطای \fcode{خطا۰۶۹}
می‌دهد: «متغیر 'باقی' با کلیدواژه‌ی 'ناپایا' تعریف شده اما هرگز تغییر نکرده
است». اما اگر آن را به اشتباه \fcode{باقی \lr{+=} ۱} بنویسید، برنامه اجرا
می‌شود و \fcode{باقی} هیچ‌گاه به صفر نمی‌رسد. ویرایشگر برخط چنین برنامه‌ای
را پس از سه ثانیه متوقف می‌کند و می‌نویسد:
«خطای زمان اجرا: زمان اجرا به پایان رسید (احتمال حلقه‌ی بی‌نهایت)».
در این حالت به دنبال خطی بگردید که باید شرط را عوض می‌کرده است.
\end{warn}

\begin{note}[کدام حلقه؟]
اگر تعداد دورها از قبل معلوم است (یا از روی یک عدد به دست می‌آید)،
\fcode{تکرار} انتخاب طبیعی است. اگر تعداد دورها به چیزی بستگی دارد که در
طول حلقه رخ می‌دهد، \fcode{تا} را برگزینید. هر کاری را با هر دو می‌شود
کرد، اما یکی از آن دو همیشه خواناتر است.
\end{note}

\section{شکستن حلقه و پریدن از یک دور}

گاهی در میانه‌ی حلقه می‌فهمیم که دیگر لازم نیست ادامه دهیم. دستور
\kwd{بشکن} حلقه را همان لحظه تمام می‌کند:

\program{پیدا کردن و متوقف شدن}
\begin{salamfa}
روال ریشه:
    تکرار ۱ تا ۱۰ در عدد:
        اگر (عدد * عدد) > ۳۰:
            سرچاپ "اولین عددی که مربعش از ۳۰ بیشتر است:", عدد
            بشکن
        پایان
    پایان
پایان
\end{salamfa}

\begin{outputfa}
اولین عددی که مربعش از ۳۰ بیشتر است: 6
\end{outputfa}

بدون \fcode{بشکن}، حلقه تا ۱۰ ادامه می‌داد و برای ۷، ۸، ۹ و ۱۰ هم پیام
چاپ می‌کرد. \fcode{بشکن} گفت: «پیدا شد، بس است.»

دستور \kwd{گذر} ملایم‌تر است: فقط از \emph{همین دور} می‌گذرد و به دور
بعد می‌رود:

\program{گذشتن از مضرب‌های سه}
\begin{salamfa}
روال ریشه:
    تکرار ۱ تا ۱۰ در عدد:
        اگر (عدد % ۳) == ۰: گذر پایان
        چاپ عدد, ""
    پایان
    سرچاپ
پایان
\end{salamfa}

\begin{outputfa}
1 2 4 5 7 8 10
\end{outputfa}

عددهای ۳، ۶ و ۹ چاپ نشدند، چون برای آن‌ها \fcode{گذر} اجرا شد و
\fcode{چاپ} هرگز به نوبت نرسید. (ترفند کوچک: \fcode{چاپ عدد, ""}
عدد را چاپ می‌کند و بعدش یک فاصله، چون میان چیزهایی که با ویرگول جدا
شده‌اند یک فاصله می‌گذارد. \fcode{سرچاپ} خالی در آخر هم خط را می‌بندد.)

\section{حلقه درون حلقه}

مثل شرط، حلقه هم می‌تواند درون حلقه‌ی دیگری باشد. حلقه‌ی درونی برای هر دور
حلقه‌ی بیرونی، \emph{کامل} اجرا می‌شود. کلاسیک‌ترین نمونه، جدول ضرب است:

\program{جدول ضرب}
\begin{salamfa}
روال ریشه:
    تکرار ۱ تا ۵ در سطر:
        ناپایا خط := ""
        تکرار ۱ تا ۵ در ستون:
            خط = خط + (سطر * ستون) + " "
        پایان
        سرچاپ خط
    پایان
پایان
\end{salamfa}

\begin{outputfa}
1 2 3 4 5
2 4 6 8 10
3 6 9 12 15
4 8 12 16 20
5 10 15 20 25
\end{outputfa}

حلقه‌ی بیرونی پنج دور دارد؛ در هر دور، حلقه‌ی درونی پنج بار می‌چرخد و یک
خط از جدول را می‌سازد. در کل ۲۵ ضرب انجام می‌شود. متغیر \fcode{خط} را
\emph{درون} حلقه‌ی بیرونی تعریف کردیم تا برای هر سطر، تازه و خالی باشد.

\program{مثلث ستاره‌ای}
\begin{salamfa}
روال ریشه:
    تکرار ۱ تا ۴ در ردیف:
        ناپایا خط := ""
        تکرار ردیف:
            خط += "*"
        پایان
        سرچاپ خط
    پایان
پایان
\end{salamfa}

\begin{outputfa}
*
**
***
****
\end{outputfa}

این‌جا تعداد دورهای حلقه‌ی درونی به شمارنده‌ی حلقه‌ی بیرونی بستگی دارد:
ردیف اول یک ستاره، ردیف دوم دو ستاره و همین‌طور تا آخر. با تغییرهای کوچک
در همین برنامه می‌توان شکل‌های گوناگونی ساخت؛ تمرین‌های پایان فصل چند
نمونه پیشنهاد می‌کنند.

\section{یک برنامه‌ی کامل‌تر}

جدول تبدیل دما از سلسیوس به فارنهایت، از صفر تا چهل درجه، ده تا ده تا:

\program{جدول تبدیل دما}
\begin{salamfa}
روال ریشه:
    سرچاپ "سلسیوس -> فارنهایت"
    سرچاپ "-------------------"
    تکرار ۰ تا ۴۰ هر ۱۰ در سلسیوس:
        فارنهایت := (((سلسیوس برگردان اعشار۶۴) * ۹.۰) / ۵.۰) + ۳۲.۰
        سرچاپ سلسیوس, "->", فارنهایت
    پایان
پایان
\end{salamfa}

\begin{outputfa}
سلسیوس -> فارنهایت
-------------------
0 -> 32
10 -> 50
20 -> 68
30 -> 86
40 -> 104
\end{outputfa}

شمارنده‌ی حلقه عددی صحیح است، پس پیش از محاسبه‌ی اعشاری آن را
\fcode{برگردان اعشار۶۴} تبدیل کردیم. متغیر \fcode{فارنهایت} درون حلقه و بدون
\fcode{ناپایا} تعریف شده است؛ در هر دور جعبه‌ای تازه ساخته می‌شود، مقدار
می‌گیرد، چاپ می‌شود و با پایان دور از بین می‌رود.

\begin{summary}
\begin{itemize}
  \item \fcode{تکرار ن:} بلوک را ن بار اجرا می‌کند. با \fcode{تکرار ن در شماره:}
        شماره‌ی دور (از صفر) را می‌گیرید.
  \item \fcode{تکرار آغاز تا انتها در شماره:} از آغاز تا انتها می‌شمارد،
        هر دو سر شامل؛ اگر آغاز بزرگ‌تر باشد، رو به پایین. \fcode{هر}
        اندازه‌ی پرش را مشخص می‌کند.
  \item \fcode{تا شرط:} تا وقتی شرط درست است تکرار می‌کند. چیزی درون
        حلقه باید شرط را سرانجام نادرست کند.
  \item \fcode{بشکن} حلقه را تمام می‌کند؛ \fcode{گذر} به دور بعد می‌پرد.
  \item حلقه‌ها می‌توانند تودرتو باشند؛ حلقه‌ی درونی برای هر دور بیرونی
        کامل اجرا می‌شود.
  \item متغیری که در حلقه جمع می‌شود، باید \fcode{ناپایا} و بیرون حلقه تعریف
        شود.
\end{itemize}
\end{summary}

\section*{تمرین‌ها}

\exercise عددهای ۱ تا ۵۰ را که بر ۷ بخش‌پذیرند، در یک خط چاپ کنید.

\exercise حاصل ضرب عددهای ۱ تا ۱۰ را با حلقه حساب و چاپ کنید.

\exercise مبلغی را در بانک با سود سالانه‌ی ۲۰ درصد می‌گذارید. با
\fcode{تا} حساب کنید بعد از چند سال مبلغ دو برابر می‌شود و در هر سال
موجودی را چاپ کنید.

\exercise برنامه‌ی مثلث ستاره‌ای را طوری تغییر دهید که مثلث وارونه بسازد
(اول چهار ستاره، بعد سه، \ldots). سپس نسخه‌ای بنویسید که یک مربع توخالی
۵ در ۵ از ستاره بسازد.

\exercise با دو حلقه‌ی تودرتو همه‌ی جفت‌عددهای میان ۱ تا ۱۰ را پیدا و چاپ
کنید که مجموعشان ۱۰ است (مثل ۱ و ۹، ۲ و ۸). هر جفت فقط یک بار چاپ شود.

\exercise برنامه‌ی زیر چند خط چاپ می‌کند؟ اول حساب کنید، بعد اجرا کنید:
\begin{salamfa}
روال ریشه:
    تکرار ۳ در الف:
        تکرار ۴ در ب:
            اگر ب == ۲: بشکن پایان
            سرچاپ الف, ب
        پایان
    پایان
پایان
\end{salamfa}

\chapter{روال‌ها}\label{ch:functions}

\begin{goals}
\begin{itemize}
  \item روال چیست و چرا برنامه را به روال‌ها تقسیم می‌کنیم
  \item تعریف روال و فراخوانی آن
  \item فرستادن مقدار به روال با پارامتر
  \item گرفتن نتیجه از روال با \fcode{برگشت}
  \item مقدار پیش‌فرض پارامترها و محدوده‌ی متغیرها
\end{itemize}
\end{goals}

\section{یک نام برای یک کار}

از فصل اول با یک روال کار می‌کنیم: \fcode{روال ریشه}. حالا وقت آن است که
ببینیم «روال» واقعاً چیست. \textbf{روال} بلوکی از دستورهاست که یک نام
دارد. هر جا آن نام را بنویسیم (به این کار \textbf{فراخوانی} می‌گویند)،
دستورهای درون روال اجرا می‌شوند. روال ریشه هم همین است، با این تفاوت که
سلام آن را خودش، در آغاز برنامه، فراخوانی می‌کند.

فرض کنید در برنامه‌ای چند بار باید یک خط جداکننده چاپ کنیم. می‌توانیم هر
بار \mbox{\scode{سرچاپ "--------------------"}} بنویسیم، یا این کار را یک بار در
روالی تعریف کنیم و هر جا لازم بود، نامش را صدا بزنیم:

\program{تعریف و فراخوانی روال}
\begin{salamfa}
روال خط جداکننده():
    سرچاپ "--------------------"
پایان

روال ریشه:
    خط جداکننده()
    سرچاپ "فهرست خرید"
    خط جداکننده()
    سرچاپ "نان، پنیر، شیر"
    خط جداکننده()
پایان
\end{salamfa}

\begin{outputfa}
--------------------
فهرست خرید
--------------------
نان، پنیر، شیر
--------------------
\end{outputfa}

تعریف روال با واژه‌ی \kwd{روال} شروع می‌شود، بعد نام روال می‌آید، بعد یک
جفت پرانتز و دونقطه، و در پایان \fcode{پایان}. نام روال هم مثل نام متغیر
می‌تواند چند واژه‌ای باشد. فراخوانی، نوشتن نام روال همراه با پرانتز است:
\scode{خط جداکننده()}. پرانتزها لازم‌اند؛ اگر آن‌ها را جا بیندازید، سلام خطا
می‌دهد: «روال \scode{'خط جداکننده'} فراخوانی نشده است: برای فراخوانی آن بنویسید
\scode{'خط جداکننده()'}».

به دو نکته دقت کنید. اول این که روال را \emph{بیرون} روال ریشه تعریف
کردیم، در بالای برنامه؛ روال‌ها همیشه در سطح بالای برنامه و کنار هم
تعریف می‌شوند، نه درون یکدیگر. دوم این که ترتیب تعریف مهم نیست؛ می‌توانستیم
\fcode{خط جداکننده} را زیر \fcode{روال ریشه} هم بنویسیم. سلام اول همه‌ی
تعریف‌ها را می‌بیند و بعد از \fcode{روال ریشه} شروع می‌کند.

\begin{note}[چرا روال؟]
سه دلیل اصلی دارد. \textbf{تکرار نکردن}: کاری که در چند جا لازم است، یک بار
نوشته می‌شود و اگر روزی باید عوض شود، یک جا عوض می‌شود. \textbf{نام
گذاشتن}: خواندن \scode{خط جداکننده()} بسیار آسان‌تر از خواندن یک ردیف خط
تیره در میان کد است؛ نام، مقصود را می‌گوید.
\textbf{تقسیم کار}: یک مسئله‌ی بزرگ را به چند روال کوچک می‌شکنیم و هر
کدام را جداگانه می‌نویسیم و آزمایش می‌کنیم.
\end{note}

\section{پارامتر: فرستادن مقدار به روال}

روال \fcode{خط جداکننده} همیشه یک کار ثابت می‌کند. اما بیشتر روال‌ها به
اطلاعاتی نیاز دارند تا کارشان را انجام دهند: نام کسی که باید به او سلام
کرد، عددی که باید مربعش را حساب کرد. این اطلاعات را از راه
\textbf{پارامتر} به روال می‌فرستیم. پارامترها را درون پرانتز تعریف روال
می‌نویسیم و برای هر یک باید \emph{گونه‌اش} را هم بگوییم:

\program{روال با پارامتر}
\begin{salamfa}
روال سلام کن(نام: رشته):
    سرچاپ "سلام", نام + "! خوش آمدی."
پایان

روال ریشه:
    سلام کن("سارا")
    سلام کن("رضا")
    دوست := "مینا"
    سلام کن(دوست)
پایان
\end{salamfa}

\begin{outputfa}
سلام سارا! خوش آمدی.
سلام رضا! خوش آمدی.
سلام مینا! خوش آمدی.
\end{outputfa}

\scode{نام: رشته} می‌گوید این روال یک پارامتر به نام \fcode{نام} از گونه‌ی
رشته می‌خواهد. در هر فراخوانی، مقداری که درون پرانتز می‌نویسیم (به آن
\textbf{آرگومان} می‌گویند) در جعبه‌ی \fcode{نام} گذاشته می‌شود و روال با
آن کار می‌کند. آرگومان می‌تواند یک مقدار مستقیم باشد، یک متغیر، یا حتی
یک محاسبه.

روال می‌تواند چند پارامتر داشته باشد که با ویرگول جدا می‌شوند. هنگام
فراخوانی، آرگومان‌ها به همان ترتیب به پارامترها می‌رسند:

\program{چند پارامتر}
\begin{salamfa}
روال ردیف علامت(تعداد: صحیح, علامت: رشته):
    ناپایا خط := ""
    تکرار تعداد:
        خط += علامت
    پایان
    سرچاپ خط
پایان

روال ریشه:
    ردیف علامت(۵, "*")
    ردیف علامت(۳, "#")
    ردیف علامت(۸, "=")
پایان
\end{salamfa}

\begin{outputfa}
*****
###
========
\end{outputfa}

\begin{warn}
گونه‌ی آرگومان باید با گونه‌ی پارامتر بخواند. فراخوانی \scode{ردیف علامت("5", "*")}
خطاست، چون \code{"5"} رشته است و پارامتر اول عدد صحیح می‌خواهد. تعداد
آرگومان‌ها هم باید درست باشد؛ نه کمتر و نه بیشتر.
\end{warn}

\section{بازگشت: گرفتن نتیجه از روال}

روال‌های بالا کاری می‌کردند (چاپ)، اما چیزی به فراخواننده پس نمی‌دادند.
بسیاری از روال‌ها باید نتیجه‌ای \emph{بدهند}: مساحت، قیمت نهایی، جواب یک
سؤال. برای این کار دو چیز لازم است: گونه‌ی نتیجه را بعد از پرانتز اعلام
کنیم، و درون روال با \kwd{برگشت} نتیجه را پس بدهیم:

\program{روال با مقدار بازگشتی}
\begin{salamfa}
روال مساحت مستطیل(طول: صحیح, عرض: صحیح): صحیح:
    برگشت طول * عرض
پایان

روال ریشه:
    سرچاپ "مساحت اتاق:", مساحت مستطیل(۴, ۵)
    سرچاپ "مساحت حیاط:", مساحت مستطیل(۱۲, ۸)
    مجموع := مساحت مستطیل(۴, ۵) + مساحت مستطیل(۱۲, ۸)
    سرچاپ "مجموع:", مجموع
پایان
\end{salamfa}

\begin{outputfa}
مساحت اتاق: 20
مساحت حیاط: 96
مجموع: 116
\end{outputfa}

به سرآغاز روال نگاه کنید:
\scode{روال مساحت مستطیل(طول: صحیح, عرض: صحیح): صحیح:}. بعد از پرانتز، یک دونقطه و گونه‌ی نتیجه (\fcode{صحیح}) آمده و بعد
دونقطه‌ی همیشگی که بلوک را باز می‌کند. دستور \scode{برگشت طول * عرض}
مقدار را حساب می‌کند، آن را به فراخواننده می‌دهد و روال همان‌جا تمام
می‌شود.

مهم‌ترین نکته این است که فراخوانی چنین روالی، خودش یک \emph{مقدار} است و
می‌توان آن را هر جا که یک مقدار می‌رود گذاشت: در \fcode{سرچاپ}، در یک
محاسبه، در تعریف یک متغیر، یا حتی به عنوان آرگومان روالی دیگر.

\subsection{روالی که بله یا نه می‌گوید}

روالی که نتیجه‌ی منطقی می‌دهد، در شرط‌ها بسیار خوش‌دست است:

\program{روال منطقی}
\begin{salamfa}
روال زوج است(عدد: صحیح): منطقی:
    برگشت (عدد % ۲) == ۰
پایان

روال ریشه:
    تکرار ۱ تا ۶ در عدد:
        اگر زوج است(عدد):
            سرچاپ عدد, "زوج"
        وگرنه:
            سرچاپ عدد, "فرد"
        پایان
    پایان
پایان
\end{salamfa}

\begin{outputfa}
1 فرد
2 زوج
3 فرد
4 زوج
5 فرد
6 زوج
\end{outputfa}

خط \scode{اگر زوج است(عدد):} تقریباً مثل فارسی خوانده می‌شود. این یکی از
زیبایی‌های نام‌گذاری خوب روال‌هاست. رسم خوبی است که نام روال‌های منطقی را
به شکل پرسش بگذاریم: \fcode{زوج است}، \fcode{معتبر است}، \fcode{خالی است}.

\subsection{چند بازگشت در یک روال}

روال می‌تواند در چند جا \fcode{برگشت} داشته باشد. به محض اجرای اولین
\fcode{برگشت}، روال تمام می‌شود و بقیه‌ی خط‌ها اجرا نمی‌شوند. این ویژگی
زنجیره‌های شرطی را کوتاه می‌کند:

\program{چند بازگشت}
\begin{salamfa}
روال برچسب نمره(نمره: صحیح): رشته:
    اگر نمره >= ۱۸: برگشت "عالی" پایان
    اگر نمره >= ۱۵: برگشت "خوب" پایان
    اگر نمره >= ۱۰: برگشت "قابل قبول" پایان
    برگشت "مردود"
پایان

روال ریشه:
    سرچاپ برچسب نمره(۱۹)
    سرچاپ برچسب نمره(۱۶)
    سرچاپ برچسب نمره(۱۲)
    سرچاپ برچسب نمره(۷)
پایان
\end{salamfa}

\begin{outputfa}
عالی
خوب
قابل قبول
مردود
\end{outputfa}

با نمره‌ی ۱۶، شرط اول نادرست است، شرط دوم درست است و روال با «خوب»
برمی‌گردد؛ خط‌های بعدی هرگز دیده نمی‌شوند. مقایسه کنید با همان کار در فصل~\ref{ch:conditions} که با \fcode{وگرنه} نوشته بودیم.

\begin{warn}
سلام با روال‌ها هم مثل متغیرها سخت‌گیر است: روالی که تعریف شده اما
هیچ‌جا فراخوانی نشده، خطاست. همچنین روالی که گونه‌ی بازگشتی اعلام کرده،
باید در \emph{همه‌ی} مسیرها چیزی برگرداند. اگر \fcode{برگشت "مردود"} آخر را
حذف کنید، سلام \emph{هشدار} می‌دهد که روال ممکن است بدون بازگرداندن مقدار به
پایان برسد. برنامه با وجود هشدار اجرا می‌شود، اما برای نمره‌ی ۷ روال دیگر
برچسبی برنمی‌گرداند و به جای آن \code{null} چاپ می‌شود؛ پس هشدارها را هم مثل خطاها جدی بگیرید. روالی مثل
\fcode{سلام کن} که فقط کاری انجام می‌دهد و نتیجه‌ای پس نمی‌دهد، به
\fcode{برگشت} نیازی ندارد.
\end{warn}

\subsection{گونه‌ی تهی}

روالی که چیزی پس نمی‌دهد هم گونه‌ی بازگشتی دارد: \kwd{تهی}، یعنی «هیچ».
وقتی گونه‌ی بازگشتی را ننویسید، همان‌طور که برای \fcode{سلام کن} ننوشتیم،
سلام آن را \fcode{تهی} می‌گیرد. می‌توانید آن را آشکارا هم بنویسید؛ برخی
برنامه‌نویسان این کار را می‌پسندند، چون روشن می‌گوید که روال مقداری
برنمی‌گرداند:

\program{روالی با گونه‌ی تهی}
\begin{salamfa}
روال خط(پهنا: صحیح): تهی:
    ناپایا متن := ""
    تکرار پهنا:
        متن += "-"
    پایان
    سرچاپ متن
پایان

روال ریشه:
    خط(۱۰)
    سرچاپ "سلام"
    خط(۱۰)
پایان
\end{salamfa}

\begin{outputfa}
----------
سلام
----------
\end{outputfa}

چون \fcode{خط} چیزی برنمی‌گرداند، فراخوانی آن مقدار نیست: نمی‌توان آن را در
\fcode{سرچاپ} یا در یک محاسبه گذاشت، فقط در خطی جداگانه.

\subsection{روال با عدد اعشاری}

\program{تبدیل دما به شکل روال}
\begin{salamfa}
روال به فارنهایت(سلسیوس: اعشار۶۴): اعشار۶۴:
    برگشت ((سلسیوس * ۹.۰) / ۵.۰) + ۳۲.۰
پایان

روال ریشه:
    سرچاپ "آب یخ می‌زند:", به فارنهایت(۰.۰)
    سرچاپ "دمای بدن:", به فارنهایت(۳۷.۰)
    سرچاپ "آب می‌جوشد:", به فارنهایت(۱۰۰.۰)
پایان
\end{salamfa}

\begin{outputfa}
آب یخ می‌زند: 32
دمای بدن: 98.6
آب می‌جوشد: 212
\end{outputfa}

فرمولی که در فصل~\ref{ch:loops} درون حلقه نوشته بودیم، حالا یک نام دارد
و هر جا لازم باشد قابل استفاده است. دقت کنید که آرگومان‌ها را با نقطه‌ی
اعشار نوشتیم (\scode{۰.۰})، چون پارامتر اعشاری است. سلام \code{۰} را هم
می‌پذیرد و خودش آن را اعشاری می‌کند، اما \scode{۰.۰} مقصود را روشن‌تر نشان
می‌دهد.

\section{مقدار پیش‌فرض برای پارامتر}

گاهی یک پارامتر بیشتر وقت‌ها مقدار مشخصی دارد و فقط گاهی عوض می‌شود.
می‌توان در تعریف روال برایش \textbf{مقدار پیش‌فرض} گذاشت؛ آن وقت
فراخواننده می‌تواند آن آرگومان را ننویسد:

\program{پارامتر با مقدار پیش‌فرض}
\begin{salamfa}
روال سلام کن(نام: رشته, پیشوند: رشته = "سلام"):
    سرچاپ پیشوند, نام
پایان

روال ریشه:
    سلام کن("سارا")
    سلام کن("دکتر اسماعیلی", "درود بر")
پایان
\end{salamfa}

\begin{outputfa}
سلام سارا
درود بر دکتر اسماعیلی
\end{outputfa}

پارامترهای پیش‌فرض‌دار باید بعد از پارامترهای معمولی بیایند.

\section{محدوده‌ی متغیرها}

متغیری که درون یک روال تعریف می‌شود، فقط درون همان روال وجود دارد. به
این ویژگی \textbf{محدوده} می‌گویند. روال ریشه نمی‌تواند متغیرهای درون
\fcode{سلام کن} را ببیند و برعکس. پارامترها هم همین‌طورند: جعبه‌هایی که
با شروع روال ساخته می‌شوند و با پایان آن از بین می‌روند.

نکته‌ی مهم‌تر این است که عددها، رشته‌ها و مقدارهای منطقی به شکل \emph{کپی}
به روال می‌رسند. روال هر کاری با نسخه‌ی خودش بکند، متغیر فراخواننده
دست‌نخورده می‌ماند (آرایه‌ها استثنا هستند؛ فصل~\ref{ch:arrays} را ببینید):

\program{روال کپی می‌گیرد}
\begin{salamfa}
روال دوبرابر کن(عدد: صحیح): صحیح:
    ناپایا نتیجه := عدد
    نتیجه *= ۲
    برگشت نتیجه
پایان

روال ریشه:
    مقدار := ۲۱
    سرچاپ "نتیجه‌ی روال:", دوبرابر کن(مقدار)
    سرچاپ "مقدار اصلی:", مقدار
پایان
\end{salamfa}

\begin{outputfa}
نتیجه‌ی روال: 42
مقدار اصلی: 21
\end{outputfa}

این همان چیزی است که در فصل~\ref{ch:variables} درباره‌ی کپی شدن مقدارها
دیدیم. اگر روال بخواهد نتیجه‌ای به بیرون بدهد، راه درستش \fcode{برگشت}
است، نه دست‌کاری آرگومان.

\subsection{پایدارهای سراسری}

گاهی مقداری در چند روال لازم است. \fcode{پایا}‌هایی که در بالای برنامه و
بیرون همه‌ی روال‌ها تعریف شوند، در همه جا در دسترس‌اند:

\program{پایدار سراسری}
\begin{salamfa}
پایا مالیات := ۹

روال قیمت نهایی(قیمت: صحیح): صحیح:
    برگشت قیمت + ((قیمت * مالیات) / ۱۰۰)
پایان

روال ریشه:
    سرچاپ قیمت نهایی(۱۰۰۰۰۰)
    سرچاپ قیمت نهایی(۲۵۰۰۰۰)
    سرچاپ "نرخ:", مالیات, "درصد"
پایان
\end{salamfa}

\begin{outputfa}
109000
272500
نرخ: 9 درصد
\end{outputfa}

بیرون از روال‌ها فقط \fcode{پایا} تعریف کنید، نه متغیر تغییرپذیر. متغیر
سراسری‌ای که هر روالی بتواند عوضش کند، پیدا کردن اشکال‌ها را بسیار سخت
می‌کند؛ بهتر است هر چه یک روال لازم دارد را از راه پارامتر بگیرد و هر چه
تولید می‌کند را از راه بازگشت پس بدهد.

\begin{deep}[روالی که خودش را صدا می‌زند]
روال می‌تواند درون خودش، خودش را فراخوانی کند. به این کار
\textbf{بازگشتی} می‌گویند و کارهای جالبی با آن می‌شود کرد. برنامه‌ی زیر
با همین ترفند شمارش معکوس می‌کند: هر فراخوانی، عددی یکی کمتر را
می‌فرستد، تا وقتی که به صفر برسد و زنجیره قطع شود. \fcode{برگشت} بدون
مقدار، روال بی‌نتیجه را همان‌جا تمام می‌کند. در این کتاب به این موضوع
بیشتر نمی‌پردازیم، اما بدانید که وجود دارد.
\end{deep}

\program{شمارش معکوس بازگشتی}
\begin{salamfa}
روال پس شماری(عدد: صحیح):
    اگر عدد == ۰:
        سرچاپ "پایان"
        برگشت
    پایان
    سرچاپ عدد
    پس شماری(عدد - ۱)
پایان

روال ریشه:
    پس شماری(۳)
پایان
\end{salamfa}
\begin{outputfa}
3
2
1
پایان
\end{outputfa}

\section{یک برنامه‌ی کامل‌تر}

صورت‌حساب یک رستوران، که هر بخش کارش را به یک روال سپرده است:

\program{صورت‌حساب رستوران}
\begin{salamfa}
پایا خدمات := ۱۰

روال مبلغ سفارش(قیمت واحد: صحیح, تعداد: صحیح): صحیح:
    برگشت قیمت واحد * تعداد
پایان

روال اعمال تخفیف(مبلغ: صحیح, درصد: صحیح): صحیح:
    برگشت مبلغ - ((مبلغ * درصد) / ۱۰۰)
پایان

روال نمایش سطر(عنوان: رشته, مبلغ: صحیح):
    سرچاپ عنوان + ":", مبلغ, "تومان"
پایان

روال ریشه:
    چلوکباب := مبلغ سفارش(۱۸۰۰۰۰, ۲)
    نوشابه := مبلغ سفارش(۱۵۰۰۰, ۳)
    جمع := چلوکباب + نوشابه
    حق خدمات := (جمع * خدمات) / ۱۰۰

    نمایش سطر("چلوکباب", چلوکباب)
    نمایش سطر("نوشابه", نوشابه)
    نمایش سطر("جمع", جمع)
    نمایش سطر("حق خدمات", حق خدمات)
    نمایش سطر("قابل پرداخت", اعمال تخفیف(جمع + حق خدمات, ۵))
پایان
\end{salamfa}

\begin{outputfa}
چلوکباب: 360000 تومان
نوشابه: 45000 تومان
جمع: 405000 تومان
حق خدمات: 40500 تومان
قابل پرداخت: 423225 تومان
\end{outputfa}

روال ریشه‌ی این برنامه را بخوانید: تقریباً مثل یک متن فارسی است، چون همه‌ی
جزئیات محاسبه در روالی با نام گویا پنهان شده است. به آخرین خط دقت کنید:
نتیجه‌ی \fcode{اعمال تخفیف} مستقیم به عنوان آرگومان به \fcode{نمایش سطر}
داده شده است. روال‌ها این‌گونه در هم می‌نشینند.

\begin{summary}
\begin{itemize}
  \item روال بلوکی نام‌دار است: \scode{روال نام(پارامترها): گونه بازگشتی:}
        \ldots\ \fcode{پایان}. فراخوانی با نام و پرانتز انجام می‌شود.
  \item روال‌ها در سطح بالای برنامه تعریف می‌شوند؛ ترتیبشان مهم نیست.
  \item پارامترها با گونه‌شان نوشته می‌شوند: \scode{نام: رشته}. آرگومان‌ها
        به ترتیب به آن‌ها می‌رسند؛ عددها و رشته‌ها به شکل کپی.
  \item \fcode{برگشت مقدار} نتیجه را پس می‌دهد و روال را تمام می‌کند.
        فراخوانی چنین روالی خودش یک مقدار است.
  \item پارامتر می‌تواند مقدار پیش‌فرض داشته باشد.
  \item متغیرهای درون روال فقط در همان روال وجود دارند. برای مقدارهای
        مشترک، \fcode{پایا} سراسری تعریف کنید.
\end{itemize}
\end{summary}

\section*{تمرین‌ها}
\markright{تمرین‌ها}

\exercise روالی به نام \fcode{بیشینه} بنویسید که دو عدد صحیح بگیرد و
بزرگ‌ترشان را برگرداند. سپس با کمک آن، بزرگ‌ترین سه عدد را پیدا کنید.

\exercise روالی بنویسید که شعاع دایره (اعشاری) را بگیرد و مساحتش را
برگرداند (عدد پی را ۳٫۱۴۱۶ بگیرید). با یک حلقه، مساحت دایره‌هایی با شعاع
۱ تا ۵ را چاپ کنید.

\exercise روالی به نام \fcode{سال کبیسه است} بنویسید که یک سال شمسی
بگیرد و بگوید آیا کبیسه است یا نه. قاعده‌ی ساده‌شده: سالی کبیسه است که
باقی‌مانده‌ی تقسیمش بر ۳۳ یکی از ۱، ۵، ۹، ۱۳، ۱۷، ۲۲، ۲۶ یا ۳۰ باشد. سپس
برای سال‌های ۱۴۰۳ تا ۱۴۰۸ نتیجه را چاپ کنید.

\exercise روالی بنویسید که تعداد ثانیه بگیرد و آن را به شکل رشته‌ای مثل
\code{"1:02:05"} (ساعت:دقیقه:ثانیه) برگرداند. راهنمایی: از محاسبه‌ی فصل~\ref{ch:operators} و چسباندن رشته‌ها استفاده کنید.

\exercise برنامه‌ی جدول ضرب فصل~\ref{ch:loops} را با روالی به نام
\fcode{سطر جدول} بازنویسی کنید که شماره‌ی سطر و تعداد ستون‌ها را بگیرد و
یک سطر جدول را چاپ کند. تعداد ستون‌ها مقدار پیش‌فرض ۱۰ داشته باشد.

\exercise روال زیر یک اشکال دارد. بدون اجرا پیدایش کنید:
\begin{salamfa}
روال قدرمطلق(عدد: صحیح): صحیح:
    اگر عدد < ۰:
        برگشت -عدد
    وگرنه عدد > ۰:
        برگشت عدد
    پایان
پایان
\end{salamfa}

\chapter{آرایه‌ها}
\label{ch:arrays}

\begin{goals}
\begin{itemize}
  \item نگه‌داری چند مقدار هم‌گونه در یک متغیر
  \item دسترسی به هر عضو با شماره (اندیس) و تغییر دادن آن
  \item پیمایش آرایه با \fcode{هر}
  \item الگوهای پایه: جمع، میانگین، بیشترین، شمارش و جست‌وجو
  \item کار با چند آرایه‌ی هم‌راستا
\end{itemize}
\end{goals}

\section{سی متغیر یا یک آرایه؟}

فرض کنید می‌خواهید نمره‌های سی دانش‌آموز را نگه دارید و میانگین بگیرید.
با آنچه تا حالا داریم، باید سی متغیر بسازید: \fcode{نمره۱}، \fcode{نمره۲}،
\ldots{} و سی بار آن‌ها را جمع کنید. و اگر کلاس سی‌ویک نفر شد؟ روشن است که این
راه به جایی نمی‌رسد.

\textbf{آرایه} پاسخ این مشکل است: یک متغیر که به جای یک مقدار، فهرستی از
مقدارهای هم‌گونه را نگه می‌دارد. به جای سی جعبه‌ی جداگانه، یک قفسه با سی
خانه‌ی شماره‌دار.

\section{ساختن آرایه و دسترسی به عضوها}

آرایه را با نوشتن مقدارها میان دو کروشه و با ویرگول می‌سازیم:

\program{ساختن آرایه}
\begin{salamfa}
روال ریشه:
    نمرات := [۱۷, ۱۵, ۱۹, ۱۲, ۲۰]
    سرچاپ نمرات
    سرچاپ "اولین نمره:", نمرات[۰]
    سرچاپ "سومین نمره:", نمرات[۲]
    سرچاپ "تعداد نمره‌ها:", نمرات.طول()
پایان
\end{salamfa}

\begin{outputfa}
[17, 15, 19, 12, 20]
اولین نمره: 17
سومین نمره: 19
تعداد نمره‌ها: 5
\end{outputfa}

چند چیز در این برنامه تازه است:
\begin{itemize}
  \item \fcode{سرچاپ نمرات} کل آرایه را با همان شکل کروشه‌ای چاپ می‌کند.
        برای نگاه سریع به محتوای آرایه بسیار مفید است.
  \item \scode{نمرات[۰]} عضو اول آرایه است. عدد داخل کروشه \textbf{اندیس}
        نام دارد و \emph{از صفر} شروع می‌شود: عضو اول اندیس ۰ دارد، عضو
        دوم اندیس ۱ و همین‌طور تا آخر. پس \scode{نمرات[۲]} عضو سوم است.
  \item \scode{نمرات.طول()} تعداد عضوهای آرایه را می‌دهد. با همان نقطه‌ای
        نوشته می‌شود که در فصل \ref{ch:types} برای \fcode{به صحیح} دیدیم.
\end{itemize}

\begin{note}[چرا از صفر؟]
اندیس را «فاصله از آغاز» بدانید نه «شماره‌ی ردیف». عضو اول، صفر خانه از
آغاز فاصله دارد. با این نگاه، شروع از صفر طبیعی است و در فصل
\ref{ch:loops} هم دیدید که شمارنده‌ی \fcode{تکرار} از صفر شروع می‌شود؛
این دو با هم جور درمی‌آیند. آخرین عضو آرایه‌ای با ۵ عضو، اندیس ۴ دارد، یعنی
همیشه \scode{طول() - ۱}.
\end{note}

\begin{warn}
اگر اندیسی بیرون از محدوده بدهید، خطا می‌گیرید. وقتی اندیس یک عدد ثابت است
(مثلاً \scode{نمرات[۵]} برای آرایه‌ای با ۵ عضو، یا \scode{نمرات[-۱]})، سلام
پیش از اجرا خبر می‌دهد: «اندیس ۵ برای آرایه‌ای به طول ۵ خارج از محدوده است».
وقتی اندیس در یک متغیر است و هنگام اجرا از محدوده بیرون می‌زند، برنامه همان‌جا
با خطای زمان اجرا متوقف می‌شود. این یکی از رایج‌ترین خطاهای کار با آرایه
است و بیشتر وقت‌ها از فراموش کردن «از صفر» می‌آید.
\end{warn}

\section{تغییر دادن عضوها}

برای عوض کردن یک عضو، با اندیس به همان عضو مقدار می‌دهیم. برای این کار
لازم نیست آرایه \fcode{ناپایا} باشد؛ \fcode{ناپایا} فقط وقتی لازم است که
بخواهید کل متغیر را با آرایه‌ی دیگری جایگزین کنید:

\program{تغییر عضو آرایه}
\begin{salamfa}
روال ریشه:
    نمرات := [۱۷, ۱۵, ۱۹, ۱۲, ۲۰]
    سرچاپ "پیش از اصلاح:", نمرات
    نمرات[۳] = ۱۴
    نمرات[۰] = نمرات[۰] + ۱
    سرچاپ "پس از اصلاح:", نمرات
پایان
\end{salamfa}

\begin{outputfa}
پیش از اصلاح: [17, 15, 19, 12, 20]
پس از اصلاح: [18, 15, 19, 14, 20]
\end{outputfa}

هر عضو آرایه مثل یک متغیر معمولی رفتار می‌کند: می‌توان خواندش، در محاسبه
گذاشتش و به آن مقدار داد. تعداد عضوها اما ثابت است؛ آرایه‌ای که با پنج عضو
ساخته شده، تا آخر پنج عضو دارد.

\section{پیمایش آرایه با هر}

قدرت واقعی آرایه وقتی پیداست که با حلقه ترکیب شود. سلام حلقه‌ای مخصوص
آرایه دارد: \kwd{هر}. این حلقه، به تعداد عضوهای آرایه دور می‌زند و در هر
دور یک عضو را در متغیری می‌گذارد:

\program{پیمایش با هر}
\begin{salamfa}
روال ریشه:
    میوه‌ها := ["سیب", "موز", "پرتقال"]
    هر میوه در میوه‌ها:
        سرچاپ "-", میوه
    پایان
پایان
\end{salamfa}

\begin{outputfa}
- سیب
- موز
- پرتقال
\end{outputfa}

\scode{هر میوه در میوه‌ها:} را این‌طور بخوانید: «برای هر میوه در میوه‌ها».
در دور اول \fcode{میوه} برابر «سیب» است، در دور دوم «موز» و در دور سوم
«پرتقال». نه اندیسی لازم است، نه شمردن؛ سلام خودش همه را یکی‌یکی می‌دهد.
آرایه‌ی رشته‌ها هم درست مثل آرایه‌ی عددهاست؛ فقط همه‌ی عضوها باید هم‌گونه
باشند.

اگر شماره‌ی هر عضو را هم بخواهید، دو نام میان پرانتز بنویسید؛ اولی اندیس
را می‌گیرد و دومی عضو را:

\program{پیمایش با اندیس}
\begin{salamfa}
روال ریشه:
    میوه‌ها := ["سیب", "موز", "پرتقال"]
    هر (شماره, میوه) در میوه‌ها:
        سرچاپ شماره + ۱, "-", میوه
    پایان
پایان
\end{salamfa}

\begin{outputfa}
1 - سیب
2 - موز
3 - پرتقال
\end{outputfa}

چون اندیس از صفر شروع می‌شود، برای نمایش شماره‌ی «انسانی» یکی به آن
اضافه کردیم.

\section{الگوهای پایه‌ی کار با آرایه}

بیشتر کارهایی که با آرایه می‌کنیم، به چند الگوی ساده برمی‌گردند. همه‌ی
آن‌ها از یک حلقه‌ی \fcode{هر} و یک یا دو متغیر تغییرپذیر ساخته می‌شوند.

\subsection{جمع و میانگین}

\program{جمع و میانگین نمره‌ها}
\begin{salamfa}
روال ریشه:
    نمرات := [۱۷, ۱۵, ۱۹, ۱۲, ۲۰]
    ناپایا مجموع := ۰
    هر نمره در نمرات:
        مجموع += نمره
    پایان
    تعداد := نمرات.طول()
    میانگین := (مجموع برگردان اعشار۶۴) / (تعداد برگردان اعشار۶۴)
    سرچاپ "مجموع:", مجموع
    سرچاپ "میانگین:", میانگین
پایان
\end{salamfa}

\begin{outputfa}
مجموع: 83
میانگین: 16.6
\end{outputfa}

همان الگوی جمع کردن فصل \ref{ch:loops} است، فقط این بار عددها از آرایه
می‌آیند. اگر فردا ده نمره‌ی دیگر به آرایه اضافه کنید، هیچ خط دیگری از
برنامه عوض نمی‌شود.

\subsection{بیشترین}

\program{گرم‌ترین روز هفته}
\begin{salamfa}
روال ریشه:
    دماها := [۲۱, ۲۵, ۱۹, ۳۰, ۲۷, ۲۴, ۲۲]
    ناپایا بیشترین := دماها[۰]
    ناپایا روز بیشترین := ۰
    هر (روز, دما) در دماها:
        اگر دما > بیشترین:
            بیشترین = دما
            روز بیشترین = روز
        پایان
    پایان
    سرچاپ "گرم‌ترین روز: روز", روز بیشترین + ۱, "با دمای", بیشترین
پایان
\end{salamfa}

\begin{outputfa}
گرم‌ترین روز: روز 4 با دمای 30
\end{outputfa}

همان الگوی «بزرگ‌ترین سه عدد» فصل \ref{ch:conditions} است: با عضو اول
شروع می‌کنیم و هر عضوی که بزرگ‌تر بود، جایش را می‌گیرد. این بار اندیس عضو
برنده را هم نگه داشتیم تا بدانیم کدام روز بوده است. برای پیدا کردن کمترین
کافی است \code{>} را به \code{<} تبدیل کنید.

\Needspace{16\baselineskip}
\subsection{شمارش}

\program{شمارش قبولی‌ها}
\begin{salamfa}
روال ریشه:
    نمرات := [۱۷, ۹, ۱۹, ۱۲, ۸, ۲۰, ۱۴]
    ناپایا قبولی := ۰
    هر نمره در نمرات:
        اگر نمره >= ۱۰: قبولی++ پایان
    پایان
    سرچاپ "قبول:", قبولی, "مردود:", نمرات.طول() - قبولی
پایان
\end{salamfa}

\begin{outputfa}
قبول: 5 مردود: 2
\end{outputfa}

\subsection{جست‌وجو}

آیا نام مشخصی در فهرست هست؟ برای این پرسش، یک متغیر منطقی می‌سازیم که
اول \fcode{نادرست} است و به محض پیدا شدن، \fcode{درست} می‌شود. وقتی پیدا
شد، با \fcode{بشکن} از ادامه‌ی جست‌وجو صرف نظر می‌کنیم:

\program{جست‌وجو در فهرست}
\begin{salamfa}
روال ریشه:
    مهمانان := ["سارا", "رضا", "مینا", "علی"]
    جستجو := "مینا"
    ناپایا پیدا شد := نادرست
    هر نام در مهمانان:
        اگر نام == جستجو:
            پیدا شد = درست
            بشکن
        پایان
    پایان
    اگر پیدا شد:
        سرچاپ جستجو, "در فهرست مهمانان هست"
    وگرنه:
        سرچاپ جستجو, "در فهرست مهمانان نیست"
    پایان
پایان
\end{salamfa}

\begin{outputfa}
مینا در فهرست مهمانان هست
\end{outputfa}

مقدار \fcode{جستجو} را به نامی که در فهرست نیست تغییر دهید و دوباره اجرا
کنید.

\section{آرایه به عنوان پارامتر روال}

آرایه را می‌توان به روال فرستاد. گونه‌ی پارامتر به شکل \scode{صحیح[:]} نوشته
می‌شود که یعنی «آرایه‌ای از عددهای صحیح با هر طولی»، و هنگام فراخوانی،
بعد از نام آرایه \code{[:]} می‌نویسیم که یعنی «همه‌ی عضوها»:

\program{روالی که آرایه می‌گیرد}
\begin{salamfa}
روال مجموع(اعداد: صحیح[:]): صحیح:
    ناپایا جمع := ۰
    هر عدد در اعداد:
        جمع += عدد
    پایان
    برگشت جمع
پایان

روال ریشه:
    کوچک := [۱, ۲, ۳]
    سرچاپ مجموع(کوچک[:])
    خریدها := [۴۵۰۰۰, ۱۲۰۰۰۰, ۳۲۰۰۰]
    سرچاپ "جمع خرید:", مجموع(خریدها[:])
پایان
\end{salamfa}

\begin{outputfa}
6
جمع خرید: 197000
\end{outputfa}

با این روال، الگوی جمع کردن را یک بار برای همیشه نوشته‌ایم و برای هر
آرایه‌ای با هر طولی به کار می‌بریم. (نوشتن \code{[:]} با اعدادی درون
کروشه، مثل \scode{خریدها[۱:۳]}، فقط بخشی از آرایه را می‌فرستد؛ در این
کتاب به آن نیازی نداریم.) تمرین‌های پایان فصل از شما می‌خواهند
همین کار را برای بیشترین و میانگین بکنید.

\begin{warn}[آرایه کپی نمی‌شود]
برخلاف عددها و رشته‌ها، آرایه‌ای که با \code{[:]} به روال فرستاده می‌شود
\emph{کپی نمی‌شود}: روال با همان آرایه‌ی اصلی کار می‌کند. اگر روال عضوی از
آن را تغییر دهد، آرایه‌ی فراخواننده هم تغییر می‌کند.
\end{warn}

\section{چند آرایه‌ی هم‌راستا}

گاهی برای هر چیز چند داده داریم: نام دانش‌آموز و نمره‌اش، نام کالا و
قیمتش. راه ساده این است که برای هر داده یک آرایه بسازیم و قرارداد
کنیم که عضو شماره‌ی \fcode{ک} در همه‌ی آرایه‌ها به یک چیز مربوط است:

\program{نام‌ها و نمره‌ها}
\begin{salamfa}
روال ریشه:
    نامها := ["سارا", "رضا", "مینا"]
    نمرات := [۱۸, ۱۲, ۱۵]
    هر (شماره, نام) در نامها:
        سرچاپ نام + ":", نمرات[شماره]
    پایان
پایان
\end{salamfa}

\begin{outputfa}
سارا: 18
رضا: 12
مینا: 15
\end{outputfa}

حلقه روی آرایه‌ی نام‌ها می‌چرخد و با اندیسی که می‌گیرد، نمره‌ی همان نفر را
از آرایه‌ی دوم برمی‌دارد. این روش ساده و کارآمد است، به شرطی که دو آرایه
همیشه هم‌طول و هم‌ترتیب بمانند.

\begin{deep}[وقتی تعداد از پیش معلوم نیست]
آرایه‌های این فصل طول ثابت دارند: نمی‌توان به آن‌ها عضو تازه افزود یا عضوی
را حذف کرد. برای فهرست‌هایی که در طول برنامه بزرگ و کوچک می‌شوند، سلام
ساختاری به نام \fcode{وکتور} دارد که با \fcode{بیفزا} عضو می‌گیرد و با
\fcode{دربیاور} عضو پس می‌دهد. همچنین برای نگه‌داری داده‌های مرتبط (مثلاً
نام و نمره و سن یک دانش‌آموز با هم) به جای چند آرایه‌ی هم‌راستا، می‌توان
\fcode{ساختار} تعریف کرد. این‌ها موضوع کتاب‌های بعدی سلام هستند؛ آنچه
در این فصل آموختید، پایه‌ی همه‌ی آن‌هاست.
\end{deep}

\section{یک برنامه‌ی کامل‌تر}

کارنامه‌ی یک کلاس کوچک: چاپ نمره و وضعیت هر نفر، میانگین کلاس و نام
بهترین دانش‌آموز.

\program{کارنامه‌ی کلاس}
\begin{salamfa}
روال ریشه:
    نامها := ["سارا", "رضا", "مینا", "علی"]
    نمرات := [۱۸, ۱۲, ۱۵, ۱۹]
    ناپایا مجموع := ۰
    ناپایا بهترین := ۰

    سرچاپ "کارنامه‌ی کلاس"
    سرچاپ "----------------"
    هر (شماره, نام) در نامها:
        نمره := نمرات[شماره]
        وضعیت := نمره >= ۱۰ ؟ "قبول" : "مردود"
        سرچاپ نام + ":", نمره, "(" + وضعیت + ")"
        مجموع += نمره
        اگر نمره > نمرات[بهترین]: بهترین = شماره پایان
    پایان
    سرچاپ "----------------"
    تعداد := نمرات.طول()
    میانگین := (مجموع برگردان اعشار۶۴) / (تعداد برگردان اعشار۶۴)
    سرچاپ "میانگین کلاس:", میانگین
    سرچاپ "بهترین دانش‌آموز:", نامها[بهترین]
پایان
\end{salamfa}

\begin{outputfa}
کارنامه‌ی کلاس
----------------
سارا: 18 (قبول)
رضا: 12 (قبول)
مینا: 15 (قبول)
علی: 19 (قبول)
----------------
میانگین کلاس: 16
بهترین دانش‌آموز: علی
\end{outputfa}

متغیر \fcode{بهترین} این بار \emph{اندیس} بهترین نمره را نگه می‌دارد نه
خود نمره را؛ به این ترتیب در پایان می‌توانیم با همان اندیس، نام او را از
آرایه‌ی نام‌ها بیرون بکشیم. مقایسه کنید با برنامه‌ی «گرم‌ترین روز» که هر
دو را نگه می‌داشت. هر دو راه درست است؛ انتخاب به این بستگی دارد که در
پایان به چه چیزی نیاز دارید.

\begin{summary}
\begin{itemize}
  \item آرایه فهرستی از مقدارهای هم‌گونه است: \scode{نمرات := [۱۷, ۱۵, ۱۹]}.
        \fcode{سرچاپ} کل آرایه را نشان می‌دهد و \scode{.طول()} تعداد عضوها
        را.
  \item اندیس از صفر شروع می‌شود: \scode{آ[۰]} اولین عضو و
        \scode{آ[آ.طول() - ۱]} آخرین عضو است. اندیس بیرون محدوده خطاست.
  \item عضوها را با اندیس تغییر می‌دهیم: \scode{آ[۲] = ۵}. \fcode{ناپایا}
        فقط برای جایگزین کردن کل آرایه لازم است.
  \item \scode{هر عضو در آرایه:} همه‌ی عضوها را یکی‌یکی می‌دهد و
        \scode{هر (اندیس, عضو) در آرایه:} اندیس را هم.
  \item الگوهای جمع، میانگین، بیشترین، شمارش و جست‌وجو همه از یک حلقه و
        یک متغیر تغییرپذیر ساخته می‌شوند.
  \item پارامتر آرایه‌ای با \scode{صحیح[:]} نوشته می‌شود و آرایه با
        \scode{نام[:]} به روال فرستاده می‌شود.
\end{itemize}
\end{summary}

\section*{تمرین‌ها}
\markright{تمرین‌ها}

\exercise آرایه‌ای از قیمت چند کالا بسازید و گران‌ترین و ارزان‌ترین را
همراه با اختلافشان چاپ کنید.

\exercise روالی به نام \fcode{میانگین} بنویسید که آرایه‌ای از عددهای صحیح
بگیرد و میانگین اعشاری آن‌ها را برگرداند. سپس روالی به نام
\fcode{بیشترین} با همان الگو بنویسید.

\exercise آرایه‌ای از دمای هفت روز هفته بسازید و چاپ کنید چند روز دما از
میانگین هفته بالاتر بوده است. (راهنمایی: به دو حلقه نیاز دارید؛ اول
میانگین، بعد شمارش.)

\exercise آرایه‌ای از نام‌ها و آرایه‌ای هم‌راستا از سن‌ها بسازید. نام همه‌ی
کسانی را که زیر ۱۸ سال دارند چاپ کنید و در پایان بگویید مسن‌ترین نفر
کیست.

\exercise آرایه‌ای از ۱۰ عدد بسازید و یک آرایه‌ی تازه با
۱۰ صفر. با حلقه، آرایه‌ی دوم را وارونه‌ی آرایه‌ی اول پر کنید (عضو آخر اولی
در عضو اول دومی، و همین‌طور). هر دو را چاپ کنید.

\exercise برنامه‌ی زیر چه چاپ می‌کند؟ اول حدس بزنید، بعد اجرا کنید:
\begin{salamfa}
روال ریشه:
    آ := [۱, ۲, ۳, ۴, ۵]
    هر (ک, عدد) در آ:
        اگر (ک % ۲) == ۰: آ[ک] = عدد * ۱۰ پایان
    پایان
    سرچاپ آ
پایان
\end{salamfa}

\chapter{کار با متن}
\label{ch:strings}

\begin{goals}
\begin{itemize}
  \item چسباندن رشته‌ها و ساختن پیام‌های ترکیبی
  \item نویسه‌های ویژه: خط جدید، جدول‌بندی و نقل قول درون رشته
  \item بریدن بخشی از رشته، شکافتن به واژه‌ها و پیراستن فاصله‌ها
  \item جست‌وجو در رشته و مقایسه‌ی رشته‌ها
  \item تبدیل رشته به عدد و برعکس
\end{itemize}
\end{goals}

\section{متن همه‌جا هست}

نام‌ها، پیام‌ها، نشانی‌ها، شماره‌تلفن‌ها، تاریخ‌ها: بخش بزرگی از داده‌های
هر برنامه متن است. در فصل‌های پیش با رشته‌ها کار کرده‌ایم، اما همیشه به
عنوان یک تکه‌ی یکپارچه. در این فصل یاد می‌گیریم درون رشته را ببینیم، از آن
تکه‌ای برداریم، آن را بشکافیم و از آن اطلاعات بیرون بکشیم. همه‌ی این کارها
با \textbf{قابلیت‌هایی} انجام می‌شوند که با یک نقطه به رشته می‌چسبند، مثل
\fcode{به صحیح} که در فصل \ref{ch:types} دیدید.

\section{چسباندن و طول}

\program{چسباندن رشته‌ها}
\begin{salamfa}
روال ریشه:
    نام := "سارا"
    نام خانوادگی := "احمدی"
    نام کامل := نام + " " + نام خانوادگی
    سرچاپ نام کامل
    شماره := "09121234567"
    سرچاپ "طول شماره:", شماره.طول()
پایان
\end{salamfa}

\begin{outputfa}
سارا احمدی
طول شماره: 11
\end{outputfa}

\scode{طول()} تعداد خانه‌هایی را می‌دهد که رشته در حافظه می‌گیرد. برای
شماره‌تلفن یا هر متنی که از رقم و حرف انگلیسی ساخته شده، هر نویسه یک خانه
است و این عدد همان چیزی است که انتظار دارید.

\begin{warn}[طول متن فارسی]
رایانه هر حرف فارسی را در \emph{دو} خانه نگه می‌دارد، در حالی که هر رقم،
حرف انگلیسی یا فاصله یک خانه می‌گیرد. \scode{طول()} تعداد خانه‌ها را
می‌شمارد؛ پس \scode{"سلام".طول()} برابر ۸ است، نه ۴. به همین دلیل در
این کتاب \scode{طول()} و \scode{زیررشته()} را بیشتر روی متن‌هایی به کار
می‌بریم که از رقم و حرف انگلیسی ساخته شده‌اند (شماره، تاریخ، کد). سلام
برای شمردن دقیق حرف‌های فارسی ابزارهای دیگری دارد که به این کتاب مربوط
نیست.
\end{warn}

\section{نویسه‌های ویژه}

بعضی چیزها را نمی‌توان مستقیم درون نقل قول نوشت: رفتن به خط بعد، یا خود
علامت نقل قول. برای این‌ها سلام ترکیب‌هایی با خط مورب وارونه \code{\textbackslash}
دارد:

\begin{center}
\begin{tabular}{cl}
\toprule
نوشتار & معنی \\
\midrule
\code{\textbackslash n} & رفتن به خط بعد \\
\code{\textbackslash t} & جدول‌بندی (یک پرش بلند) \\
\code{\textbackslash "} & خود علامت نقل قول \\
\code{\textbackslash\textbackslash} & خود خط مورب وارونه \\
\bottomrule
\end{tabular}
\end{center}

\program{نویسه‌های ویژه}
\begin{salamfa}
روال ریشه:
    سرچاپ "خط اول\nخط دوم"
    سرچاپ "نام: \t سارا"
    سرچاپ "او گفت: \"سلام\""
پایان
\end{salamfa}

\begin{outputfa}
خط اول
خط دوم
نام:      سارا
او گفت: "سلام"
\end{outputfa}

با \code{\textbackslash n} یک \fcode{سرچاپ} می‌تواند چند خط بنویسد. این
ترکیب‌ها هر کدام \emph{یک} نویسه به حساب می‌آیند، هرچند با دو علامت نوشته
می‌شوند.

\section{چند جور دیگر نقل قول}

در این کتاب همیشه از \code{"..."} استفاده می‌کنیم؛ آن را یاد بگیرید و کافی
است. اما سلام چهار نشانه‌ی دیگر هم برای نوشتن متن ثابت در کد دارد که ممکن
است در برنامه‌ی دیگران ببینید، پس بد نیست بشناسیدشان:

\begin{center}
\begin{tabular}{cll}
\toprule
نشانه & معنی & دنباله‌های گریز \\
\midrule
\code{"..."} & رشته، با هر طولی & دارد \\
\code{«...»} & دقیقاً مثل \code{"..."} (فصل \ref{ch:first} را ببینید) & دارد \\
\code{'...'} & فقط \emph{یک} نویسه، نه یک رشته & دارد \\
\code{\textasciigrave...\textasciigrave} & رشته‌ی خام؛ هیچ دنباله‌ای پردازش نمی‌شود & ندارد \\
\code{u'...'} & دقیقاً یک نویسه‌ی یونیکد (فارسی یا هر زبان دیگر) & ندارد \\
\bottomrule
\end{tabular}
\end{center}

\program{سه جور دیگر}
\begin{salamfa}
روال ریشه:
    حرف := 'A'
    سرچاپ حرف
    نویسه := u'س'
    سرچاپ نویسه
    مسیر := `C:\Users\Sara`
    سرچاپ مسیر
پایان
\end{salamfa}

\begin{outputfa}
A
س
C:\Users\Sara
\end{outputfa}

خط سوم خروجی را ببینید: در رشته‌ی خام میان \code{\textasciigrave} و
\code{\textasciigrave}، خط مورب وارونه دیگر آغاز یک دنباله‌ی گریز نیست و
همان‌طور که نوشته شده چاپ می‌شود؛ برای نوشتن یک مسیر ویندوزی یا یک الگوی
جست‌وجو که پر از \code{\textbackslash} است، این خیلی به کار می‌آید.

\begin{warn}[\code{'...'} با فارسی کار نمی‌کند]
\code{'...'} فقط \emph{یک بایت} می‌پذیرد، نه یک نویسه‌ی کامل. حرف‌های
انگلیسی هر کدام یک بایت‌اند، پس \fcode{'A'} درست است. اما هر حرف فارسی دو
بایت جا می‌گیرد، پس \fcode{'س'} به این خطا می‌خورد: «نویسه‌ای که میان
\code{'...'} نوشته می‌شود فقط یک بایت جا دارد». اگر
یک نویسه‌ی یونیکد (فارسی یا هر زبان دیگر) لازم دارید، به جای \code{'...'} از
\code{u'...'} استفاده کنید، مثل خط چهارم برنامه‌ی بالا.
\end{warn}

\section{بریدن بخشی از رشته}

\scode{زیررشته(آغاز, تعداد)} از جایگاه \fcode{آغاز} (که مثل آرایه‌ها از صفر
شمرده می‌شود) به تعداد \fcode{تعداد} نویسه برمی‌دارد:

\program{تکه‌های یک تاریخ}
\begin{salamfa}
روال ریشه:
    تاریخ := "1404/06/28"
    سال := تاریخ.زیررشته(۰, ۴)
    ماه := تاریخ.زیررشته(۵, ۲)
    روز := تاریخ.زیررشته(۸, ۲)
    سرچاپ "سال:", سال, "ماه:", ماه, "روز:", روز
    سرچاپ "سال بعد:", سال.به صحیح() + ۱
پایان
\end{salamfa}

\begin{outputfa}
سال: 1404 ماه: 06 روز: 28
سال بعد: 1405
\end{outputfa}

جایگاه‌ها را بشمارید: چهار رقم سال در جایگاه‌های ۰ تا ۳ هستند، بعد یک
خط مورب در جایگاه ۴، بعد ماه در ۵ و ۶، خط مورب در ۷ و روز در ۸ و ۹.
نتیجه‌ی \fcode{زیررشته} یک رشته‌ی تازه است؛ رشته‌ی اصلی دست‌نخورده می‌ماند.

با همین ابزار می‌توان یک شماره‌تلفن را بررسی کرد:

\program{بررسی شماره‌ی همراه}
\begin{salamfa}
روال ریشه:
    شماره := "09121234567"
    پیش شماره := شماره.زیررشته(۰, ۴)
    اگر (شماره.طول() == ۱۱) و (پیش شماره == "0912"):
        سرچاپ "شماره‌ی همراه اول است"
    وگرنه:
        سرچاپ "شماره‌ی ناشناخته"
    پایان
پایان
\end{salamfa}

\begin{outputfa}
شماره‌ی همراه اول است
\end{outputfa}

\section{شکافتن به واژه‌ها}

\scode{بشکاف(جداکننده)} رشته را در هر جایی که جداکننده هست می‌برد و
فهرستی از تکه‌ها می‌دهد. این فهرست شبیه آرایه است: عضوها را با
\scode{بگیر(اندیس)} (یا مثل آرایه با \scode{[اندیس]}) می‌خوانیم و با \fcode{هر}
هم می‌توان آن را پیمود:

\program{شمارش واژه‌های یک جمله}
\begin{salamfa}
روال ریشه:
    جمله := "زبان سلام برای همه است"
    کلمات := جمله.بشکاف(" ")
    سرچاپ "تعداد واژه‌ها:", کلمات.طول()
    تکرار کلمات.طول() در شماره:
        سرچاپ شماره + ۱, "-", کلمات.بگیر(شماره)
    پایان
پایان
\end{salamfa}

\begin{outputfa}
تعداد واژه‌ها: 5
1 - زبان
2 - سلام
3 - برای
4 - همه
5 - است
\end{outputfa}

شکافتن با فاصله روی متن فارسی هم درست کار می‌کند، چون فاصله همیشه یک
نویسه است. جداکننده می‌تواند هر رشته‌ای باشد: \scode{بشکاف("،")} برای
فهرستی که با ویرگول فارسی جدا شده، یا \scode{بشکاف("/")} برای همان تاریخ
بالا.

\section{پیراستن فاصله‌های اضافی}

متنی که از کاربر یا از یک پرونده می‌آید، اغلب فاصله‌های اضافی در آغاز و
پایان دارد. \scode{پیراست()} آن‌ها را می‌زداید:

\program{پیراستن}
\begin{salamfa}
روال ریشه:
    متن خام := "   قبول   "
    پاک := متن خام.پیراست()
    سرچاپ "[" + متن خام + "]"
    سرچاپ "[" + پاک + "]"
    سرچاپ پاک == "قبول"
پایان
\end{salamfa}

\begin{outputfa}
[   قبول   ]
[قبول]
true
\end{outputfa}

کروشه‌ها را برای این گذاشتیم که فاصله‌های نامرئی دیده شوند. بدون
پیراستن، مقایسه با \code{"قبول"} نادرست می‌شد، چون فاصله‌ها هم بخشی از
رشته‌اند. قاعده‌ی خوب: هر متنی را که از بیرون می‌آید، پیش از مقایسه
بپیرایید.

\section{جست‌وجو در رشته}

\scode{بیاب(متن)} می‌گوید متن داده‌شده در کجای رشته شروع می‌شود، و اگر
اصلاً در رشته نباشد، \scode{-۱} برمی‌گرداند. رایج‌ترین کاربردش همین پرسش
«آیا هست یا نه» است:

\program{جست‌وجو در پیام}
\begin{salamfa}
روال ریشه:
    پیام := "پرداخت با موفقیت انجام شد"
    اگر پیام.بیاب("موفقیت") != -۱:
        سرچاپ "عملیات موفق بود"
    وگرنه:
        سرچاپ "عملیات ناموفق بود"
    پایان
    سرچاپ پیام.بیاب("خطا")
پایان
\end{salamfa}

\begin{outputfa}
عملیات موفق بود
-1
\end{outputfa}

\section{رشته و عدد}

در فصل \ref{ch:types} دیدیم که \code{+} میان رشته و عدد، عدد را به رشته
تبدیل می‌کند و می‌چسباند، و \scode{به صحیح()} و \scode{به اعشار()} راه
عکس را می‌روند:

\program{از متن به عدد}
\begin{salamfa}
روال ریشه:
    قیمت متن := "25000"
    تعداد متن := "3"
    قیمت := قیمت متن.به صحیح()
    تعداد := تعداد متن.به صحیح()
    سرچاپ "مبلغ:", قیمت * تعداد
    وزن := "2.5".به اعشار()
    سرچاپ "وزن دو کارتن:", وزن * ۲.۰
پایان
\end{salamfa}

\begin{outputfa}
مبلغ: 75000
وزن دو کارتن: 5
\end{outputfa}

\begin{warn}
\scode{به صحیح()} و \scode{به اعشار()} هم رقم‌های انگلیسی و هم رقم‌های فارسی
را می‌فهمند: \scode{"۲۵۰۰۰".به صحیح()} برابر ۲۵۰۰۰ است. اما اگر رشته عدد
نباشد (مثل \code{"abc"})، نتیجه صفر است؛ پس پیش از تبدیل مطمئن شوید که متن
واقعاً عدد است.
\end{warn}

\section{ساختن رشته با حلقه}

در فصل \ref{ch:loops} رشته‌ای را قطعه‌قطعه ساختیم (مثلث ستاره‌ای). این کار را
می‌توان در روالی جمع کرد که هر رشته‌ای را هر تعداد بار تکرار کند:

\program{تکرار یک رشته}
\begin{salamfa}
روال چندبار(متن: رشته, تعداد: صحیح): رشته:
    ناپایا نتیجه := ""
    تکرار تعداد:
        نتیجه += متن
    پایان
    برگشت نتیجه
پایان

روال ریشه:
    سرچاپ چندبار("-=", ۱۰)
    سرچاپ چندبار("*", ۵) + " پنج ستاره " + چندبار("*", ۵)
    سرچاپ چندبار("-=", ۱۰)
پایان
\end{salamfa}

\begin{outputfa}
-=-=-=-=-=-=-=-=-=-=
***** پنج ستاره *****
-=-=-=-=-=-=-=-=-=-=
\end{outputfa}

\section{یک برنامه‌ی کامل‌تر}

تاریخی به شکل \code{"1404/06/28"} را می‌گیریم و آن را به شکل خوانا، با
نام ماه، چاپ می‌کنیم:

\program{تاریخ خوانا}
\begin{salamfa}
روال نام ماه(شماره: صحیح): رشته:
    برگشت همخوان شماره:
        ۱ => "فروردین"
        ۲ => "اردیبهشت"
        ۳ => "خرداد"
        ۴ => "تیر"
        ۵ => "مرداد"
        ۶ => "شهریور"
        ۷ => "مهر"
        ۸ => "آبان"
        ۹ => "آذر"
        ۱۰ => "دی"
        ۱۱ => "بهمن"
        ۱۲ => "اسفند"
        وگرنه => "نامعتبر"
    پایان
پایان

روال ریشه:
    تاریخ := "1404/06/28"
    سال := تاریخ.زیررشته(۰, ۴)
    ماه := تاریخ.زیررشته(۵, ۲).به صحیح()
    روز := تاریخ.زیررشته(۸, ۲).به صحیح()
    سرچاپ روز, نام ماه(ماه), سال
پایان
\end{salamfa}

\begin{outputfa}
28 شهریور 1404
\end{outputfa}

به خط \scode{تاریخ.زیررشته(۵, ۲).به صحیح()} دقت کنید: نتیجه‌ی
\fcode{زیررشته} یک رشته است و می‌توان بلافاصله قابلیت دیگری را روی آن
صدا زد. این «زنجیره کردن» در کار با رشته‌ها بسیار رایج است.

\begin{summary}
\begin{itemize}
  \item \code{+} رشته‌ها را می‌چسباند و \scode{طول()} تعداد نویسه‌ها را
        می‌دهد (هر حرف فارسی دو تا شمرده می‌شود).
  \item \code{\textbackslash n}، \code{\textbackslash t} و
        \code{\textbackslash "} نویسه‌های ویژه‌اند.
  \item \scode{زیررشته(آغاز, تعداد)} تکه‌ای را می‌برد،
        \scode{بشکاف(جداکننده)} رشته را به فهرستی از تکه‌ها می‌شکافد،
        \scode{پیراست()} فاصله‌های دو سر را می‌زداید و \scode{بیاب(متن)}
        جایگاه یک متن را (یا \scode{-۱}) می‌دهد.
  \item \scode{به صحیح()} و \scode{به اعشار()} رشته را به عدد تبدیل
        می‌کنند، چه رقم‌هایش فارسی باشد و چه انگلیسی.
  \item قابلیت‌ها را می‌توان زنجیره کرد: \scode{متن.زیررشته(۰, ۴).به صحیح()}.
\end{itemize}
\end{summary}

\section*{تمرین‌ها}
\markright{تمرین‌ها}

\exercise نام و نام خانوادگی را در دو متغیر بگذارید و با یک \fcode{سرچاپ}
سه خط چاپ کنید: نام کامل، نام خانوادگی به تنهایی، و جمله‌ای که هر دو را
در خود دارد. (راهنمایی: \code{\textbackslash n}.)

\exercise کد ملی ده‌رقمی را در رشته‌ای بگذارید و بررسی کنید که دقیقاً ده
نویسه دارد. سپس سه رقم اول آن را جدا کنید و چاپ کنید.

\exercise جمله‌ای را در متغیری بگذارید و درازترین واژه‌ی آن را پیدا و
چاپ کنید. (چون واژه‌ها فارسی‌اند، \scode{طول()} هر واژه دو برابر شمار
حرف‌های آن است؛ برای مقایسه‌ی درازی این مشکلی ایجاد نمی‌کند.)

\exercise فهرستی از نمره‌ها را به شکل یک رشته‌ی جداشده با ویرگول
(\code{"17,15,19,12"}) در متغیری بگذارید. آن را بشکافید، هر تکه را به
عدد تبدیل کنید و مجموع و میانگین را چاپ کنید.

\exercise روالی بنویسید که یک رشته و یک عرض بگیرد و اگر رشته کوتاه‌تر از
عرض بود، به سمت راست آن آن‌قدر فاصله بچسباند که به آن عرض برسد. با آن
یک جدول دوستونی مرتب از نام کالا و قیمت چاپ کنید. (از رشته‌های انگلیسی
یا رقم استفاده کنید تا \scode{طول()} دقیق باشد.)

\chapter{برنامه‌های کوچک کاربردی}\label{ch:projects}

\begin{goals}
\begin{itemize}
  \item کنار هم گذاشتن همه‌ی آموخته‌ها در چند برنامه‌ی کامل
  \item شکستن یک مسئله‌ی روزمره به قدم‌های کوچک
  \item خواندن یک برنامه‌ی چندروالی و فهمیدن جریان آن
  \item ایده‌هایی برای ادامه‌ی راه
\end{itemize}
\end{goals}

\section{از ابزار تا برنامه}

در فصل‌های گذشته ابزارهای پایه‌ی برنامه‌نویسی را یکی‌یکی شناختید: چاپ،
متغیر، گونه‌ها، محاسبه، شرط، حلقه، روال، آرایه و رشته. هر برنامه‌ای که در
جهان نوشته شده، هر قدر بزرگ، از همین ابزارها ساخته شده است. در این فصل
چند برنامه‌ی کوچک اما کامل می‌نویسیم؛ برنامه‌هایی که شبیه کارهای واقعی‌اند
و در هر کدام چند ابزار با هم کار می‌کنند.

برای هر برنامه اول مسئله را می‌گوییم، بعد آن را به قدم‌های کوچک می‌شکنیم و
بعد برنامه را می‌نویسیم. پیشنهاد می‌کنیم پیش از خواندن برنامه‌ی ما، خودتان
تلاش کنید آن را بنویسید؛ سپس دو نسخه را مقایسه کنید. هیچ اشکالی ندارد که
راه شما با راه ما فرق داشته باشد.

\section{قبض برق}

\textbf{مسئله.} شرکت برق مصرف ماهانه را پلکانی حساب می‌کند: ۱۰۰ کیلووات‌ساعت
اول هر واحد ۵۰۰ تومان، ۲۰۰ واحد بعدی هر واحد ۱٬۰۰۰ تومان و بیش از آن هر واحد
۲٬۰۰۰ تومان. مبلغ قبض چند مشترک را حساب کنید.

\textbf{قدم‌ها.} روالی می‌خواهیم که مصرف را بگیرد و مبلغ را برگرداند. مصرف
را پله‌پله کم می‌کنیم: اول تا ۱۰۰ واحد را با نرخ اول، بعد تا ۲۰۰ واحد را با
نرخ دوم و بقیه را با نرخ سوم. سپس برای آرایه‌ای از مصرف‌ها، مبلغ هر یک را
چاپ می‌کنیم.

\program{قبض برق پلکانی}
\begin{salamfa}
روال مبلغ قبض(مصرف: صحیح): صحیح:
    ناپایا باقی := مصرف
    ناپایا مبلغ := ۰

    // پله‌ی اول: تا ۱۰۰ واحد، هر واحد ۵۰۰ تومان
    پله یک := باقی > ۱۰۰ ؟ ۱۰۰ : باقی
    مبلغ += پله یک * ۵۰۰
    باقی -= پله یک

    // پله‌ی دوم: تا ۲۰۰ واحد بعدی، هر واحد ۱۰۰۰ تومان
    پله دو := باقی > ۲۰۰ ؟ ۲۰۰ : باقی
    مبلغ += پله دو * ۱۰۰۰
    باقی -= پله دو

    // پله‌ی سوم: هر چه ماند، هر واحد ۲۰۰۰ تومان
    مبلغ += باقی * ۲۰۰۰
    برگشت مبلغ
پایان

روال ریشه:
    مصرف مشترکان := [۸۰, ۱۰۰, ۲۵۰, ۴۰۰]
    هر مصرف در مصرف مشترکان:
        قبض := مبلغ قبض(مصرف)
        سرچاپ "مصرف", مصرف, "واحد:", قبض, "تومان"
    پایان
پایان
\end{salamfa}

\begin{outputfa}
مصرف 80 واحد: 40000 تومان
مصرف 100 واحد: 50000 تومان
مصرف 250 واحد: 200000 تومان
مصرف 400 واحد: 450000 تومان
\end{outputfa}

مصرف ۲۵۰ واحد را دنبال کنید: ۱۰۰ واحد در پله‌ی اول (۵۰ هزار)، ۱۵۰ واحد در
پله‌ی دوم (۱۵۰ هزار)، و چیزی برای پله‌ی سوم نمی‌ماند؛ جمعاً ۲۰۰ هزار تومان.
کارور \scode{؟ :} در هر پله می‌گوید «همه‌ی ظرفیت پله، یا هر چه مانده، هر کدام
که کمتر است».

\begin{tip}[ایده برای ادامه]
نرخ‌ها را به \fcode{پایا} تبدیل کنید و پله‌ی چهارمی اضافه کنید. سپس در
روال ریشه جمع کل درآمد شرکت را هم حساب کنید.
\end{tip}

\section{نمودار نمره‌ها}

\textbf{مسئله.} نمره‌های یک کلاس را داریم و می‌خواهیم ببینیم چند نفر در هر
بازه (زیر ۱۰، ۱۰ تا ۱۴، ۱۵ تا ۱۷، ۱۸ تا ۲۰) هستند و این را به شکل یک
نمودار ستاره‌ای نشان دهیم.

\textbf{قدم‌ها.} چهار شمارنده لازم داریم؛ برای هر نمره تصمیم می‌گیریم که به
کدام شمارنده اضافه شود. سپس برای هر بازه، یک خط چاپ می‌کنیم که به تعداد
شمارنده ستاره دارد. چاپ ستاره‌ها را به یک روال می‌سپاریم.

\program{نمودار ستاره‌ای نمره‌ها}
\begin{salamfa}
روال سطر نمودار(عنوان: رشته, تعداد: صحیح):
    ناپایا ستاره ها := ""
    تکرار تعداد:
        ستاره ها += "*"
    پایان
    سرچاپ عنوان + ":", ستاره ها, "(" + تعداد + ")"
پایان

روال ریشه:
    نمرات := [۱۲, ۱۹, ۸, ۱۵, ۲۰, ۱۷, ۱۱, ۹, ۱۶, ۱۸, ۱۴, ۱۳]
    ناپایا مردود := ۰
    ناپایا متوسط := ۰
    ناپایا خوب := ۰
    ناپایا عالی := ۰

    هر نمره در نمرات:
        اگر نمره < ۱۰: مردود++
        وگرنه نمره < ۱۵: متوسط++
        وگرنه نمره < ۱۸: خوب++
        وگرنه: عالی++
        پایان
    پایان

    سرچاپ "نمودار نمره‌های کلاس"
    سطر نمودار("زیر ۱۰   ", مردود)
    سطر نمودار("۱۰ تا ۱۴ ", متوسط)
    سطر نمودار("۱۵ تا ۱۷ ", خوب)
    سطر نمودار("۱۸ تا ۲۰ ", عالی)
پایان
\end{salamfa}

\begin{outputfa}
نمودار نمره‌های کلاس
زیر ۱۰   : ** (2)
۱۰ تا ۱۴ : **** (4)
۱۵ تا ۱۷ : *** (3)
۱۸ تا ۲۰ : *** (3)
\end{outputfa}

به شکل فشرده‌ی زنجیره‌ی شرط دقت کنید: هر شاخه یک دستور کوتاه دارد و در
همان خط نوشته شده و فقط یک \fcode{پایان} کل زنجیره را می‌بندد. برای
تراز شدن ستون ستاره‌ها، به عنوان‌ها فاصله‌ی اضافی چسبانده‌ایم.

\begin{tip}[ایده برای ادامه]
به جای چهار شمارنده‌ی جداگانه، یک آرایه با چهار صفر بسازید
و با اندیس، شمارنده‌ی مناسب را افزایش دهید. سپس با یک حلقه‌ی \fcode{هر} و
یک آرایه‌ی هم‌راستا از عنوان‌ها، نمودار را چاپ کنید.
\end{tip}

\section{صورت‌حساب فروشگاه}

\textbf{مسئله.} فهرست کالاهای سبد خرید یک مشتری را داریم: نام، قیمت واحد و
تعداد. باید صورت‌حسابی چاپ کنیم که مبلغ هر ردیف، جمع کل، تخفیف (اگر خرید
از ۵۰۰ هزار تومان بیشتر بود ۱۰ درصد) و مبلغ نهایی را نشان دهد.

\textbf{قدم‌ها.} سه آرایه‌ی هم‌راستا داریم. با یک حلقه روی نام‌ها می‌رویم،
مبلغ هر ردیف را حساب و چاپ می‌کنیم و به جمع کل می‌افزاییم. بعد از حلقه،
تخفیف را با یک شرط تعیین می‌کنیم.

\program{صورت‌حساب فروشگاه}
\begin{salamfa}
پایا آستانه := ۵۰۰۰۰۰
پایا تخفیف := ۱۰

روال ریشه:
    کالاها := ["برنج", "روغن", "چای", "شکر"]
    قیمتها := [۱۸۰۰۰۰, ۹۵۰۰۰, ۱۲۰۰۰۰, ۴۰۰۰۰]
    تعدادها := [۲, ۱, ۱, ۳]
    ناپایا جمع := ۰

    سرچاپ "صورت‌حساب"
    سرچاپ "-----------------------------"
    هر (شماره, کالا) در کالاها:
        مبلغ := قیمتها[شماره] * تعدادها[شماره]
        سرچاپ کالا, "x", تعدادها[شماره], "=", مبلغ
        جمع += مبلغ
    پایان
    سرچاپ "-----------------------------"
    سرچاپ "جمع کل:", جمع

    اگر جمع > آستانه:
        مبلغ تخفیف := (جمع * تخفیف) / ۱۰۰
        سرچاپ "تخفیف", تخفیف, "درصد:", مبلغ تخفیف
        سرچاپ "قابل پرداخت:", جمع - مبلغ تخفیف
    وگرنه:
        سرچاپ "قابل پرداخت:", جمع
        مانده := آستانه - جمع
        سرچاپ "با", مانده, "تومان خرید بیشتر، ۱۰ درصد تخفیف"
    پایان
پایان
\end{salamfa}

\begin{outputfa}
صورت‌حساب
-----------------------------
برنج x 2 = 360000
روغن x 1 = 95000
چای x 1 = 120000
شکر x 3 = 120000
-----------------------------
جمع کل: 695000
تخفیف 10 درصد: 69500
قابل پرداخت: 625500
\end{outputfa}

تعداد برنج و شکر را هر دو به ۱ تغییر دهید تا شاخه‌ی \fcode{وگرنه} را هم ببینید. متغیر
\fcode{مبلغ} درون حلقه و بدون \fcode{ناپایا} تعریف شده، چون در هر دور یک
بار مقدار می‌گیرد و همان است؛ فقط \fcode{جمع} است که در طول حلقه تغییر
می‌کند.

\begin{tip}[ایده برای ادامه]
یک آرایه‌ی چهارم برای «کالای تخفیف‌دار» (منطقی) اضافه کنید که روی آن
کالاها ۵ درصد تخفیف جداگانه اعمال شود. یا گران‌ترین ردیف صورت‌حساب را
پیدا کنید و در پایان چاپ کنید.
\end{tip}

\section{تقویم یک ماه}

\textbf{مسئله.} می‌خواهیم تقویم یک ماه را چاپ کنیم: هفت ستون از شنبه تا
جمعه، با روزهای ماه در جای درست. برای این کار باید بدانیم ماه چند روز
دارد و روز اول آن چه روزی از هفته است (۰ برای شنبه تا ۶ برای جمعه).

\textbf{قدم‌ها.} اول سرستون‌ها را چاپ می‌کنیم. بعد یک رشته می‌سازیم و به
تعداد روزهایی که پیش از روز اول ماه در هفته گذشته، جای خالی می‌گذاریم.
سپس روزها را یکی‌یکی می‌چسبانیم و هر بار که به جمعه رسیدیم، خط را چاپ
می‌کنیم و خط تازه‌ای شروع می‌کنیم. برای این که ستون‌ها تراز بمانند،
عددهای یک‌رقمی را با یک فاصله‌ی اضافی چاپ می‌کنیم.

\program{تقویم ماه}
\begin{salamfa}
روال دو رقمی(عدد: صحیح): رشته:
    اگر عدد < ۱۰: برگشت " " + عدد پایان
    برگشت "" + عدد
پایان

روال ریشه:
    تعداد روز := ۳۱
    // ۰ شنبه ... ۶ جمعه؛ این ماه از سه‌شنبه شروع می‌شود
    روز اول := ۳

    سرچاپ "شنبه یکشنبه دوشنبه سه‌شنبه چهارشنبه پنجشنبه جمعه"
    سرچاپ " ش   ی   د   س   چ   پ   ج"
    ناپایا خط := ""
    ناپایا ستون := ۰

    // جاهای خالی پیش از روز اول
    تکرار روز اول:
        خط += "    "
        ستون++
    پایان

    تکرار ۱ تا تعداد روز در روز:
        خط += دو رقمی(روز) + "  "
        ستون++
        اگر ستون == ۷:
            سرچاپ خط
            خط = ""
            ستون = ۰
        پایان
    پایان
    اگر ستون > ۰: سرچاپ خط پایان
پایان
\end{salamfa}

% The calendar is typeset as a table so its columns line up the way the
% editor's right-to-left output pane shows them.
\salamoutbox{%
شنبه یکشنبه دوشنبه سه‌شنبه چهارشنبه پنجشنبه جمعه\par
{\latincodefont\small
\begin{tabular}{@{}*{7}{>{\centering\arraybackslash}p{2em}}@{}}
{\salamcodeFAfont ش} & {\salamcodeFAfont ی} & {\salamcodeFAfont د} & {\salamcodeFAfont س} & {\salamcodeFAfont چ} & {\salamcodeFAfont پ} & {\salamcodeFAfont ج} \\
 & & & 1 & 2 & 3 & 4 \\
5 & 6 & 7 & 8 & 9 & 10 & 11 \\
12 & 13 & 14 & 15 & 16 & 17 & 18 \\
19 & 20 & 21 & 22 & 23 & 24 & 25 \\
26 & 27 & 28 & 29 & 30 & 31 & \\
\end{tabular}}}

شرط پایانی \scode{اگر ستون > ۰} برای این است که هفته‌ی ناتمام آخر هم
چاپ شود؛ بدون آن، روزهای ۲۶ تا ۳۱ در \fcode{خط} می‌ماندند و هرگز چاپ
نمی‌شدند. این نوع «کار ناتمام بعد از حلقه» یکی از جاهایی است که
برنامه‌نویسان زیاد فراموش می‌کنند.

\begin{tip}[ایده برای ادامه]
با \fcode{همخوان}، از روی شماره‌ی ماه تعداد روزهایش را به دست آورید، و روز
اول هر ماه را از روز اول ماه پیش حساب کنید (باقی‌مانده‌ی جمع بر ۷). آن وقت
می‌توانید تقویم کل سال را چاپ کنید.
\end{tip}

\section{تحلیل یک جمله}

\textbf{مسئله.} جمله‌ای داریم و می‌خواهیم بدانیم چند واژه دارد، درازترین
واژه‌اش کدام است و یک واژه‌ی مشخص چند بار در آن آمده است.

\textbf{قدم‌ها.} جمله را با فاصله می‌شکافیم. با یک حلقه روی واژه‌ها
می‌رویم: تعداد را می‌شماریم، درازترین را با الگوی «بیشترین» پیدا می‌کنیم و
هر جا واژه با واژه‌ی مورد نظر برابر بود، شمارنده‌ای را افزایش می‌دهیم.

\program{تحلیل جمله}
\begin{salamfa}
روال ریشه:
    جمله := "سلام زبان برنامه‌نویسی است و سلام یعنی درود"
    واژه هدف := "سلام"

    واژه ها := جمله.بشکاف(" ")
    ناپایا درازترین := ""
    ناپایا شمار هدف := ۰

    تکرار واژه ها.طول() در شماره:
        واژه := واژه ها.بگیر(شماره)
        اگر واژه.طول() > درازترین.طول():
            درازترین = واژه
        پایان
        اگر واژه == واژه هدف:
            شمار هدف++
        پایان
    پایان

    سرچاپ "تعداد واژه‌ها:", واژه ها.طول()
    سرچاپ "درازترین واژه:", درازترین
    سرچاپ "واژه‌ی «" + واژه هدف + "»", شمار هدف, "بار آمده است"
پایان
\end{salamfa}

\begin{outputfa}
تعداد واژه‌ها: 8
درازترین واژه: برنامه‌نویسی
واژه‌ی «سلام» 2 بار آمده است
\end{outputfa}

\fcode{درازترین} با رشته‌ی خالی شروع می‌شود که طولش صفر است؛ پس اولین
واژه حتماً جایش را می‌گیرد و از آن به بعد فقط واژه‌های درازتر. همان الگوی
«بیشترین» است، این بار برای رشته‌ها.

\begin{tip}[ایده برای ادامه]
کوتاه‌ترین واژه را هم پیدا کنید. سپس جمله‌ی تازه‌ای بسازید که در آن هر
واژه‌ی هدف با واژه‌ی دیگری (مثلاً «درود») جایگزین شده باشد.
\end{tip}

\section{و بعد از این؟}

اگر تا این‌جا آمده‌اید و برنامه‌های این فصل را فهمیده‌اید، دیگر
برنامه‌نویس هستید؛ شاید تازه‌کار، اما برنامه‌نویس. از این‌جا به بعد چند
راه پیش رو دارید:

\begin{itemize}
  \item \textbf{بیشتر بنویسید.} هر کار کوچک روزمره را که با عدد یا متن
        سروکار دارد، به برنامه تبدیل کنید: محاسبه‌ی سود وام، برنامه‌ی
        هفتگی درس، آمار مصرف آب خانه. تمرین‌های پایان فصل‌ها را دوباره
        ببینید و آن‌هایی را که رد کرده‌اید بنویسید.
  \item \textbf{برنامه‌های نمونه‌ی ویرایشگر برخط را بخوانید.} حالا همه‌ی
        آن‌ها را می‌فهمید. بعضی از آن‌ها چیزهایی را نشان می‌دهند که در این
        کتاب نبود، مثل \fcode{ساختار} برای دسته‌بندی داده‌های مرتبط، یا
        \fcode{وکتور} برای فهرست‌هایی که بزرگ و کوچک می‌شوند.
  \item \textbf{کتاب‌های بعدی سلام.} این کتاب پایه بود. سلام امکانات
        بسیار بیشتری دارد: کتابخانه‌ای بزرگ از ابزارهای آماده، ساختن
        صفحه‌های وب، کار با پرونده‌ها و شبکه، و بسیاری چیزهای دیگر. همه‌ی
        این‌ها بر همان پایه‌ای بنا می‌شوند که این‌جا آموختید.
\end{itemize}

مهم‌تر از همه: از اشتباه نترسید. هر پیام خطایی که می‌بینید، یک درس کوچک
است؛ و هر برنامه‌ای که بالاخره درست کار می‌کند، یک پیروزی کوچک. لذت
برنامه‌نویسی در همین پیروزی‌های کوچک است.

\begin{summary}
\begin{itemize}
  \item هر مسئله را به قدم‌های کوچک بشکنید و برای هر قدم روشن، یک روال
        یا یک بلوک بنویسید.
  \item الگوهایی که آموختید (شمارنده، جمع، بیشترین، جست‌وجو، ساختن رشته
        در حلقه) در برنامه‌های واقعی بارها تکرار می‌شوند.
  \item بعد از هر حلقه بپرسید: آیا کار ناتمامی مانده که باید انجام شود؟
  \item برنامه‌ای که کار می‌کند نقطه‌ی پایان نیست؛ آن را بهتر، کوتاه‌تر و
        خواناتر کنید.
\end{itemize}
\end{summary}

\section*{تمرین‌ها}
\markright{تمرین‌ها}

\exercise برنامه‌ای بنویسید که برای یک وام با مبلغ، نرخ سود سالانه و مدت
(به ماه)، قسط ماهانه‌ی ساده (مبلغ به علاوه‌ی سود کل، تقسیم بر تعداد ماه)
را حساب کند و جدول ماه‌به‌ماه مانده‌ی بدهی را چاپ کند.

\exercise دمای هر روز یک ماه را در آرایه‌ای بگذارید و گزارشی چاپ کنید
شامل: میانگین ماه، گرم‌ترین و سردترین روز، تعداد روزهای بالای ۳۰ درجه و
طولانی‌ترین رشته‌ی روزهای پیاپی بالای ۳۰ درجه.

\exercise فهرستی از نام‌ها و ساعت‌های کار هفتگی کارکنان را در دو آرایه‌ی
هم‌راستا بگذارید. با نرخ ساعتی ثابت، دستمزد هر نفر را حساب کنید؛ ساعت‌های
بیش از ۴۰ در هفته، اضافه‌کاری با نرخ یک‌ونیم برابر حساب می‌شود. جمع کل
دستمزدها را هم چاپ کنید.

\exercise یک «تبدیل‌گر واحد» بنویسید: روالی که مقداری و نام واحد مبدأ و
مقصد (مثلاً \code{"km"} و \code{"m"}) را بگیرد و با \fcode{همخوان} تبدیل
را انجام دهد. دست‌کم پنج جفت واحد را پشتیبانی کنید.

\exercise یک متن چندخطی (با \code{\textbackslash n}) را در رشته‌ای بگذارید،
آن را با \scode{بشکاف("\textbackslash n")} به خط‌ها بشکافید و برای هر خط،
شماره‌ی خط و تعداد واژه‌هایش را چاپ کنید. در پایان بگویید کدام خط بیشترین
واژه را دارد.


\beginappendices
\chapter{واژه‌های زبان سلام}
\label{app:keywords}

واژه‌های جدول نخست و بخش «واژه‌های دیگر زبان» برای خود زبان سلام رزرو
شده‌اند و نمی‌توان از آن‌ها به عنوان نام متغیر یا روال استفاده کرد؛ نه به
تنهایی و نه به عنوان بخشی از یک نام چندواژه‌ای (تنها استثنا \fcode{ریشه} است
که واژه‌ی رزرو شده نیست، اما چون نام روال آغاز برنامه است، متغیری هم‌نام آن
نمی‌توان ساخت). جدول نخست واژه‌هایی است که در این
کتاب به کار رفته‌اند و بخش «واژه‌های دیگر زبان» واژه‌هایی که در کتاب‌های بعدی
با آن‌ها آشنا می‌شوید.

\section{واژه‌های به‌کاررفته در این کتاب}

\begin{longtable}{lp{8.5cm}l}
\toprule
واژه & معنی و کاربرد & فصل \\
\midrule
\endhead
\fcode{روال} & تعریف روال؛ \fcode{روال ریشه} نقطه‌ی شروع برنامه & \ref{ch:first}، \ref{ch:functions} \\
\fcode{ریشه} & نام روال آغاز برنامه (رزرو شده نیست) & \ref{ch:first} \\
\fcode{پایان} & پایان هر بلوک (روال، شرط، حلقه) & \ref{ch:first} \\
\fcode{سرچاپ} & نوشتن روی خروجی و رفتن به خط بعد & \ref{ch:first} \\
\fcode{چاپ} & نوشتن روی خروجی بدون رفتن به خط بعد & \ref{ch:first} \\
\fcode{ناپایا} & تعریف متغیری که مقدارش عوض می‌شود & \ref{ch:variables} \\
\fcode{پایا} & تعریف یک پایا & \ref{ch:variables} \\
\fcode{درست} & مقدار منطقی درست & \ref{ch:types} \\
\fcode{نادرست} & مقدار منطقی نادرست & \ref{ch:types} \\
\fcode{برگردان} & تبدیل گونه & \ref{ch:types} \\
\fcode{و} & «و» منطقی & \ref{ch:operators} \\
\fcode{یا} & «یا» منطقی & \ref{ch:operators} \\
\fcode{وارونه} & «نه» منطقی (نقیض) & \ref{ch:operators} \\
\fcode{اگر} & شرط & \ref{ch:conditions} \\
\fcode{وگرنه} & راه دیگر شرط؛ با شرط یا بی‌شرط & \ref{ch:conditions} \\
\fcode{همخوان} & مقایسه‌ی یک مقدار با چند حالت & \ref{ch:conditions} \\
\fcode{تکرار} & حلقه با تعداد مشخص یا بازه & \ref{ch:loops} \\
\fcode{تا} & پایان بازه در \fcode{تکرار} & \ref{ch:loops} \\
\fcode{هر} & اندازه‌ی پرش در \fcode{تکرار} (\fcode{تکرار ۱ تا ۱۰ هر ۲ در ک}) & \ref{ch:loops} \\
\fcode{تا} & حلقه تا وقتی شرط برقرار است & \ref{ch:loops} \\
\fcode{بشکن} & پایان دادن به حلقه & \ref{ch:loops} \\
\fcode{گذر} & پریدن به دور بعدی حلقه & \ref{ch:loops} \\
\fcode{برگشت} & پس دادن نتیجه از روال & \ref{ch:functions} \\
\fcode{هر} & پیمایش عضوهای آرایه & \ref{ch:arrays} \\
\fcode{در} & مشخص کردن آرایه در حلقه‌ی \fcode{هر} و نام‌گذاری شمارنده‌ی \fcode{تکرار} & \ref{ch:arrays}، \ref{ch:loops} \\
\fcode{برابر} & مقایسه‌ی برابری، معادل \fcode{==} & \ref{ch:operators} \\
\fcode{نابرابر} & مقایسه‌ی نابرابری، معادل \fcode{!=} & \ref{ch:operators} \\
\bottomrule
\end{longtable}
\chaptermark{واژه‌های زبان سلام}

\section{نام گونه‌ها}

\begin{center}
\begin{tabular}{ll}
\toprule
نام & گونه \\
\midrule
\fcode{صحیح} & عدد صحیح \\
\fcode{اعشار۶۴} & عدد اعشاری \\
\fcode{رشته} & متن \\
\fcode{منطقی} & درست یا نادرست \\
\bottomrule
\end{tabular}
\end{center}

سلام گونه‌های دیگری هم دارد (مثل \fcode{صحیح۶۴} برای عددهای بسیار بزرگ،
\fcode{نویسه} برای یک نویسه‌ی تک‌بایتی مانند یک حرف انگلیسی، و \fcode{یونیکد}
برای یک حرف فارسی) که در این کتاب به آن‌ها نیازی نبود. نام گونه‌ها رزرو شده
نیستند.

\section{واژه‌های دیگر زبان}

این واژه‌ها هم رزرو شده‌اند و نباید در نام‌ها به کار روند:
\fcode{ساختار}، \fcode{جداشمار}، \fcode{گونه}، \fcode{واردسازی}، \fcode{بسته}،
\fcode{پوچ}، \fcode{این}، \fcode{چیدمان}، \fcode{بر}، \fcode{همگانی}،
\fcode{میانجی}، \fcode{کاربست}، \fcode{بخش}، \fcode{ترابرد}،
\fcode{فراخوانی}، \fcode{درون‌داد}، \fcode{برون‌داد}، \fcode{دیرکن}، \fcode{ورودی}، \fcode{نادرست‌چاپ}،
\fcode{نادرست‌سرچاپ}، \fcode{ناب}، \fcode{درخط}، \fcode{نادرخط}،
\fcode{نابرگشت} و \fcode{بی‌کاره}.

واژه‌های \fcode{پیوند}، \fcode{ایستا}، \fcode{پویا} و \fcode{چارچوب} رزرو شده
نیستند؛ آن‌ها فقط در دستوری که کتابخانه‌های بیرونی را به برنامه پیوند می‌دهد
معنای ویژه دارند و در جاهای دیگر می‌توانید از آن‌ها برای نام‌گذاری استفاده کنید.

\begin{note}
اگر نامی که انتخاب کرده‌اید خطای «کلمه‌ی رزرو شده» یا خطای دستوری عجیبی
مثل «یک عبارت انتظار می‌رفت» گرفت، شاید یکی از واژه‌هایش رزرو شده باشد.
ساده‌ترین راه این است که آن را کمی تغییر دهید: به جای \fcode{بسته} بنویسید
\fcode{کارتن} یا \fcode{جعبه}، به جای \fcode{قیمت هر واحد} بنویسید
\fcode{قیمت هرواحد}.
\end{note}

\section{قابلیت‌های رشته‌ها و آرایه‌ها}

این‌ها واژه‌ی رزرو شده نیستند، اما با نقطه به رشته‌ها و آرایه‌ها می‌چسبند و
در این کتاب به کار رفته‌اند:

\begin{center}
\begin{tabular}{ll}
\toprule
قابلیت & کار \\
\midrule
\scode{طول()} & تعداد خانه‌های رشته (فصل \ref{ch:strings}) یا عضوهای آرایه \\
\scode{زیررشته(آغاز, تعداد)} & بریدن تکه‌ای از رشته \\
\scode{بشکاف(جداکننده)} & شکافتن رشته به فهرستی از تکه‌ها \\
\scode{بگیر(اندیس)} & خواندن یک عضو از فهرست حاصل از \fcode{بشکاف} \\
\scode{پیراست()} & زدودن فاصله‌های آغاز و پایان \\
\scode{بیاب(متن)} & جایگاه یک متن در رشته، یا \scode{-۱} \\
\scode{به صحیح()} & تبدیل رشته به عدد صحیح \\
\scode{به اعشار()} & تبدیل رشته به عدد اعشاری \\
\bottomrule
\end{tabular}
\end{center}

\chapter{خطاهای رایج و راه حل آن‌ها}
\label{app:errors}

پیام خطا دوست شماست؛ اما در آغاز، خواندنش کمی تمرین می‌خواهد. پیام خطایی که
سلام پیش از اجرای برنامه می‌دهد معمولاً سه بخش دارد: خود پیام، جای خطا (نام پرونده، شماره‌ی خط و شماره‌ی ستون؛
مثلاً \code{main.salam:3:5} یعنی خط ۳، ستون ۵)، و
تکه‌ای از همان خط با یک نشانه‌ی \code{\textasciicircum} زیر جایی که مشکل
دیده شده است. بیشتر پیام‌ها شماره‌ای هم دارند که در کروشه، پس از واژه‌ی
«خطا» می‌آید؛ مثلاً \scode{خطا[خطا۰۰۱]}. برخی پیام‌ها هم در پایان، خطی با
\scode{= راهنما:} دارند که پیشنهاد مشخصی برای رفع خطا می‌دهد. هر پیامی که با
\fcode{هشدار} آغاز شود، خطا نیست: برنامه اجرا می‌شود، اما سلام می‌گوید جایی از
آن احتمالاً اشتباه است. خطاهای زمان اجرا، مثل تمام شدن زمان اجرا، هنگام
اجرای برنامه پیش می‌آیند و جای خطا و نشانه‌ی \code{\textasciicircum} ندارند.
در این پیوست رایج‌ترین پیام‌هایی را که تازه‌کارها
می‌بینند، همراه با علت و راه حل آورده‌ایم؛ برای کوتاهی، تکه‌ی برنامه و
\code{\textasciicircum} را ننوشته‌ایم.

\section*{رشته‌ی پایان‌نیافته}

\begin{progoutfa}
خطا: رشته‌ی پایان‌نیافته
\end{progoutfa}

\textbf{علت.} نقل قول پایانی یک رشته فراموش شده است.
\textbf{راه حل.} در خطی که پیام می‌گوید، هر رشته را نگاه کنید و مطمئن شوید
با \code{"} باز و بسته شده است.

\section*{\scode{'پایان'} برای بستن بلوک انتظار می‌رفت}

\begin{progoutfa}
خطا: 'پایان' برای بستن بلوک انتظار می‌رفت (در پایان پرونده)
\end{progoutfa}

\textbf{علت.} یک بلوک (روال، \fcode{اگر}، حلقه) با دونقطه باز شده اما
\fcode{پایان} آن نوشته نشده است. معمولاً شماره‌ی خط به پایان پرونده اشاره
می‌کند، چون سلام تا آخر دنبال \fcode{پایان} گشته و پیدا نکرده است.
\textbf{راه حل.} تعداد دونقطه‌های آغاز بلوک و تعداد \fcode{پایان}‌ها را
بشمارید. تورفتگی درست، پیدا کردن جای گمشده را آسان می‌کند.

\section*{\scode{':'} برای باز کردن بلوک انتظار می‌رفت}

\begin{progoutfa}
خطا: ':' برای باز کردن بلوک انتظار می‌رفت (در پایان خط)
\end{progoutfa}

\textbf{علت.} دونقطه‌ی پایان خط \fcode{روال ریشه}، \fcode{تکرار} یا \fcode{تا}
فراموش شده است. اگر دونقطه‌ی پس از شرط \fcode{اگر} یا \fcode{وگرنه} جا بیفتد،
پیام کمی فرق دارد:

\begin{progoutfa}
خطا: ':' پس از شرط اگر انتظار می‌رفت (در پایان خط)
\end{progoutfa}

\textbf{راه حل.} دونقطه را در پایان همان خط بگذارید.

\section*{شناسه‌ی ناشناخته}

\begin{progoutfa}
خطا[خطا۰۰۱]: شناسه‌ی ناشناخته 'سنن'
\end{progoutfa}

\textbf{علت.} نامی به کار رفته که تعریف نشده است. بیشتر وقت‌ها یک غلط
املایی است (\fcode{سنن} به جای \fcode{سن})، یا متغیری که درون بلوک
دیگری تعریف شده و این‌جا در دسترس نیست.
\textbf{راه حل.} املای نام را با تعریفش مقایسه کنید. اگر متغیر درون یک
حلقه یا روال دیگر تعریف شده، آن را به جای مناسب منتقل کنید.

\section*{روال فراخوانی نشده است}

\begin{progoutfa}
خطا[خطا۱۱۰]: روال 'خط جداکننده' فراخوانی نشده است: برای فراخوانی آن بنویسید 'خط جداکننده()'
\end{progoutfa}

\textbf{علت.} نام روال بدون پرانتز نوشته شده است.
\textbf{راه حل.} پس از نام روال پرانتز بگذارید: \mbox{\scode{خط جداکننده()}}.

\section*{میان مقدارهایی که چاپ می‌شوند \scode{','} انتظار می‌رفت}

\begin{progoutfa}
خطا[خطا۱۱۰]: میان مقدارهایی که چاپ می‌شوند ',' انتظار می‌رفت
\end{progoutfa}

\textbf{علت.} در یک \fcode{سرچاپ} یا \fcode{چاپ} ویرگول میان دو مقدار فراموش
شده است؛ مثلاً \mbox{\scode{سرچاپ "سن:" سن}}.
\textbf{راه حل.} میان هر دو مقدار ویرگول بگذارید: \mbox{\scode{سرچاپ "سن:", سن}}.
همین شماره‌ی خطا برای هر دستوری هم می‌آید که هیچ کاری نمی‌کند؛ مثلاً خطی که
فقط \scode{الف + ۱} است و نتیجه‌اش را در جایی نمی‌گذارد.

\section*{نام نمی‌تواند با رقم آغاز شود}

\begin{progoutfa}
خطا: نام نمی‌تواند با رقم آغاز شود؛ ...
\end{progoutfa}

\textbf{علت.} نام متغیری مثل \fcode{۲نفر} با رقم شروع شده است.
\textbf{راه حل.} رقم را پس از یک حرف بگذارید: \fcode{نفر۲}.

\section*{متغیر استفاده‌نشده}

\begin{progoutfa}
خطا[خطا۰۵۹]: متغیر استفاده‌نشده 'سن': ...
\end{progoutfa}

\textbf{علت.} متغیری تعریف شده اما هیچ‌جا خوانده نشده است.
\textbf{راه حل.} یا از آن استفاده کنید، یا خط تعریفش را حذف کنید. اگر
فقط برای آزمایش موقت آن را نگه می‌دارید، می‌توانید نامش را با
\code{\_} شروع کنید تا سلام از آن بگذرد؛ اما همان‌طور که خود پیام می‌گوید،
این کار توصیه نمی‌شود و پیش از پایان کار باید آن خط را پاک کنید.

\section*{روال استفاده نشده است}

\begin{progoutfa}
خطا[خطا۰۶۶]: روال 'خط جداکننده' استفاده نشده است: خصوصی است و هیچ‌جا فراخوانی نمی‌شود؛ ...
\end{progoutfa}

\textbf{علت.} روالی تعریف شده اما هیچ‌جا فراخوانی نشده است.
\textbf{راه حل.} آن را در جایی فراخوانی کنید یا حذفش کنید. فراموش نکنید
که فراخوانی پرانتز می‌خواهد: \mbox{\scode{خط جداکننده()}}.

\section*{نمی‌توان متغیر تغییرناپذیر را دوباره مقداردهی کرد}

\begin{progoutfa}
خطا[خطا۰۱۳]: نمی‌توان متغیر تغییرناپذیر 'سن' را دوباره مقداردهی کرد؛ ...
\end{progoutfa}

\textbf{علت.} به متغیری که بدون \fcode{ناپایا} تعریف شده، با \code{=} مقدار
تازه داده‌اید.
\textbf{راه حل.} در خط تعریف، پیش از نام، \fcode{ناپایا} بنویسید. اگر واقعاً
قرار نبوده عوض شود، خطی را که مقدار تازه می‌دهد بازبینی کنید؛ شاید
اشتباه از آن‌جاست.

\section*{کلمه‌ی رزرو شده}

\begin{progoutfa}
خطا: 'بسته' یک کلمه‌ی رزرو شده است و نمی‌توان آن را به عنوان نام به کار برد
  = راهنما: نام دیگری انتخاب کنید، برای نمونه 'بسته_'
\end{progoutfa}

\textbf{علت.} نام متغیر یا روال، یا یکی از واژه‌های آن، جزو واژه‌های خود
زبان است (پیوست \ref{app:keywords}).
\textbf{راه حل.} نام دیگری انتخاب کنید یا دو واژه را به هم بچسبانید.

\section*{نمی‌توان گونه‌ی الف را به گونه‌ی ب نسبت داد}

\begin{progoutfa}
خطا[خطا۰۰۲]: نمی‌توان 'رشته' را به 'صحیح' نسبت داد
\end{progoutfa}

\textbf{علت.} مقداری از یک گونه در جای مقداری از گونه‌ی دیگر گذاشته شده: رشته
در متغیر عددی، عدد اعشاری در متغیر صحیح، آرگومان با گونه‌ی اشتباه به روال.
نام گونه‌ها در پیام، همان نام‌هایی است که در برنامه می‌نویسید:
\fcode{رشته}، \fcode{صحیح}، \fcode{اعشار۶۴} و \fcode{منطقی}. اگر آرگومانی با
گونه‌ی اشتباه به روال بدهید، پیام مشابهی با شماره‌ی \fcode{خطا۰۱۲} می‌گیرید:
«آرگومان ۱: نمی‌توان \scode{'رشته'} را به \scode{'صحیح'} پاس داد».
\textbf{راه حل.} گونه‌ها را یکسان کنید: با \fcode{برگردان}، با
\scode{به صحیح()} یا با اصلاح مقدار.

\section*{شرط باید بولی باشد}

\begin{progoutfa}
خطا[خطا۰۲۱]: شرط 'اگر' باید بولی باشد، اما 'صحیح' دریافت شد
\end{progoutfa}

\textbf{علت.} در \fcode{اگر}، \fcode{تا} یا کارور \scode{؟ :} چیزی نوشته
شده که مقدار منطقی نیست؛ مثلاً یک عدد.
\textbf{راه حل.} شرط را به یک مقایسه تبدیل کنید: به جای \scode{اگر تعداد:}
بنویسید \scode{اگر تعداد > ۰:}.

\section*{مقداردهی را نمی‌توان به‌جای یک مقدار به کار برد}

\begin{progoutfa}
خطا[خطا۰۱۳]: مقداردهی را نمی‌توان به‌جای یک مقدار به کار برد؛
برای مقایسه‌ی دو مقدار به‌جای '=' بنویسید '==' یا 'برابر'
\end{progoutfa}

\textbf{علت.} در شرط، به جای \code{==} یک \code{=} نوشته شده است؛ مثلاً
\scode{اگر سن = ۲۰:}. \code{=} مقدار می‌دهد و پرسشی نمی‌پرسد.
\textbf{راه حل.} بنویسید \scode{اگر سن == ۲۰:} یا
\scode{اگر سن برابر ۲۰:}.

\section*{هشدار: روال ممکن است بدون بازگشت تمام شود}

\begin{progoutfa}
هشدار[هشدار۰۷۱]: روال 'قدرمطلق' ممکن است بدون بازگرداندن مقداری از گونه 'صحیح'
به پایان بدنه‌ی خود برسد؛ یک 'برگشت' نهایی اضافه کنید
\end{progoutfa}

\textbf{علت.} روالی که گونه‌ی بازگشتی اعلام کرده، در یکی از مسیرها (مثلاً
وقتی هیچ شرطی برقرار نیست) به \fcode{برگشت} نمی‌رسد. این یک هشدار است، نه
خطا: برنامه اجرا می‌شود، اما اگر اجرا به همان مسیر برسد، روال مقدار درستی
برنمی‌گرداند.
\textbf{راه حل.} در پایان روال یک \fcode{برگشت} برای حالت پیش‌فرض
بگذارید.

\section*{برنامه اجرا می‌شود اما تمام نمی‌شود}

\begin{progoutfa}
خطای زمان اجرا: زمان اجرا به پایان رسید (احتمال حلقه‌ی بی‌نهایت)
\end{progoutfa}

\textbf{علت.} حلقه‌ی بی‌پایان: شرط یک \fcode{تا} هرگز نادرست نمی‌شود.
ویرایشگر برخط برنامه‌ای را که بیش از سه ثانیه طول بکشد، با این پیام متوقف
می‌کند.
\textbf{راه حل.} در حلقه به دنبال متغیری بگردید که باید در هر دور تغییر
می‌کرده و نکرده است. اگر شمارنده‌ی حلقه را با \fcode{ناپایا} تعریف کرده اما
هیچ‌جا تغییرش نداده باشید، سلام پیش از اجرا خبرتان می‌کند:

\begin{progoutfa}
خطا[خطا۰۶۹]: متغیر 'شمارنده' با کلیدواژه‌ی 'ناپایا' تعریف شده اما هرگز تغییر نکرده است؛ ...
\end{progoutfa}

\section*{برنامه اجرا می‌شود اما خروجی اشتباه است}

سخت‌ترین نوع اشکال، اشکالی است که سلام نمی‌تواند تشخیص دهد، چون برنامه
از نظر زبان درست است و فقط چیزی غیر از آنچه شما می‌خواستید انجام می‌دهد.
چند راهنمایی:
\begin{itemize}
  \item در چند جای برنامه \fcode{سرچاپ} موقت بگذارید و مقدار متغیرها را در
        میانه‌ی کار ببینید. بعد از پیدا کردن اشکال، آن‌ها را پاک کنید.
  \item به تقسیم عددهای صحیح مشکوک باشید (\scode{۷ / ۲} برابر \code{۳}
        است) و به ترتیب انجام کارور‌ها.
  \item مطمئن شوید که در شرط‌ها \code{==} نوشته‌اید، نه \code{=}.
  \item در حلقه‌ها بررسی کنید که شمارنده از کجا شروع و کجا تمام می‌شود؛
        اشتباه‌های «یکی کم یا یکی زیاد» بسیار رایج‌اند.
  \item مطمئن شوید که متغیرهای جمع‌کننده بیرون حلقه تعریف شده‌اند.
  \item برنامه را با کوچک‌ترین ورودی ممکن (یک عضو، صفر، عدد منفی)
        امتحان کنید.
\end{itemize}

\chapter{پاسخ تمرین‌های منتخب}
\label{app:solutions}

در این پیوست پاسخ شماری از تمرین‌های کتاب آمده است. پاسخ‌ها یک راه حل‌اند،
نه تنها راه حل؛ اگر برنامه‌ی شما خروجی درست می‌دهد و خوانا است، درست است،
حتی اگر با این‌جا فرق داشته باشد. خواهش می‌کنیم پیش از نگاه کردن به هر
پاسخ، خودتان تلاش کنید.

\section*{فصل \ref{ch:first}}

\subsection*{تمرین ۱.۱}
\begin{salamfa}
روال ریشه:
    سرچاپ "سارا"
    سرچاپ "احمدی"
    سرچاپ "کاشان"
پایان
\end{salamfa}

\subsection*{تمرین ۱.۳}
\begin{salamfa}
روال ریشه:
    سرچاپ "حاصل جمع ۱۲۵ و ۳۷۵:", ۱۲۵ + ۳۷۵
    سرچاپ "حاصل ضرب ۲۴ در ۷:", ۲۴ * ۷
    سرچاپ "ساعت‌های یک هفته:", ۷ * ۲۴
پایان
\end{salamfa}
\begin{outputfa}
حاصل جمع ۱۲۵ و ۳۷۵: 500
حاصل ضرب ۲۴ در ۷: 168
ساعت‌های یک هفته: 168
\end{outputfa}

\subsection*{تمرین ۱.۴}
سه اشتباه: دونقطه‌ی بعد از \fcode{روال ریشه} فراموش شده، در خط دوم
\fcode{print} انگلیسی به جای \fcode{سرچاپ} نوشته شده، و رشته‌ی خط سوم نقل قول
پایانی ندارد.

\section*{فصل \ref{ch:variables}}

\subsection*{تمرین ۲.۱}
\begin{salamfa}
روال ریشه:
    طول := ۵
    عرض := ۴
    ارتفاع := ۳
    سرچاپ "مساحت کف:", طول * عرض, "متر مربع"
    سرچاپ "حجم اتاق:", طول * عرض * ارتفاع, "متر مکعب"
پایان
\end{salamfa}
\begin{outputfa}
مساحت کف: 20 متر مربع
حجم اتاق: 60 متر مکعب
\end{outputfa}

\subsection*{تمرین ۲.۳}
خروجی \code{۷ ۱} است. دنبال کنید: \fcode{الف} برابر ۳ و \fcode{ب} برابر
۶ می‌شود. سپس \fcode{الف} برابر ۷ می‌شود. در آخر \fcode{ب} برابر
\code{۷ - ۶} یعنی ۱ می‌شود.

\subsection*{تمرین ۲.۴}
\fcode{امتیاز} بدون \fcode{ناپایا} تعریف شده اما در خط بعد عوض می‌شود؛ باید
\fcode{ناپایا امتیاز \lr{:=} ۱۰} باشد. \fcode{نام بازیکن} تعریف شده اما
استفاده نشده؛ یا در \fcode{سرچاپ} به کارش ببرید یا خطش را حذف کنید.

\section*{فصل \ref{ch:types}}

\subsection*{تمرین ۳.۳}
\begin{salamfa}
روال ریشه:
    دقیقه ها := ۱۳۵
    ساعت := دقیقه ها / ۶۰
    دقیقه := دقیقه ها % ۶۰
    سرچاپ ساعت, "ساعت و", دقیقه, "دقیقه"
پایان
\end{salamfa}
\begin{outputfa}
2 ساعت و 15 دقیقه
\end{outputfa}

\section*{فصل \ref{ch:operators}}

\subsection*{تمرین ۴.۲}
\begin{salamfa}
روال ریشه:
    عدد := ۴۷۲۹
    یکان := عدد % ۱۰
    دهگان := (عدد / ۱۰) % ۱۰
    صدگان := (عدد / ۱۰۰) % ۱۰
    هزارگان := عدد / ۱۰۰۰
    سرچاپ "مجموع رقم‌ها:", یکان + دهگان + صدگان + هزارگان
پایان
\end{salamfa}
\begin{outputfa}
مجموع رقم‌ها: 22
\end{outputfa}

\section*{فصل \ref{ch:conditions}}

\subsection*{تمرین ۵.۳}
\begin{salamfa}
روال ریشه:
    ساعت := ۱۵
    اگر ساعت < ۱۲:
        سرچاپ "صبح بخیر"
    وگرنه ساعت < ۱۷:
        سرچاپ "ظهر بخیر"
    وگرنه ساعت < ۲۰:
        سرچاپ "عصر بخیر"
    وگرنه:
        سرچاپ "شب بخیر"
    پایان
پایان
\end{salamfa}
\begin{outputfa}
ظهر بخیر
\end{outputfa}

\subsection*{تمرین ۵.۶}
سه خط چاپ می‌شود: «بزرگ‌تر از ده»، «مضرب پنج» و «مضرب سه». زنجیره‌ی اول
با نخستین شرط درست تمام می‌شود و «بزرگ‌تر از پنج» چاپ نمی‌شود؛ اما دو
\fcode{اگر} بعدی جداگانه‌اند و هر دو بررسی می‌شوند.

\section*{فصل \ref{ch:loops}}

\subsection*{تمرین ۶.۱}
\begin{salamfa}
روال ریشه:
    ناپایا خط := ""
    تکرار ۱ تا ۵۰ در عدد:
        اگر (عدد % ۷) == ۰:
            خط = خط + عدد + " "
        پایان
    پایان
    سرچاپ خط
پایان
\end{salamfa}
\begin{outputfa}
7 14 21 28 35 42 49
\end{outputfa}

\subsection*{تمرین ۶.۲}
\begin{salamfa}
روال ریشه:
    ناپایا حاصل := ۱
    تکرار ۱ تا ۱۰ در عدد:
        حاصل *= عدد
    پایان
    سرچاپ حاصل
پایان
\end{salamfa}
\begin{outputfa}
3628800
\end{outputfa}

\subsection*{تمرین ۶.۶}
شش خط. حلقه‌ی بیرونی سه دور دارد و در هر دور، حلقه‌ی درونی فقط برای
\fcode{ب} برابر ۰ و ۱ چاپ می‌کند و در ۲ می‌شکند.

\section*{فصل \ref{ch:functions}}

\subsection*{تمرین ۷.۱}
\begin{salamfa}
روال بزرگتر(الف: صحیح, ب: صحیح): صحیح:
    اگر الف > ب: برگشت الف پایان
    برگشت ب
پایان

روال ریشه:
    دوتای_اول := بزرگتر(۱۲, ۲۵)
    سرچاپ بزرگتر(دوتای_اول, ۱۹)
پایان
\end{salamfa}
\begin{outputfa}
25
\end{outputfa}
بزرگ‌ترین سه عدد، بزرگ‌تر از «بزرگ‌تر دو تای اول» و سومی است.

\subsection*{تمرین ۷.۶}
وقتی \fcode{عدد} صفر است، هیچ‌کدام از دو شرط برقرار نیست و روال بدون
\fcode{برگشت} تمام می‌شود؛ سلام این را خطا می‌گیرد. راه حل: شرط دوم را
به \fcode{وگرنه:} بی‌شرط تبدیل کنید، یا در پایان روال
\fcode{برگشت عدد} بگذارید.

\section*{فصل \ref{ch:arrays}}

\subsection*{تمرین ۸.۲}
\begin{salamfa}
روال میانگین(اعداد: صحیح[:]): اعشار۶۴:
    ناپایا جمع := ۰
    هر عدد در اعداد:
        جمع += عدد
    پایان
    تعداد := اعداد.طول()
    برگشت (جمع برگردان اعشار۶۴) / (تعداد برگردان اعشار۶۴)
پایان

روال بیشترین(اعداد: صحیح[:]): صحیح:
    ناپایا بزرگترین := اعداد[۰]
    هر عدد در اعداد:
        اگر عدد > بزرگترین: بزرگترین = عدد پایان
    پایان
    برگشت بزرگترین
پایان

روال ریشه:
    نمرات := [۱۷, ۱۵, ۱۹, ۱۲, ۲۰]
    سرچاپ "میانگین:", میانگین(نمرات[:])
    سرچاپ "بیشترین:", بیشترین(نمرات[:])
پایان
\end{salamfa}
\begin{outputfa}
میانگین: 16.6
بیشترین: 20
\end{outputfa}

\subsection*{تمرین ۸.۶}
خروجی \code{[۱۰, ۲, ۳۰, ۴, ۵۰]} است: عضوهایی که اندیس زوج دارند (۰، ۲
و ۴) ده برابر می‌شوند.

\section*{فصل \ref{ch:strings}}

\subsection*{تمرین ۹.۴}
\begin{salamfa}
روال ریشه:
    متن := "17,15,19,12"
    تکه ها := متن.بشکاف(",")
    ناپایا مجموع := ۰
    تکرار تکه ها.طول() در ک:
        مجموع += تکه ها.بگیر(ک).به صحیح()
    پایان
    تعداد := تکه ها.طول()
    سرچاپ "مجموع:", مجموع
    سرچاپ "میانگین:", (مجموع برگردان اعشار۶۴) / (تعداد برگردان اعشار۶۴)
پایان
\end{salamfa}
\begin{outputfa}
مجموع: 63
میانگین: 15.75
\end{outputfa}

\chapter*{سخن پایانی}
\addcontentsline{toc}{chapter}{سخن پایانی}
\markboth{سخن پایانی}{سخن پایانی}

این کتاب با یک جمله شروع شد: «سلام، دنیا!». آن روز یک ویرایشگر خالی بود و
سه خط کد. امروز می‌دانید یک برنامه چگونه مقدارها را در حافظه نگه می‌دارد،
چگونه حساب می‌کند و تصمیم می‌گیرد، چگونه یک کار را بی‌خستگی تکرار می‌کند و
چگونه کار بزرگ را میان روال‌های کوچک و نام‌دار تقسیم می‌کند. با آرایه‌ها
مجموعه‌ای از داده‌ها را و با رشته‌ها واژه‌ها و جمله‌ها را در دست گرفته‌اید و
در فصل آخر همه‌ی این‌ها را کنار هم گذاشتید تا برنامه‌هایی بسازید که به
کار می‌آیند.

مهم‌تر از هر ابزار، شیوه‌ی نگاه است که آموختید: مسئله‌ی بزرگ را به قدم‌های
کوچک بشکنید، هر قدم را جداگانه آزمایش کنید، و وقتی برنامه درست کار نکرد،
به پیام خطا به چشم یک راهنما نگاه کنید، نه یک ایراد.

\section*{این کتاب همچنان کنار شماست}

سه پیوست پایانی برای همین ساخته شده‌اند که پس از خواندن کتاب هم به کارتان
بیایند. هر وقت مطمئن نبودید یک واژه رزرو شده است یا نه، به پیوست
\ref{app:keywords} سر بزنید. هر وقت برنامه‌تان پیامی داد که نمی‌فهمیدید،
پیوست \ref{app:errors} را باز کنید. و هر وقت تمرینی را نوشتید و خواستید
راه‌حل خود را با راه‌حلی دیگر بسنجید، پاسخ‌ها در پیوست \ref{app:solutions}
هستند. یادتان باشد که یک تمرین بیش از یک پاسخ درست دارد.

و اگر می‌خواهید همین حالا چیزی بنویسید، ویرایشگر برخط سلام همیشه باز است:
\begin{center}
\url{https://editor.salamlang.ir}
\end{center}

\Needspace{18\baselineskip}
\section*{آخرین برنامه}

هر سفری را باید با یک برنامه‌ی سلام بست. این هم آخرین برنامه‌ی کتاب، که
همان برنامه‌ی نخست است با یک خط بیشتر:

\begin{salamfa}
روال ریشه:
    سرچاپ "سلام، دنیا!"
    سرچاپ "تا دیدار، برنامه‌نویس!"
پایان
\end{salamfa}

\begin{outputfa}
سلام، دنیا!
تا دیدار، برنامه‌نویس!
\end{outputfa}

هر برنامه‌نویسی روزی با همین «سلام، دنیا!» شروع کرده است. شما هم شروع
کرده‌اید؛ ادامه‌ی راه با شماست.

\vspace{1.5em}
\begin{center}
{\color{brand}\rule{2.5cm}{1pt}}\par\vspace{0.6em}
{\salamcodefont\bfseries\Large\color{salamKeyword}پایان}
\end{center}


\end{document}
